% This file is part of Background Music.
%
% Background Music is free software: you can redistribute it and/or
% modify it under the terms of the GNU General Public License as
% published by the Free Software Foundation, either version 2 of the
% License, or (at your option) any later version.
%
% Background Music is distributed in the hope that it will be useful,
% but WITHOUT ANY WARRANTY; without even the implied warranty of
% MERCHANTABILITY or FITNESS FOR A PARTICULAR PURPOSE. See the
% GNU General Public License for more details.
%
% You should have received a copy of the GNU General Public License
% along with Background Music. If not, see <http://www.gnu.org/licenses/>.

%
% VolumeIcons.tex
%
% Build with XeTeX:
%   xelatex -jobname=Volume0 '\def\UseOption{}% This file is part of Background Music.
%
% Background Music is free software: you can redistribute it and/or
% modify it under the terms of the GNU General Public License as
% published by the Free Software Foundation, either version 2 of the
% License, or (at your option) any later version.
%
% Background Music is distributed in the hope that it will be useful,
% but WITHOUT ANY WARRANTY; without even the implied warranty of
% MERCHANTABILITY or FITNESS FOR A PARTICULAR PURPOSE. See the
% GNU General Public License for more details.
%
% You should have received a copy of the GNU General Public License
% along with Background Music. If not, see <http://www.gnu.org/licenses/>.

%
% VolumeIcons.tex
%
% Build with XeTeX:
%   xelatex -jobname=Volume0 '\def\UseOption{}% This file is part of Background Music.
%
% Background Music is free software: you can redistribute it and/or
% modify it under the terms of the GNU General Public License as
% published by the Free Software Foundation, either version 2 of the
% License, or (at your option) any later version.
%
% Background Music is distributed in the hope that it will be useful,
% but WITHOUT ANY WARRANTY; without even the implied warranty of
% MERCHANTABILITY or FITNESS FOR A PARTICULAR PURPOSE. See the
% GNU General Public License for more details.
%
% You should have received a copy of the GNU General Public License
% along with Background Music. If not, see <http://www.gnu.org/licenses/>.

%
% VolumeIcons.tex
%
% Build with XeTeX:
%   xelatex -jobname=Volume0 '\def\UseOption{}% This file is part of Background Music.
%
% Background Music is free software: you can redistribute it and/or
% modify it under the terms of the GNU General Public License as
% published by the Free Software Foundation, either version 2 of the
% License, or (at your option) any later version.
%
% Background Music is distributed in the hope that it will be useful,
% but WITHOUT ANY WARRANTY; without even the implied warranty of
% MERCHANTABILITY or FITNESS FOR A PARTICULAR PURPOSE. See the
% GNU General Public License for more details.
%
% You should have received a copy of the GNU General Public License
% along with Background Music. If not, see <http://www.gnu.org/licenses/>.

%
% VolumeIcons.tex
%
% Build with XeTeX:
%   xelatex -jobname=Volume0 '\def\UseOption{}\input{VolumeIcons.tex}'
%   xelatex -jobname=Volume1 '\def\UseOption{w1}\input{VolumeIcons.tex}'
%   xelatex -jobname=Volume2 '\def\UseOption{w1,w2}\input{VolumeIcons.tex}'
%   xelatex -jobname=Volume3 '\def\UseOption{w1,w2,w3}\input{VolumeIcons.tex}'
%   for n in 0 1 2 3; do mv Volume$n.pdf ../BGMApp/BGMApp/Images.xcassets/Volume$n.imageset/; done
%
% Might build correctly with regular LaTeX. I haven't tried it.
%

\documentclass[tikz]{standalone}
\usepackage{tikz}
% "dummyOption" prevents "Package optional Warning: No options were selected,
% so all optional text will be printed" when building Volume0.pdf.
\usepackage[dummyOption]{optional}

\begin{document}
\begin{tikzpicture}

% Speaker (Rounded box and triangle)
\fill[rounded corners=5mm]
    (0mm, 62.5mm) rectangle (25mm, 37.5mm) {};
\draw[rounded corners=2.5mm,fill=black]
    (3mm, 50mm)--(34mm, 76.5mm)--(34mm, 23.5mm)--cycle;

% First sound wave (Curved line)
\opt{w1}{
    \draw[line width=4.3mm,line cap=round]
        (44mm, 36.5mm) to[out=46,in=-46] (44mm, 63.5mm);
}

% Second sound wave (Curved line)
\opt{w2}{
    \draw[line width=4.3mm,line cap=round]
        (57.5mm, 27.5mm) to[out=46,in=-46] (57.5mm, 72.5mm);
}

% Third sound wave (Curved line)
\opt{w3}{
    \draw[line width=4.3mm,line cap=round]
        (72mm, 18.5mm) to[out=46,in=-46] (72mm, 81.5mm);
}

% Always draw a transparent copy of the third wave so the images will all have
% the same width.
\draw[line width=4.3mm,line cap=round,opacity=0]
    (72mm, 18.5mm) to[out=46,in=-46] (72mm, 81.5mm);

\end{tikzpicture}
\end{document}

'
%   xelatex -jobname=Volume1 '\def\UseOption{w1}% This file is part of Background Music.
%
% Background Music is free software: you can redistribute it and/or
% modify it under the terms of the GNU General Public License as
% published by the Free Software Foundation, either version 2 of the
% License, or (at your option) any later version.
%
% Background Music is distributed in the hope that it will be useful,
% but WITHOUT ANY WARRANTY; without even the implied warranty of
% MERCHANTABILITY or FITNESS FOR A PARTICULAR PURPOSE. See the
% GNU General Public License for more details.
%
% You should have received a copy of the GNU General Public License
% along with Background Music. If not, see <http://www.gnu.org/licenses/>.

%
% VolumeIcons.tex
%
% Build with XeTeX:
%   xelatex -jobname=Volume0 '\def\UseOption{}\input{VolumeIcons.tex}'
%   xelatex -jobname=Volume1 '\def\UseOption{w1}\input{VolumeIcons.tex}'
%   xelatex -jobname=Volume2 '\def\UseOption{w1,w2}\input{VolumeIcons.tex}'
%   xelatex -jobname=Volume3 '\def\UseOption{w1,w2,w3}\input{VolumeIcons.tex}'
%   for n in 0 1 2 3; do mv Volume$n.pdf ../BGMApp/BGMApp/Images.xcassets/Volume$n.imageset/; done
%
% Might build correctly with regular LaTeX. I haven't tried it.
%

\documentclass[tikz]{standalone}
\usepackage{tikz}
% "dummyOption" prevents "Package optional Warning: No options were selected,
% so all optional text will be printed" when building Volume0.pdf.
\usepackage[dummyOption]{optional}

\begin{document}
\begin{tikzpicture}

% Speaker (Rounded box and triangle)
\fill[rounded corners=5mm]
    (0mm, 62.5mm) rectangle (25mm, 37.5mm) {};
\draw[rounded corners=2.5mm,fill=black]
    (3mm, 50mm)--(34mm, 76.5mm)--(34mm, 23.5mm)--cycle;

% First sound wave (Curved line)
\opt{w1}{
    \draw[line width=4.3mm,line cap=round]
        (44mm, 36.5mm) to[out=46,in=-46] (44mm, 63.5mm);
}

% Second sound wave (Curved line)
\opt{w2}{
    \draw[line width=4.3mm,line cap=round]
        (57.5mm, 27.5mm) to[out=46,in=-46] (57.5mm, 72.5mm);
}

% Third sound wave (Curved line)
\opt{w3}{
    \draw[line width=4.3mm,line cap=round]
        (72mm, 18.5mm) to[out=46,in=-46] (72mm, 81.5mm);
}

% Always draw a transparent copy of the third wave so the images will all have
% the same width.
\draw[line width=4.3mm,line cap=round,opacity=0]
    (72mm, 18.5mm) to[out=46,in=-46] (72mm, 81.5mm);

\end{tikzpicture}
\end{document}

'
%   xelatex -jobname=Volume2 '\def\UseOption{w1,w2}% This file is part of Background Music.
%
% Background Music is free software: you can redistribute it and/or
% modify it under the terms of the GNU General Public License as
% published by the Free Software Foundation, either version 2 of the
% License, or (at your option) any later version.
%
% Background Music is distributed in the hope that it will be useful,
% but WITHOUT ANY WARRANTY; without even the implied warranty of
% MERCHANTABILITY or FITNESS FOR A PARTICULAR PURPOSE. See the
% GNU General Public License for more details.
%
% You should have received a copy of the GNU General Public License
% along with Background Music. If not, see <http://www.gnu.org/licenses/>.

%
% VolumeIcons.tex
%
% Build with XeTeX:
%   xelatex -jobname=Volume0 '\def\UseOption{}\input{VolumeIcons.tex}'
%   xelatex -jobname=Volume1 '\def\UseOption{w1}\input{VolumeIcons.tex}'
%   xelatex -jobname=Volume2 '\def\UseOption{w1,w2}\input{VolumeIcons.tex}'
%   xelatex -jobname=Volume3 '\def\UseOption{w1,w2,w3}\input{VolumeIcons.tex}'
%   for n in 0 1 2 3; do mv Volume$n.pdf ../BGMApp/BGMApp/Images.xcassets/Volume$n.imageset/; done
%
% Might build correctly with regular LaTeX. I haven't tried it.
%

\documentclass[tikz]{standalone}
\usepackage{tikz}
% "dummyOption" prevents "Package optional Warning: No options were selected,
% so all optional text will be printed" when building Volume0.pdf.
\usepackage[dummyOption]{optional}

\begin{document}
\begin{tikzpicture}

% Speaker (Rounded box and triangle)
\fill[rounded corners=5mm]
    (0mm, 62.5mm) rectangle (25mm, 37.5mm) {};
\draw[rounded corners=2.5mm,fill=black]
    (3mm, 50mm)--(34mm, 76.5mm)--(34mm, 23.5mm)--cycle;

% First sound wave (Curved line)
\opt{w1}{
    \draw[line width=4.3mm,line cap=round]
        (44mm, 36.5mm) to[out=46,in=-46] (44mm, 63.5mm);
}

% Second sound wave (Curved line)
\opt{w2}{
    \draw[line width=4.3mm,line cap=round]
        (57.5mm, 27.5mm) to[out=46,in=-46] (57.5mm, 72.5mm);
}

% Third sound wave (Curved line)
\opt{w3}{
    \draw[line width=4.3mm,line cap=round]
        (72mm, 18.5mm) to[out=46,in=-46] (72mm, 81.5mm);
}

% Always draw a transparent copy of the third wave so the images will all have
% the same width.
\draw[line width=4.3mm,line cap=round,opacity=0]
    (72mm, 18.5mm) to[out=46,in=-46] (72mm, 81.5mm);

\end{tikzpicture}
\end{document}

'
%   xelatex -jobname=Volume3 '\def\UseOption{w1,w2,w3}% This file is part of Background Music.
%
% Background Music is free software: you can redistribute it and/or
% modify it under the terms of the GNU General Public License as
% published by the Free Software Foundation, either version 2 of the
% License, or (at your option) any later version.
%
% Background Music is distributed in the hope that it will be useful,
% but WITHOUT ANY WARRANTY; without even the implied warranty of
% MERCHANTABILITY or FITNESS FOR A PARTICULAR PURPOSE. See the
% GNU General Public License for more details.
%
% You should have received a copy of the GNU General Public License
% along with Background Music. If not, see <http://www.gnu.org/licenses/>.

%
% VolumeIcons.tex
%
% Build with XeTeX:
%   xelatex -jobname=Volume0 '\def\UseOption{}\input{VolumeIcons.tex}'
%   xelatex -jobname=Volume1 '\def\UseOption{w1}\input{VolumeIcons.tex}'
%   xelatex -jobname=Volume2 '\def\UseOption{w1,w2}\input{VolumeIcons.tex}'
%   xelatex -jobname=Volume3 '\def\UseOption{w1,w2,w3}\input{VolumeIcons.tex}'
%   for n in 0 1 2 3; do mv Volume$n.pdf ../BGMApp/BGMApp/Images.xcassets/Volume$n.imageset/; done
%
% Might build correctly with regular LaTeX. I haven't tried it.
%

\documentclass[tikz]{standalone}
\usepackage{tikz}
% "dummyOption" prevents "Package optional Warning: No options were selected,
% so all optional text will be printed" when building Volume0.pdf.
\usepackage[dummyOption]{optional}

\begin{document}
\begin{tikzpicture}

% Speaker (Rounded box and triangle)
\fill[rounded corners=5mm]
    (0mm, 62.5mm) rectangle (25mm, 37.5mm) {};
\draw[rounded corners=2.5mm,fill=black]
    (3mm, 50mm)--(34mm, 76.5mm)--(34mm, 23.5mm)--cycle;

% First sound wave (Curved line)
\opt{w1}{
    \draw[line width=4.3mm,line cap=round]
        (44mm, 36.5mm) to[out=46,in=-46] (44mm, 63.5mm);
}

% Second sound wave (Curved line)
\opt{w2}{
    \draw[line width=4.3mm,line cap=round]
        (57.5mm, 27.5mm) to[out=46,in=-46] (57.5mm, 72.5mm);
}

% Third sound wave (Curved line)
\opt{w3}{
    \draw[line width=4.3mm,line cap=round]
        (72mm, 18.5mm) to[out=46,in=-46] (72mm, 81.5mm);
}

% Always draw a transparent copy of the third wave so the images will all have
% the same width.
\draw[line width=4.3mm,line cap=round,opacity=0]
    (72mm, 18.5mm) to[out=46,in=-46] (72mm, 81.5mm);

\end{tikzpicture}
\end{document}

'
%   for n in 0 1 2 3; do mv Volume$n.pdf ../BGMApp/BGMApp/Images.xcassets/Volume$n.imageset/; done
%
% Might build correctly with regular LaTeX. I haven't tried it.
%

\documentclass[tikz]{standalone}
\usepackage{tikz}
% "dummyOption" prevents "Package optional Warning: No options were selected,
% so all optional text will be printed" when building Volume0.pdf.
\usepackage[dummyOption]{optional}

\begin{document}
\begin{tikzpicture}

% Speaker (Rounded box and triangle)
\fill[rounded corners=5mm]
    (0mm, 62.5mm) rectangle (25mm, 37.5mm) {};
\draw[rounded corners=2.5mm,fill=black]
    (3mm, 50mm)--(34mm, 76.5mm)--(34mm, 23.5mm)--cycle;

% First sound wave (Curved line)
\opt{w1}{
    \draw[line width=4.3mm,line cap=round]
        (44mm, 36.5mm) to[out=46,in=-46] (44mm, 63.5mm);
}

% Second sound wave (Curved line)
\opt{w2}{
    \draw[line width=4.3mm,line cap=round]
        (57.5mm, 27.5mm) to[out=46,in=-46] (57.5mm, 72.5mm);
}

% Third sound wave (Curved line)
\opt{w3}{
    \draw[line width=4.3mm,line cap=round]
        (72mm, 18.5mm) to[out=46,in=-46] (72mm, 81.5mm);
}

% Always draw a transparent copy of the third wave so the images will all have
% the same width.
\draw[line width=4.3mm,line cap=round,opacity=0]
    (72mm, 18.5mm) to[out=46,in=-46] (72mm, 81.5mm);

\end{tikzpicture}
\end{document}

'
%   xelatex -jobname=Volume1 '\def\UseOption{w1}% This file is part of Background Music.
%
% Background Music is free software: you can redistribute it and/or
% modify it under the terms of the GNU General Public License as
% published by the Free Software Foundation, either version 2 of the
% License, or (at your option) any later version.
%
% Background Music is distributed in the hope that it will be useful,
% but WITHOUT ANY WARRANTY; without even the implied warranty of
% MERCHANTABILITY or FITNESS FOR A PARTICULAR PURPOSE. See the
% GNU General Public License for more details.
%
% You should have received a copy of the GNU General Public License
% along with Background Music. If not, see <http://www.gnu.org/licenses/>.

%
% VolumeIcons.tex
%
% Build with XeTeX:
%   xelatex -jobname=Volume0 '\def\UseOption{}% This file is part of Background Music.
%
% Background Music is free software: you can redistribute it and/or
% modify it under the terms of the GNU General Public License as
% published by the Free Software Foundation, either version 2 of the
% License, or (at your option) any later version.
%
% Background Music is distributed in the hope that it will be useful,
% but WITHOUT ANY WARRANTY; without even the implied warranty of
% MERCHANTABILITY or FITNESS FOR A PARTICULAR PURPOSE. See the
% GNU General Public License for more details.
%
% You should have received a copy of the GNU General Public License
% along with Background Music. If not, see <http://www.gnu.org/licenses/>.

%
% VolumeIcons.tex
%
% Build with XeTeX:
%   xelatex -jobname=Volume0 '\def\UseOption{}\input{VolumeIcons.tex}'
%   xelatex -jobname=Volume1 '\def\UseOption{w1}\input{VolumeIcons.tex}'
%   xelatex -jobname=Volume2 '\def\UseOption{w1,w2}\input{VolumeIcons.tex}'
%   xelatex -jobname=Volume3 '\def\UseOption{w1,w2,w3}\input{VolumeIcons.tex}'
%   for n in 0 1 2 3; do mv Volume$n.pdf ../BGMApp/BGMApp/Images.xcassets/Volume$n.imageset/; done
%
% Might build correctly with regular LaTeX. I haven't tried it.
%

\documentclass[tikz]{standalone}
\usepackage{tikz}
% "dummyOption" prevents "Package optional Warning: No options were selected,
% so all optional text will be printed" when building Volume0.pdf.
\usepackage[dummyOption]{optional}

\begin{document}
\begin{tikzpicture}

% Speaker (Rounded box and triangle)
\fill[rounded corners=5mm]
    (0mm, 62.5mm) rectangle (25mm, 37.5mm) {};
\draw[rounded corners=2.5mm,fill=black]
    (3mm, 50mm)--(34mm, 76.5mm)--(34mm, 23.5mm)--cycle;

% First sound wave (Curved line)
\opt{w1}{
    \draw[line width=4.3mm,line cap=round]
        (44mm, 36.5mm) to[out=46,in=-46] (44mm, 63.5mm);
}

% Second sound wave (Curved line)
\opt{w2}{
    \draw[line width=4.3mm,line cap=round]
        (57.5mm, 27.5mm) to[out=46,in=-46] (57.5mm, 72.5mm);
}

% Third sound wave (Curved line)
\opt{w3}{
    \draw[line width=4.3mm,line cap=round]
        (72mm, 18.5mm) to[out=46,in=-46] (72mm, 81.5mm);
}

% Always draw a transparent copy of the third wave so the images will all have
% the same width.
\draw[line width=4.3mm,line cap=round,opacity=0]
    (72mm, 18.5mm) to[out=46,in=-46] (72mm, 81.5mm);

\end{tikzpicture}
\end{document}

'
%   xelatex -jobname=Volume1 '\def\UseOption{w1}% This file is part of Background Music.
%
% Background Music is free software: you can redistribute it and/or
% modify it under the terms of the GNU General Public License as
% published by the Free Software Foundation, either version 2 of the
% License, or (at your option) any later version.
%
% Background Music is distributed in the hope that it will be useful,
% but WITHOUT ANY WARRANTY; without even the implied warranty of
% MERCHANTABILITY or FITNESS FOR A PARTICULAR PURPOSE. See the
% GNU General Public License for more details.
%
% You should have received a copy of the GNU General Public License
% along with Background Music. If not, see <http://www.gnu.org/licenses/>.

%
% VolumeIcons.tex
%
% Build with XeTeX:
%   xelatex -jobname=Volume0 '\def\UseOption{}\input{VolumeIcons.tex}'
%   xelatex -jobname=Volume1 '\def\UseOption{w1}\input{VolumeIcons.tex}'
%   xelatex -jobname=Volume2 '\def\UseOption{w1,w2}\input{VolumeIcons.tex}'
%   xelatex -jobname=Volume3 '\def\UseOption{w1,w2,w3}\input{VolumeIcons.tex}'
%   for n in 0 1 2 3; do mv Volume$n.pdf ../BGMApp/BGMApp/Images.xcassets/Volume$n.imageset/; done
%
% Might build correctly with regular LaTeX. I haven't tried it.
%

\documentclass[tikz]{standalone}
\usepackage{tikz}
% "dummyOption" prevents "Package optional Warning: No options were selected,
% so all optional text will be printed" when building Volume0.pdf.
\usepackage[dummyOption]{optional}

\begin{document}
\begin{tikzpicture}

% Speaker (Rounded box and triangle)
\fill[rounded corners=5mm]
    (0mm, 62.5mm) rectangle (25mm, 37.5mm) {};
\draw[rounded corners=2.5mm,fill=black]
    (3mm, 50mm)--(34mm, 76.5mm)--(34mm, 23.5mm)--cycle;

% First sound wave (Curved line)
\opt{w1}{
    \draw[line width=4.3mm,line cap=round]
        (44mm, 36.5mm) to[out=46,in=-46] (44mm, 63.5mm);
}

% Second sound wave (Curved line)
\opt{w2}{
    \draw[line width=4.3mm,line cap=round]
        (57.5mm, 27.5mm) to[out=46,in=-46] (57.5mm, 72.5mm);
}

% Third sound wave (Curved line)
\opt{w3}{
    \draw[line width=4.3mm,line cap=round]
        (72mm, 18.5mm) to[out=46,in=-46] (72mm, 81.5mm);
}

% Always draw a transparent copy of the third wave so the images will all have
% the same width.
\draw[line width=4.3mm,line cap=round,opacity=0]
    (72mm, 18.5mm) to[out=46,in=-46] (72mm, 81.5mm);

\end{tikzpicture}
\end{document}

'
%   xelatex -jobname=Volume2 '\def\UseOption{w1,w2}% This file is part of Background Music.
%
% Background Music is free software: you can redistribute it and/or
% modify it under the terms of the GNU General Public License as
% published by the Free Software Foundation, either version 2 of the
% License, or (at your option) any later version.
%
% Background Music is distributed in the hope that it will be useful,
% but WITHOUT ANY WARRANTY; without even the implied warranty of
% MERCHANTABILITY or FITNESS FOR A PARTICULAR PURPOSE. See the
% GNU General Public License for more details.
%
% You should have received a copy of the GNU General Public License
% along with Background Music. If not, see <http://www.gnu.org/licenses/>.

%
% VolumeIcons.tex
%
% Build with XeTeX:
%   xelatex -jobname=Volume0 '\def\UseOption{}\input{VolumeIcons.tex}'
%   xelatex -jobname=Volume1 '\def\UseOption{w1}\input{VolumeIcons.tex}'
%   xelatex -jobname=Volume2 '\def\UseOption{w1,w2}\input{VolumeIcons.tex}'
%   xelatex -jobname=Volume3 '\def\UseOption{w1,w2,w3}\input{VolumeIcons.tex}'
%   for n in 0 1 2 3; do mv Volume$n.pdf ../BGMApp/BGMApp/Images.xcassets/Volume$n.imageset/; done
%
% Might build correctly with regular LaTeX. I haven't tried it.
%

\documentclass[tikz]{standalone}
\usepackage{tikz}
% "dummyOption" prevents "Package optional Warning: No options were selected,
% so all optional text will be printed" when building Volume0.pdf.
\usepackage[dummyOption]{optional}

\begin{document}
\begin{tikzpicture}

% Speaker (Rounded box and triangle)
\fill[rounded corners=5mm]
    (0mm, 62.5mm) rectangle (25mm, 37.5mm) {};
\draw[rounded corners=2.5mm,fill=black]
    (3mm, 50mm)--(34mm, 76.5mm)--(34mm, 23.5mm)--cycle;

% First sound wave (Curved line)
\opt{w1}{
    \draw[line width=4.3mm,line cap=round]
        (44mm, 36.5mm) to[out=46,in=-46] (44mm, 63.5mm);
}

% Second sound wave (Curved line)
\opt{w2}{
    \draw[line width=4.3mm,line cap=round]
        (57.5mm, 27.5mm) to[out=46,in=-46] (57.5mm, 72.5mm);
}

% Third sound wave (Curved line)
\opt{w3}{
    \draw[line width=4.3mm,line cap=round]
        (72mm, 18.5mm) to[out=46,in=-46] (72mm, 81.5mm);
}

% Always draw a transparent copy of the third wave so the images will all have
% the same width.
\draw[line width=4.3mm,line cap=round,opacity=0]
    (72mm, 18.5mm) to[out=46,in=-46] (72mm, 81.5mm);

\end{tikzpicture}
\end{document}

'
%   xelatex -jobname=Volume3 '\def\UseOption{w1,w2,w3}% This file is part of Background Music.
%
% Background Music is free software: you can redistribute it and/or
% modify it under the terms of the GNU General Public License as
% published by the Free Software Foundation, either version 2 of the
% License, or (at your option) any later version.
%
% Background Music is distributed in the hope that it will be useful,
% but WITHOUT ANY WARRANTY; without even the implied warranty of
% MERCHANTABILITY or FITNESS FOR A PARTICULAR PURPOSE. See the
% GNU General Public License for more details.
%
% You should have received a copy of the GNU General Public License
% along with Background Music. If not, see <http://www.gnu.org/licenses/>.

%
% VolumeIcons.tex
%
% Build with XeTeX:
%   xelatex -jobname=Volume0 '\def\UseOption{}\input{VolumeIcons.tex}'
%   xelatex -jobname=Volume1 '\def\UseOption{w1}\input{VolumeIcons.tex}'
%   xelatex -jobname=Volume2 '\def\UseOption{w1,w2}\input{VolumeIcons.tex}'
%   xelatex -jobname=Volume3 '\def\UseOption{w1,w2,w3}\input{VolumeIcons.tex}'
%   for n in 0 1 2 3; do mv Volume$n.pdf ../BGMApp/BGMApp/Images.xcassets/Volume$n.imageset/; done
%
% Might build correctly with regular LaTeX. I haven't tried it.
%

\documentclass[tikz]{standalone}
\usepackage{tikz}
% "dummyOption" prevents "Package optional Warning: No options were selected,
% so all optional text will be printed" when building Volume0.pdf.
\usepackage[dummyOption]{optional}

\begin{document}
\begin{tikzpicture}

% Speaker (Rounded box and triangle)
\fill[rounded corners=5mm]
    (0mm, 62.5mm) rectangle (25mm, 37.5mm) {};
\draw[rounded corners=2.5mm,fill=black]
    (3mm, 50mm)--(34mm, 76.5mm)--(34mm, 23.5mm)--cycle;

% First sound wave (Curved line)
\opt{w1}{
    \draw[line width=4.3mm,line cap=round]
        (44mm, 36.5mm) to[out=46,in=-46] (44mm, 63.5mm);
}

% Second sound wave (Curved line)
\opt{w2}{
    \draw[line width=4.3mm,line cap=round]
        (57.5mm, 27.5mm) to[out=46,in=-46] (57.5mm, 72.5mm);
}

% Third sound wave (Curved line)
\opt{w3}{
    \draw[line width=4.3mm,line cap=round]
        (72mm, 18.5mm) to[out=46,in=-46] (72mm, 81.5mm);
}

% Always draw a transparent copy of the third wave so the images will all have
% the same width.
\draw[line width=4.3mm,line cap=round,opacity=0]
    (72mm, 18.5mm) to[out=46,in=-46] (72mm, 81.5mm);

\end{tikzpicture}
\end{document}

'
%   for n in 0 1 2 3; do mv Volume$n.pdf ../BGMApp/BGMApp/Images.xcassets/Volume$n.imageset/; done
%
% Might build correctly with regular LaTeX. I haven't tried it.
%

\documentclass[tikz]{standalone}
\usepackage{tikz}
% "dummyOption" prevents "Package optional Warning: No options were selected,
% so all optional text will be printed" when building Volume0.pdf.
\usepackage[dummyOption]{optional}

\begin{document}
\begin{tikzpicture}

% Speaker (Rounded box and triangle)
\fill[rounded corners=5mm]
    (0mm, 62.5mm) rectangle (25mm, 37.5mm) {};
\draw[rounded corners=2.5mm,fill=black]
    (3mm, 50mm)--(34mm, 76.5mm)--(34mm, 23.5mm)--cycle;

% First sound wave (Curved line)
\opt{w1}{
    \draw[line width=4.3mm,line cap=round]
        (44mm, 36.5mm) to[out=46,in=-46] (44mm, 63.5mm);
}

% Second sound wave (Curved line)
\opt{w2}{
    \draw[line width=4.3mm,line cap=round]
        (57.5mm, 27.5mm) to[out=46,in=-46] (57.5mm, 72.5mm);
}

% Third sound wave (Curved line)
\opt{w3}{
    \draw[line width=4.3mm,line cap=round]
        (72mm, 18.5mm) to[out=46,in=-46] (72mm, 81.5mm);
}

% Always draw a transparent copy of the third wave so the images will all have
% the same width.
\draw[line width=4.3mm,line cap=round,opacity=0]
    (72mm, 18.5mm) to[out=46,in=-46] (72mm, 81.5mm);

\end{tikzpicture}
\end{document}

'
%   xelatex -jobname=Volume2 '\def\UseOption{w1,w2}% This file is part of Background Music.
%
% Background Music is free software: you can redistribute it and/or
% modify it under the terms of the GNU General Public License as
% published by the Free Software Foundation, either version 2 of the
% License, or (at your option) any later version.
%
% Background Music is distributed in the hope that it will be useful,
% but WITHOUT ANY WARRANTY; without even the implied warranty of
% MERCHANTABILITY or FITNESS FOR A PARTICULAR PURPOSE. See the
% GNU General Public License for more details.
%
% You should have received a copy of the GNU General Public License
% along with Background Music. If not, see <http://www.gnu.org/licenses/>.

%
% VolumeIcons.tex
%
% Build with XeTeX:
%   xelatex -jobname=Volume0 '\def\UseOption{}% This file is part of Background Music.
%
% Background Music is free software: you can redistribute it and/or
% modify it under the terms of the GNU General Public License as
% published by the Free Software Foundation, either version 2 of the
% License, or (at your option) any later version.
%
% Background Music is distributed in the hope that it will be useful,
% but WITHOUT ANY WARRANTY; without even the implied warranty of
% MERCHANTABILITY or FITNESS FOR A PARTICULAR PURPOSE. See the
% GNU General Public License for more details.
%
% You should have received a copy of the GNU General Public License
% along with Background Music. If not, see <http://www.gnu.org/licenses/>.

%
% VolumeIcons.tex
%
% Build with XeTeX:
%   xelatex -jobname=Volume0 '\def\UseOption{}\input{VolumeIcons.tex}'
%   xelatex -jobname=Volume1 '\def\UseOption{w1}\input{VolumeIcons.tex}'
%   xelatex -jobname=Volume2 '\def\UseOption{w1,w2}\input{VolumeIcons.tex}'
%   xelatex -jobname=Volume3 '\def\UseOption{w1,w2,w3}\input{VolumeIcons.tex}'
%   for n in 0 1 2 3; do mv Volume$n.pdf ../BGMApp/BGMApp/Images.xcassets/Volume$n.imageset/; done
%
% Might build correctly with regular LaTeX. I haven't tried it.
%

\documentclass[tikz]{standalone}
\usepackage{tikz}
% "dummyOption" prevents "Package optional Warning: No options were selected,
% so all optional text will be printed" when building Volume0.pdf.
\usepackage[dummyOption]{optional}

\begin{document}
\begin{tikzpicture}

% Speaker (Rounded box and triangle)
\fill[rounded corners=5mm]
    (0mm, 62.5mm) rectangle (25mm, 37.5mm) {};
\draw[rounded corners=2.5mm,fill=black]
    (3mm, 50mm)--(34mm, 76.5mm)--(34mm, 23.5mm)--cycle;

% First sound wave (Curved line)
\opt{w1}{
    \draw[line width=4.3mm,line cap=round]
        (44mm, 36.5mm) to[out=46,in=-46] (44mm, 63.5mm);
}

% Second sound wave (Curved line)
\opt{w2}{
    \draw[line width=4.3mm,line cap=round]
        (57.5mm, 27.5mm) to[out=46,in=-46] (57.5mm, 72.5mm);
}

% Third sound wave (Curved line)
\opt{w3}{
    \draw[line width=4.3mm,line cap=round]
        (72mm, 18.5mm) to[out=46,in=-46] (72mm, 81.5mm);
}

% Always draw a transparent copy of the third wave so the images will all have
% the same width.
\draw[line width=4.3mm,line cap=round,opacity=0]
    (72mm, 18.5mm) to[out=46,in=-46] (72mm, 81.5mm);

\end{tikzpicture}
\end{document}

'
%   xelatex -jobname=Volume1 '\def\UseOption{w1}% This file is part of Background Music.
%
% Background Music is free software: you can redistribute it and/or
% modify it under the terms of the GNU General Public License as
% published by the Free Software Foundation, either version 2 of the
% License, or (at your option) any later version.
%
% Background Music is distributed in the hope that it will be useful,
% but WITHOUT ANY WARRANTY; without even the implied warranty of
% MERCHANTABILITY or FITNESS FOR A PARTICULAR PURPOSE. See the
% GNU General Public License for more details.
%
% You should have received a copy of the GNU General Public License
% along with Background Music. If not, see <http://www.gnu.org/licenses/>.

%
% VolumeIcons.tex
%
% Build with XeTeX:
%   xelatex -jobname=Volume0 '\def\UseOption{}\input{VolumeIcons.tex}'
%   xelatex -jobname=Volume1 '\def\UseOption{w1}\input{VolumeIcons.tex}'
%   xelatex -jobname=Volume2 '\def\UseOption{w1,w2}\input{VolumeIcons.tex}'
%   xelatex -jobname=Volume3 '\def\UseOption{w1,w2,w3}\input{VolumeIcons.tex}'
%   for n in 0 1 2 3; do mv Volume$n.pdf ../BGMApp/BGMApp/Images.xcassets/Volume$n.imageset/; done
%
% Might build correctly with regular LaTeX. I haven't tried it.
%

\documentclass[tikz]{standalone}
\usepackage{tikz}
% "dummyOption" prevents "Package optional Warning: No options were selected,
% so all optional text will be printed" when building Volume0.pdf.
\usepackage[dummyOption]{optional}

\begin{document}
\begin{tikzpicture}

% Speaker (Rounded box and triangle)
\fill[rounded corners=5mm]
    (0mm, 62.5mm) rectangle (25mm, 37.5mm) {};
\draw[rounded corners=2.5mm,fill=black]
    (3mm, 50mm)--(34mm, 76.5mm)--(34mm, 23.5mm)--cycle;

% First sound wave (Curved line)
\opt{w1}{
    \draw[line width=4.3mm,line cap=round]
        (44mm, 36.5mm) to[out=46,in=-46] (44mm, 63.5mm);
}

% Second sound wave (Curved line)
\opt{w2}{
    \draw[line width=4.3mm,line cap=round]
        (57.5mm, 27.5mm) to[out=46,in=-46] (57.5mm, 72.5mm);
}

% Third sound wave (Curved line)
\opt{w3}{
    \draw[line width=4.3mm,line cap=round]
        (72mm, 18.5mm) to[out=46,in=-46] (72mm, 81.5mm);
}

% Always draw a transparent copy of the third wave so the images will all have
% the same width.
\draw[line width=4.3mm,line cap=round,opacity=0]
    (72mm, 18.5mm) to[out=46,in=-46] (72mm, 81.5mm);

\end{tikzpicture}
\end{document}

'
%   xelatex -jobname=Volume2 '\def\UseOption{w1,w2}% This file is part of Background Music.
%
% Background Music is free software: you can redistribute it and/or
% modify it under the terms of the GNU General Public License as
% published by the Free Software Foundation, either version 2 of the
% License, or (at your option) any later version.
%
% Background Music is distributed in the hope that it will be useful,
% but WITHOUT ANY WARRANTY; without even the implied warranty of
% MERCHANTABILITY or FITNESS FOR A PARTICULAR PURPOSE. See the
% GNU General Public License for more details.
%
% You should have received a copy of the GNU General Public License
% along with Background Music. If not, see <http://www.gnu.org/licenses/>.

%
% VolumeIcons.tex
%
% Build with XeTeX:
%   xelatex -jobname=Volume0 '\def\UseOption{}\input{VolumeIcons.tex}'
%   xelatex -jobname=Volume1 '\def\UseOption{w1}\input{VolumeIcons.tex}'
%   xelatex -jobname=Volume2 '\def\UseOption{w1,w2}\input{VolumeIcons.tex}'
%   xelatex -jobname=Volume3 '\def\UseOption{w1,w2,w3}\input{VolumeIcons.tex}'
%   for n in 0 1 2 3; do mv Volume$n.pdf ../BGMApp/BGMApp/Images.xcassets/Volume$n.imageset/; done
%
% Might build correctly with regular LaTeX. I haven't tried it.
%

\documentclass[tikz]{standalone}
\usepackage{tikz}
% "dummyOption" prevents "Package optional Warning: No options were selected,
% so all optional text will be printed" when building Volume0.pdf.
\usepackage[dummyOption]{optional}

\begin{document}
\begin{tikzpicture}

% Speaker (Rounded box and triangle)
\fill[rounded corners=5mm]
    (0mm, 62.5mm) rectangle (25mm, 37.5mm) {};
\draw[rounded corners=2.5mm,fill=black]
    (3mm, 50mm)--(34mm, 76.5mm)--(34mm, 23.5mm)--cycle;

% First sound wave (Curved line)
\opt{w1}{
    \draw[line width=4.3mm,line cap=round]
        (44mm, 36.5mm) to[out=46,in=-46] (44mm, 63.5mm);
}

% Second sound wave (Curved line)
\opt{w2}{
    \draw[line width=4.3mm,line cap=round]
        (57.5mm, 27.5mm) to[out=46,in=-46] (57.5mm, 72.5mm);
}

% Third sound wave (Curved line)
\opt{w3}{
    \draw[line width=4.3mm,line cap=round]
        (72mm, 18.5mm) to[out=46,in=-46] (72mm, 81.5mm);
}

% Always draw a transparent copy of the third wave so the images will all have
% the same width.
\draw[line width=4.3mm,line cap=round,opacity=0]
    (72mm, 18.5mm) to[out=46,in=-46] (72mm, 81.5mm);

\end{tikzpicture}
\end{document}

'
%   xelatex -jobname=Volume3 '\def\UseOption{w1,w2,w3}% This file is part of Background Music.
%
% Background Music is free software: you can redistribute it and/or
% modify it under the terms of the GNU General Public License as
% published by the Free Software Foundation, either version 2 of the
% License, or (at your option) any later version.
%
% Background Music is distributed in the hope that it will be useful,
% but WITHOUT ANY WARRANTY; without even the implied warranty of
% MERCHANTABILITY or FITNESS FOR A PARTICULAR PURPOSE. See the
% GNU General Public License for more details.
%
% You should have received a copy of the GNU General Public License
% along with Background Music. If not, see <http://www.gnu.org/licenses/>.

%
% VolumeIcons.tex
%
% Build with XeTeX:
%   xelatex -jobname=Volume0 '\def\UseOption{}\input{VolumeIcons.tex}'
%   xelatex -jobname=Volume1 '\def\UseOption{w1}\input{VolumeIcons.tex}'
%   xelatex -jobname=Volume2 '\def\UseOption{w1,w2}\input{VolumeIcons.tex}'
%   xelatex -jobname=Volume3 '\def\UseOption{w1,w2,w3}\input{VolumeIcons.tex}'
%   for n in 0 1 2 3; do mv Volume$n.pdf ../BGMApp/BGMApp/Images.xcassets/Volume$n.imageset/; done
%
% Might build correctly with regular LaTeX. I haven't tried it.
%

\documentclass[tikz]{standalone}
\usepackage{tikz}
% "dummyOption" prevents "Package optional Warning: No options were selected,
% so all optional text will be printed" when building Volume0.pdf.
\usepackage[dummyOption]{optional}

\begin{document}
\begin{tikzpicture}

% Speaker (Rounded box and triangle)
\fill[rounded corners=5mm]
    (0mm, 62.5mm) rectangle (25mm, 37.5mm) {};
\draw[rounded corners=2.5mm,fill=black]
    (3mm, 50mm)--(34mm, 76.5mm)--(34mm, 23.5mm)--cycle;

% First sound wave (Curved line)
\opt{w1}{
    \draw[line width=4.3mm,line cap=round]
        (44mm, 36.5mm) to[out=46,in=-46] (44mm, 63.5mm);
}

% Second sound wave (Curved line)
\opt{w2}{
    \draw[line width=4.3mm,line cap=round]
        (57.5mm, 27.5mm) to[out=46,in=-46] (57.5mm, 72.5mm);
}

% Third sound wave (Curved line)
\opt{w3}{
    \draw[line width=4.3mm,line cap=round]
        (72mm, 18.5mm) to[out=46,in=-46] (72mm, 81.5mm);
}

% Always draw a transparent copy of the third wave so the images will all have
% the same width.
\draw[line width=4.3mm,line cap=round,opacity=0]
    (72mm, 18.5mm) to[out=46,in=-46] (72mm, 81.5mm);

\end{tikzpicture}
\end{document}

'
%   for n in 0 1 2 3; do mv Volume$n.pdf ../BGMApp/BGMApp/Images.xcassets/Volume$n.imageset/; done
%
% Might build correctly with regular LaTeX. I haven't tried it.
%

\documentclass[tikz]{standalone}
\usepackage{tikz}
% "dummyOption" prevents "Package optional Warning: No options were selected,
% so all optional text will be printed" when building Volume0.pdf.
\usepackage[dummyOption]{optional}

\begin{document}
\begin{tikzpicture}

% Speaker (Rounded box and triangle)
\fill[rounded corners=5mm]
    (0mm, 62.5mm) rectangle (25mm, 37.5mm) {};
\draw[rounded corners=2.5mm,fill=black]
    (3mm, 50mm)--(34mm, 76.5mm)--(34mm, 23.5mm)--cycle;

% First sound wave (Curved line)
\opt{w1}{
    \draw[line width=4.3mm,line cap=round]
        (44mm, 36.5mm) to[out=46,in=-46] (44mm, 63.5mm);
}

% Second sound wave (Curved line)
\opt{w2}{
    \draw[line width=4.3mm,line cap=round]
        (57.5mm, 27.5mm) to[out=46,in=-46] (57.5mm, 72.5mm);
}

% Third sound wave (Curved line)
\opt{w3}{
    \draw[line width=4.3mm,line cap=round]
        (72mm, 18.5mm) to[out=46,in=-46] (72mm, 81.5mm);
}

% Always draw a transparent copy of the third wave so the images will all have
% the same width.
\draw[line width=4.3mm,line cap=round,opacity=0]
    (72mm, 18.5mm) to[out=46,in=-46] (72mm, 81.5mm);

\end{tikzpicture}
\end{document}

'
%   xelatex -jobname=Volume3 '\def\UseOption{w1,w2,w3}% This file is part of Background Music.
%
% Background Music is free software: you can redistribute it and/or
% modify it under the terms of the GNU General Public License as
% published by the Free Software Foundation, either version 2 of the
% License, or (at your option) any later version.
%
% Background Music is distributed in the hope that it will be useful,
% but WITHOUT ANY WARRANTY; without even the implied warranty of
% MERCHANTABILITY or FITNESS FOR A PARTICULAR PURPOSE. See the
% GNU General Public License for more details.
%
% You should have received a copy of the GNU General Public License
% along with Background Music. If not, see <http://www.gnu.org/licenses/>.

%
% VolumeIcons.tex
%
% Build with XeTeX:
%   xelatex -jobname=Volume0 '\def\UseOption{}% This file is part of Background Music.
%
% Background Music is free software: you can redistribute it and/or
% modify it under the terms of the GNU General Public License as
% published by the Free Software Foundation, either version 2 of the
% License, or (at your option) any later version.
%
% Background Music is distributed in the hope that it will be useful,
% but WITHOUT ANY WARRANTY; without even the implied warranty of
% MERCHANTABILITY or FITNESS FOR A PARTICULAR PURPOSE. See the
% GNU General Public License for more details.
%
% You should have received a copy of the GNU General Public License
% along with Background Music. If not, see <http://www.gnu.org/licenses/>.

%
% VolumeIcons.tex
%
% Build with XeTeX:
%   xelatex -jobname=Volume0 '\def\UseOption{}\input{VolumeIcons.tex}'
%   xelatex -jobname=Volume1 '\def\UseOption{w1}\input{VolumeIcons.tex}'
%   xelatex -jobname=Volume2 '\def\UseOption{w1,w2}\input{VolumeIcons.tex}'
%   xelatex -jobname=Volume3 '\def\UseOption{w1,w2,w3}\input{VolumeIcons.tex}'
%   for n in 0 1 2 3; do mv Volume$n.pdf ../BGMApp/BGMApp/Images.xcassets/Volume$n.imageset/; done
%
% Might build correctly with regular LaTeX. I haven't tried it.
%

\documentclass[tikz]{standalone}
\usepackage{tikz}
% "dummyOption" prevents "Package optional Warning: No options were selected,
% so all optional text will be printed" when building Volume0.pdf.
\usepackage[dummyOption]{optional}

\begin{document}
\begin{tikzpicture}

% Speaker (Rounded box and triangle)
\fill[rounded corners=5mm]
    (0mm, 62.5mm) rectangle (25mm, 37.5mm) {};
\draw[rounded corners=2.5mm,fill=black]
    (3mm, 50mm)--(34mm, 76.5mm)--(34mm, 23.5mm)--cycle;

% First sound wave (Curved line)
\opt{w1}{
    \draw[line width=4.3mm,line cap=round]
        (44mm, 36.5mm) to[out=46,in=-46] (44mm, 63.5mm);
}

% Second sound wave (Curved line)
\opt{w2}{
    \draw[line width=4.3mm,line cap=round]
        (57.5mm, 27.5mm) to[out=46,in=-46] (57.5mm, 72.5mm);
}

% Third sound wave (Curved line)
\opt{w3}{
    \draw[line width=4.3mm,line cap=round]
        (72mm, 18.5mm) to[out=46,in=-46] (72mm, 81.5mm);
}

% Always draw a transparent copy of the third wave so the images will all have
% the same width.
\draw[line width=4.3mm,line cap=round,opacity=0]
    (72mm, 18.5mm) to[out=46,in=-46] (72mm, 81.5mm);

\end{tikzpicture}
\end{document}

'
%   xelatex -jobname=Volume1 '\def\UseOption{w1}% This file is part of Background Music.
%
% Background Music is free software: you can redistribute it and/or
% modify it under the terms of the GNU General Public License as
% published by the Free Software Foundation, either version 2 of the
% License, or (at your option) any later version.
%
% Background Music is distributed in the hope that it will be useful,
% but WITHOUT ANY WARRANTY; without even the implied warranty of
% MERCHANTABILITY or FITNESS FOR A PARTICULAR PURPOSE. See the
% GNU General Public License for more details.
%
% You should have received a copy of the GNU General Public License
% along with Background Music. If not, see <http://www.gnu.org/licenses/>.

%
% VolumeIcons.tex
%
% Build with XeTeX:
%   xelatex -jobname=Volume0 '\def\UseOption{}\input{VolumeIcons.tex}'
%   xelatex -jobname=Volume1 '\def\UseOption{w1}\input{VolumeIcons.tex}'
%   xelatex -jobname=Volume2 '\def\UseOption{w1,w2}\input{VolumeIcons.tex}'
%   xelatex -jobname=Volume3 '\def\UseOption{w1,w2,w3}\input{VolumeIcons.tex}'
%   for n in 0 1 2 3; do mv Volume$n.pdf ../BGMApp/BGMApp/Images.xcassets/Volume$n.imageset/; done
%
% Might build correctly with regular LaTeX. I haven't tried it.
%

\documentclass[tikz]{standalone}
\usepackage{tikz}
% "dummyOption" prevents "Package optional Warning: No options were selected,
% so all optional text will be printed" when building Volume0.pdf.
\usepackage[dummyOption]{optional}

\begin{document}
\begin{tikzpicture}

% Speaker (Rounded box and triangle)
\fill[rounded corners=5mm]
    (0mm, 62.5mm) rectangle (25mm, 37.5mm) {};
\draw[rounded corners=2.5mm,fill=black]
    (3mm, 50mm)--(34mm, 76.5mm)--(34mm, 23.5mm)--cycle;

% First sound wave (Curved line)
\opt{w1}{
    \draw[line width=4.3mm,line cap=round]
        (44mm, 36.5mm) to[out=46,in=-46] (44mm, 63.5mm);
}

% Second sound wave (Curved line)
\opt{w2}{
    \draw[line width=4.3mm,line cap=round]
        (57.5mm, 27.5mm) to[out=46,in=-46] (57.5mm, 72.5mm);
}

% Third sound wave (Curved line)
\opt{w3}{
    \draw[line width=4.3mm,line cap=round]
        (72mm, 18.5mm) to[out=46,in=-46] (72mm, 81.5mm);
}

% Always draw a transparent copy of the third wave so the images will all have
% the same width.
\draw[line width=4.3mm,line cap=round,opacity=0]
    (72mm, 18.5mm) to[out=46,in=-46] (72mm, 81.5mm);

\end{tikzpicture}
\end{document}

'
%   xelatex -jobname=Volume2 '\def\UseOption{w1,w2}% This file is part of Background Music.
%
% Background Music is free software: you can redistribute it and/or
% modify it under the terms of the GNU General Public License as
% published by the Free Software Foundation, either version 2 of the
% License, or (at your option) any later version.
%
% Background Music is distributed in the hope that it will be useful,
% but WITHOUT ANY WARRANTY; without even the implied warranty of
% MERCHANTABILITY or FITNESS FOR A PARTICULAR PURPOSE. See the
% GNU General Public License for more details.
%
% You should have received a copy of the GNU General Public License
% along with Background Music. If not, see <http://www.gnu.org/licenses/>.

%
% VolumeIcons.tex
%
% Build with XeTeX:
%   xelatex -jobname=Volume0 '\def\UseOption{}\input{VolumeIcons.tex}'
%   xelatex -jobname=Volume1 '\def\UseOption{w1}\input{VolumeIcons.tex}'
%   xelatex -jobname=Volume2 '\def\UseOption{w1,w2}\input{VolumeIcons.tex}'
%   xelatex -jobname=Volume3 '\def\UseOption{w1,w2,w3}\input{VolumeIcons.tex}'
%   for n in 0 1 2 3; do mv Volume$n.pdf ../BGMApp/BGMApp/Images.xcassets/Volume$n.imageset/; done
%
% Might build correctly with regular LaTeX. I haven't tried it.
%

\documentclass[tikz]{standalone}
\usepackage{tikz}
% "dummyOption" prevents "Package optional Warning: No options were selected,
% so all optional text will be printed" when building Volume0.pdf.
\usepackage[dummyOption]{optional}

\begin{document}
\begin{tikzpicture}

% Speaker (Rounded box and triangle)
\fill[rounded corners=5mm]
    (0mm, 62.5mm) rectangle (25mm, 37.5mm) {};
\draw[rounded corners=2.5mm,fill=black]
    (3mm, 50mm)--(34mm, 76.5mm)--(34mm, 23.5mm)--cycle;

% First sound wave (Curved line)
\opt{w1}{
    \draw[line width=4.3mm,line cap=round]
        (44mm, 36.5mm) to[out=46,in=-46] (44mm, 63.5mm);
}

% Second sound wave (Curved line)
\opt{w2}{
    \draw[line width=4.3mm,line cap=round]
        (57.5mm, 27.5mm) to[out=46,in=-46] (57.5mm, 72.5mm);
}

% Third sound wave (Curved line)
\opt{w3}{
    \draw[line width=4.3mm,line cap=round]
        (72mm, 18.5mm) to[out=46,in=-46] (72mm, 81.5mm);
}

% Always draw a transparent copy of the third wave so the images will all have
% the same width.
\draw[line width=4.3mm,line cap=round,opacity=0]
    (72mm, 18.5mm) to[out=46,in=-46] (72mm, 81.5mm);

\end{tikzpicture}
\end{document}

'
%   xelatex -jobname=Volume3 '\def\UseOption{w1,w2,w3}% This file is part of Background Music.
%
% Background Music is free software: you can redistribute it and/or
% modify it under the terms of the GNU General Public License as
% published by the Free Software Foundation, either version 2 of the
% License, or (at your option) any later version.
%
% Background Music is distributed in the hope that it will be useful,
% but WITHOUT ANY WARRANTY; without even the implied warranty of
% MERCHANTABILITY or FITNESS FOR A PARTICULAR PURPOSE. See the
% GNU General Public License for more details.
%
% You should have received a copy of the GNU General Public License
% along with Background Music. If not, see <http://www.gnu.org/licenses/>.

%
% VolumeIcons.tex
%
% Build with XeTeX:
%   xelatex -jobname=Volume0 '\def\UseOption{}\input{VolumeIcons.tex}'
%   xelatex -jobname=Volume1 '\def\UseOption{w1}\input{VolumeIcons.tex}'
%   xelatex -jobname=Volume2 '\def\UseOption{w1,w2}\input{VolumeIcons.tex}'
%   xelatex -jobname=Volume3 '\def\UseOption{w1,w2,w3}\input{VolumeIcons.tex}'
%   for n in 0 1 2 3; do mv Volume$n.pdf ../BGMApp/BGMApp/Images.xcassets/Volume$n.imageset/; done
%
% Might build correctly with regular LaTeX. I haven't tried it.
%

\documentclass[tikz]{standalone}
\usepackage{tikz}
% "dummyOption" prevents "Package optional Warning: No options were selected,
% so all optional text will be printed" when building Volume0.pdf.
\usepackage[dummyOption]{optional}

\begin{document}
\begin{tikzpicture}

% Speaker (Rounded box and triangle)
\fill[rounded corners=5mm]
    (0mm, 62.5mm) rectangle (25mm, 37.5mm) {};
\draw[rounded corners=2.5mm,fill=black]
    (3mm, 50mm)--(34mm, 76.5mm)--(34mm, 23.5mm)--cycle;

% First sound wave (Curved line)
\opt{w1}{
    \draw[line width=4.3mm,line cap=round]
        (44mm, 36.5mm) to[out=46,in=-46] (44mm, 63.5mm);
}

% Second sound wave (Curved line)
\opt{w2}{
    \draw[line width=4.3mm,line cap=round]
        (57.5mm, 27.5mm) to[out=46,in=-46] (57.5mm, 72.5mm);
}

% Third sound wave (Curved line)
\opt{w3}{
    \draw[line width=4.3mm,line cap=round]
        (72mm, 18.5mm) to[out=46,in=-46] (72mm, 81.5mm);
}

% Always draw a transparent copy of the third wave so the images will all have
% the same width.
\draw[line width=4.3mm,line cap=round,opacity=0]
    (72mm, 18.5mm) to[out=46,in=-46] (72mm, 81.5mm);

\end{tikzpicture}
\end{document}

'
%   for n in 0 1 2 3; do mv Volume$n.pdf ../BGMApp/BGMApp/Images.xcassets/Volume$n.imageset/; done
%
% Might build correctly with regular LaTeX. I haven't tried it.
%

\documentclass[tikz]{standalone}
\usepackage{tikz}
% "dummyOption" prevents "Package optional Warning: No options were selected,
% so all optional text will be printed" when building Volume0.pdf.
\usepackage[dummyOption]{optional}

\begin{document}
\begin{tikzpicture}

% Speaker (Rounded box and triangle)
\fill[rounded corners=5mm]
    (0mm, 62.5mm) rectangle (25mm, 37.5mm) {};
\draw[rounded corners=2.5mm,fill=black]
    (3mm, 50mm)--(34mm, 76.5mm)--(34mm, 23.5mm)--cycle;

% First sound wave (Curved line)
\opt{w1}{
    \draw[line width=4.3mm,line cap=round]
        (44mm, 36.5mm) to[out=46,in=-46] (44mm, 63.5mm);
}

% Second sound wave (Curved line)
\opt{w2}{
    \draw[line width=4.3mm,line cap=round]
        (57.5mm, 27.5mm) to[out=46,in=-46] (57.5mm, 72.5mm);
}

% Third sound wave (Curved line)
\opt{w3}{
    \draw[line width=4.3mm,line cap=round]
        (72mm, 18.5mm) to[out=46,in=-46] (72mm, 81.5mm);
}

% Always draw a transparent copy of the third wave so the images will all have
% the same width.
\draw[line width=4.3mm,line cap=round,opacity=0]
    (72mm, 18.5mm) to[out=46,in=-46] (72mm, 81.5mm);

\end{tikzpicture}
\end{document}

'
%   for n in 0 1 2 3; do mv Volume$n.pdf ../BGMApp/BGMApp/Images.xcassets/Volume$n.imageset/; done
%
% Might build correctly with regular LaTeX. I haven't tried it.
%

\documentclass[tikz]{standalone}
\usepackage{tikz}
% "dummyOption" prevents "Package optional Warning: No options were selected,
% so all optional text will be printed" when building Volume0.pdf.
\usepackage[dummyOption]{optional}

\begin{document}
\begin{tikzpicture}

% Speaker (Rounded box and triangle)
\fill[rounded corners=5mm]
    (0mm, 62.5mm) rectangle (25mm, 37.5mm) {};
\draw[rounded corners=2.5mm,fill=black]
    (3mm, 50mm)--(34mm, 76.5mm)--(34mm, 23.5mm)--cycle;

% First sound wave (Curved line)
\opt{w1}{
    \draw[line width=4.3mm,line cap=round]
        (44mm, 36.5mm) to[out=46,in=-46] (44mm, 63.5mm);
}

% Second sound wave (Curved line)
\opt{w2}{
    \draw[line width=4.3mm,line cap=round]
        (57.5mm, 27.5mm) to[out=46,in=-46] (57.5mm, 72.5mm);
}

% Third sound wave (Curved line)
\opt{w3}{
    \draw[line width=4.3mm,line cap=round]
        (72mm, 18.5mm) to[out=46,in=-46] (72mm, 81.5mm);
}

% Always draw a transparent copy of the third wave so the images will all have
% the same width.
\draw[line width=4.3mm,line cap=round,opacity=0]
    (72mm, 18.5mm) to[out=46,in=-46] (72mm, 81.5mm);

\end{tikzpicture}
\end{document}

'
%   xelatex -jobname=Volume1 '\def\UseOption{w1}% This file is part of Background Music.
%
% Background Music is free software: you can redistribute it and/or
% modify it under the terms of the GNU General Public License as
% published by the Free Software Foundation, either version 2 of the
% License, or (at your option) any later version.
%
% Background Music is distributed in the hope that it will be useful,
% but WITHOUT ANY WARRANTY; without even the implied warranty of
% MERCHANTABILITY or FITNESS FOR A PARTICULAR PURPOSE. See the
% GNU General Public License for more details.
%
% You should have received a copy of the GNU General Public License
% along with Background Music. If not, see <http://www.gnu.org/licenses/>.

%
% VolumeIcons.tex
%
% Build with XeTeX:
%   xelatex -jobname=Volume0 '\def\UseOption{}% This file is part of Background Music.
%
% Background Music is free software: you can redistribute it and/or
% modify it under the terms of the GNU General Public License as
% published by the Free Software Foundation, either version 2 of the
% License, or (at your option) any later version.
%
% Background Music is distributed in the hope that it will be useful,
% but WITHOUT ANY WARRANTY; without even the implied warranty of
% MERCHANTABILITY or FITNESS FOR A PARTICULAR PURPOSE. See the
% GNU General Public License for more details.
%
% You should have received a copy of the GNU General Public License
% along with Background Music. If not, see <http://www.gnu.org/licenses/>.

%
% VolumeIcons.tex
%
% Build with XeTeX:
%   xelatex -jobname=Volume0 '\def\UseOption{}% This file is part of Background Music.
%
% Background Music is free software: you can redistribute it and/or
% modify it under the terms of the GNU General Public License as
% published by the Free Software Foundation, either version 2 of the
% License, or (at your option) any later version.
%
% Background Music is distributed in the hope that it will be useful,
% but WITHOUT ANY WARRANTY; without even the implied warranty of
% MERCHANTABILITY or FITNESS FOR A PARTICULAR PURPOSE. See the
% GNU General Public License for more details.
%
% You should have received a copy of the GNU General Public License
% along with Background Music. If not, see <http://www.gnu.org/licenses/>.

%
% VolumeIcons.tex
%
% Build with XeTeX:
%   xelatex -jobname=Volume0 '\def\UseOption{}\input{VolumeIcons.tex}'
%   xelatex -jobname=Volume1 '\def\UseOption{w1}\input{VolumeIcons.tex}'
%   xelatex -jobname=Volume2 '\def\UseOption{w1,w2}\input{VolumeIcons.tex}'
%   xelatex -jobname=Volume3 '\def\UseOption{w1,w2,w3}\input{VolumeIcons.tex}'
%   for n in 0 1 2 3; do mv Volume$n.pdf ../BGMApp/BGMApp/Images.xcassets/Volume$n.imageset/; done
%
% Might build correctly with regular LaTeX. I haven't tried it.
%

\documentclass[tikz]{standalone}
\usepackage{tikz}
% "dummyOption" prevents "Package optional Warning: No options were selected,
% so all optional text will be printed" when building Volume0.pdf.
\usepackage[dummyOption]{optional}

\begin{document}
\begin{tikzpicture}

% Speaker (Rounded box and triangle)
\fill[rounded corners=5mm]
    (0mm, 62.5mm) rectangle (25mm, 37.5mm) {};
\draw[rounded corners=2.5mm,fill=black]
    (3mm, 50mm)--(34mm, 76.5mm)--(34mm, 23.5mm)--cycle;

% First sound wave (Curved line)
\opt{w1}{
    \draw[line width=4.3mm,line cap=round]
        (44mm, 36.5mm) to[out=46,in=-46] (44mm, 63.5mm);
}

% Second sound wave (Curved line)
\opt{w2}{
    \draw[line width=4.3mm,line cap=round]
        (57.5mm, 27.5mm) to[out=46,in=-46] (57.5mm, 72.5mm);
}

% Third sound wave (Curved line)
\opt{w3}{
    \draw[line width=4.3mm,line cap=round]
        (72mm, 18.5mm) to[out=46,in=-46] (72mm, 81.5mm);
}

% Always draw a transparent copy of the third wave so the images will all have
% the same width.
\draw[line width=4.3mm,line cap=round,opacity=0]
    (72mm, 18.5mm) to[out=46,in=-46] (72mm, 81.5mm);

\end{tikzpicture}
\end{document}

'
%   xelatex -jobname=Volume1 '\def\UseOption{w1}% This file is part of Background Music.
%
% Background Music is free software: you can redistribute it and/or
% modify it under the terms of the GNU General Public License as
% published by the Free Software Foundation, either version 2 of the
% License, or (at your option) any later version.
%
% Background Music is distributed in the hope that it will be useful,
% but WITHOUT ANY WARRANTY; without even the implied warranty of
% MERCHANTABILITY or FITNESS FOR A PARTICULAR PURPOSE. See the
% GNU General Public License for more details.
%
% You should have received a copy of the GNU General Public License
% along with Background Music. If not, see <http://www.gnu.org/licenses/>.

%
% VolumeIcons.tex
%
% Build with XeTeX:
%   xelatex -jobname=Volume0 '\def\UseOption{}\input{VolumeIcons.tex}'
%   xelatex -jobname=Volume1 '\def\UseOption{w1}\input{VolumeIcons.tex}'
%   xelatex -jobname=Volume2 '\def\UseOption{w1,w2}\input{VolumeIcons.tex}'
%   xelatex -jobname=Volume3 '\def\UseOption{w1,w2,w3}\input{VolumeIcons.tex}'
%   for n in 0 1 2 3; do mv Volume$n.pdf ../BGMApp/BGMApp/Images.xcassets/Volume$n.imageset/; done
%
% Might build correctly with regular LaTeX. I haven't tried it.
%

\documentclass[tikz]{standalone}
\usepackage{tikz}
% "dummyOption" prevents "Package optional Warning: No options were selected,
% so all optional text will be printed" when building Volume0.pdf.
\usepackage[dummyOption]{optional}

\begin{document}
\begin{tikzpicture}

% Speaker (Rounded box and triangle)
\fill[rounded corners=5mm]
    (0mm, 62.5mm) rectangle (25mm, 37.5mm) {};
\draw[rounded corners=2.5mm,fill=black]
    (3mm, 50mm)--(34mm, 76.5mm)--(34mm, 23.5mm)--cycle;

% First sound wave (Curved line)
\opt{w1}{
    \draw[line width=4.3mm,line cap=round]
        (44mm, 36.5mm) to[out=46,in=-46] (44mm, 63.5mm);
}

% Second sound wave (Curved line)
\opt{w2}{
    \draw[line width=4.3mm,line cap=round]
        (57.5mm, 27.5mm) to[out=46,in=-46] (57.5mm, 72.5mm);
}

% Third sound wave (Curved line)
\opt{w3}{
    \draw[line width=4.3mm,line cap=round]
        (72mm, 18.5mm) to[out=46,in=-46] (72mm, 81.5mm);
}

% Always draw a transparent copy of the third wave so the images will all have
% the same width.
\draw[line width=4.3mm,line cap=round,opacity=0]
    (72mm, 18.5mm) to[out=46,in=-46] (72mm, 81.5mm);

\end{tikzpicture}
\end{document}

'
%   xelatex -jobname=Volume2 '\def\UseOption{w1,w2}% This file is part of Background Music.
%
% Background Music is free software: you can redistribute it and/or
% modify it under the terms of the GNU General Public License as
% published by the Free Software Foundation, either version 2 of the
% License, or (at your option) any later version.
%
% Background Music is distributed in the hope that it will be useful,
% but WITHOUT ANY WARRANTY; without even the implied warranty of
% MERCHANTABILITY or FITNESS FOR A PARTICULAR PURPOSE. See the
% GNU General Public License for more details.
%
% You should have received a copy of the GNU General Public License
% along with Background Music. If not, see <http://www.gnu.org/licenses/>.

%
% VolumeIcons.tex
%
% Build with XeTeX:
%   xelatex -jobname=Volume0 '\def\UseOption{}\input{VolumeIcons.tex}'
%   xelatex -jobname=Volume1 '\def\UseOption{w1}\input{VolumeIcons.tex}'
%   xelatex -jobname=Volume2 '\def\UseOption{w1,w2}\input{VolumeIcons.tex}'
%   xelatex -jobname=Volume3 '\def\UseOption{w1,w2,w3}\input{VolumeIcons.tex}'
%   for n in 0 1 2 3; do mv Volume$n.pdf ../BGMApp/BGMApp/Images.xcassets/Volume$n.imageset/; done
%
% Might build correctly with regular LaTeX. I haven't tried it.
%

\documentclass[tikz]{standalone}
\usepackage{tikz}
% "dummyOption" prevents "Package optional Warning: No options were selected,
% so all optional text will be printed" when building Volume0.pdf.
\usepackage[dummyOption]{optional}

\begin{document}
\begin{tikzpicture}

% Speaker (Rounded box and triangle)
\fill[rounded corners=5mm]
    (0mm, 62.5mm) rectangle (25mm, 37.5mm) {};
\draw[rounded corners=2.5mm,fill=black]
    (3mm, 50mm)--(34mm, 76.5mm)--(34mm, 23.5mm)--cycle;

% First sound wave (Curved line)
\opt{w1}{
    \draw[line width=4.3mm,line cap=round]
        (44mm, 36.5mm) to[out=46,in=-46] (44mm, 63.5mm);
}

% Second sound wave (Curved line)
\opt{w2}{
    \draw[line width=4.3mm,line cap=round]
        (57.5mm, 27.5mm) to[out=46,in=-46] (57.5mm, 72.5mm);
}

% Third sound wave (Curved line)
\opt{w3}{
    \draw[line width=4.3mm,line cap=round]
        (72mm, 18.5mm) to[out=46,in=-46] (72mm, 81.5mm);
}

% Always draw a transparent copy of the third wave so the images will all have
% the same width.
\draw[line width=4.3mm,line cap=round,opacity=0]
    (72mm, 18.5mm) to[out=46,in=-46] (72mm, 81.5mm);

\end{tikzpicture}
\end{document}

'
%   xelatex -jobname=Volume3 '\def\UseOption{w1,w2,w3}% This file is part of Background Music.
%
% Background Music is free software: you can redistribute it and/or
% modify it under the terms of the GNU General Public License as
% published by the Free Software Foundation, either version 2 of the
% License, or (at your option) any later version.
%
% Background Music is distributed in the hope that it will be useful,
% but WITHOUT ANY WARRANTY; without even the implied warranty of
% MERCHANTABILITY or FITNESS FOR A PARTICULAR PURPOSE. See the
% GNU General Public License for more details.
%
% You should have received a copy of the GNU General Public License
% along with Background Music. If not, see <http://www.gnu.org/licenses/>.

%
% VolumeIcons.tex
%
% Build with XeTeX:
%   xelatex -jobname=Volume0 '\def\UseOption{}\input{VolumeIcons.tex}'
%   xelatex -jobname=Volume1 '\def\UseOption{w1}\input{VolumeIcons.tex}'
%   xelatex -jobname=Volume2 '\def\UseOption{w1,w2}\input{VolumeIcons.tex}'
%   xelatex -jobname=Volume3 '\def\UseOption{w1,w2,w3}\input{VolumeIcons.tex}'
%   for n in 0 1 2 3; do mv Volume$n.pdf ../BGMApp/BGMApp/Images.xcassets/Volume$n.imageset/; done
%
% Might build correctly with regular LaTeX. I haven't tried it.
%

\documentclass[tikz]{standalone}
\usepackage{tikz}
% "dummyOption" prevents "Package optional Warning: No options were selected,
% so all optional text will be printed" when building Volume0.pdf.
\usepackage[dummyOption]{optional}

\begin{document}
\begin{tikzpicture}

% Speaker (Rounded box and triangle)
\fill[rounded corners=5mm]
    (0mm, 62.5mm) rectangle (25mm, 37.5mm) {};
\draw[rounded corners=2.5mm,fill=black]
    (3mm, 50mm)--(34mm, 76.5mm)--(34mm, 23.5mm)--cycle;

% First sound wave (Curved line)
\opt{w1}{
    \draw[line width=4.3mm,line cap=round]
        (44mm, 36.5mm) to[out=46,in=-46] (44mm, 63.5mm);
}

% Second sound wave (Curved line)
\opt{w2}{
    \draw[line width=4.3mm,line cap=round]
        (57.5mm, 27.5mm) to[out=46,in=-46] (57.5mm, 72.5mm);
}

% Third sound wave (Curved line)
\opt{w3}{
    \draw[line width=4.3mm,line cap=round]
        (72mm, 18.5mm) to[out=46,in=-46] (72mm, 81.5mm);
}

% Always draw a transparent copy of the third wave so the images will all have
% the same width.
\draw[line width=4.3mm,line cap=round,opacity=0]
    (72mm, 18.5mm) to[out=46,in=-46] (72mm, 81.5mm);

\end{tikzpicture}
\end{document}

'
%   for n in 0 1 2 3; do mv Volume$n.pdf ../BGMApp/BGMApp/Images.xcassets/Volume$n.imageset/; done
%
% Might build correctly with regular LaTeX. I haven't tried it.
%

\documentclass[tikz]{standalone}
\usepackage{tikz}
% "dummyOption" prevents "Package optional Warning: No options were selected,
% so all optional text will be printed" when building Volume0.pdf.
\usepackage[dummyOption]{optional}

\begin{document}
\begin{tikzpicture}

% Speaker (Rounded box and triangle)
\fill[rounded corners=5mm]
    (0mm, 62.5mm) rectangle (25mm, 37.5mm) {};
\draw[rounded corners=2.5mm,fill=black]
    (3mm, 50mm)--(34mm, 76.5mm)--(34mm, 23.5mm)--cycle;

% First sound wave (Curved line)
\opt{w1}{
    \draw[line width=4.3mm,line cap=round]
        (44mm, 36.5mm) to[out=46,in=-46] (44mm, 63.5mm);
}

% Second sound wave (Curved line)
\opt{w2}{
    \draw[line width=4.3mm,line cap=round]
        (57.5mm, 27.5mm) to[out=46,in=-46] (57.5mm, 72.5mm);
}

% Third sound wave (Curved line)
\opt{w3}{
    \draw[line width=4.3mm,line cap=round]
        (72mm, 18.5mm) to[out=46,in=-46] (72mm, 81.5mm);
}

% Always draw a transparent copy of the third wave so the images will all have
% the same width.
\draw[line width=4.3mm,line cap=round,opacity=0]
    (72mm, 18.5mm) to[out=46,in=-46] (72mm, 81.5mm);

\end{tikzpicture}
\end{document}

'
%   xelatex -jobname=Volume1 '\def\UseOption{w1}% This file is part of Background Music.
%
% Background Music is free software: you can redistribute it and/or
% modify it under the terms of the GNU General Public License as
% published by the Free Software Foundation, either version 2 of the
% License, or (at your option) any later version.
%
% Background Music is distributed in the hope that it will be useful,
% but WITHOUT ANY WARRANTY; without even the implied warranty of
% MERCHANTABILITY or FITNESS FOR A PARTICULAR PURPOSE. See the
% GNU General Public License for more details.
%
% You should have received a copy of the GNU General Public License
% along with Background Music. If not, see <http://www.gnu.org/licenses/>.

%
% VolumeIcons.tex
%
% Build with XeTeX:
%   xelatex -jobname=Volume0 '\def\UseOption{}% This file is part of Background Music.
%
% Background Music is free software: you can redistribute it and/or
% modify it under the terms of the GNU General Public License as
% published by the Free Software Foundation, either version 2 of the
% License, or (at your option) any later version.
%
% Background Music is distributed in the hope that it will be useful,
% but WITHOUT ANY WARRANTY; without even the implied warranty of
% MERCHANTABILITY or FITNESS FOR A PARTICULAR PURPOSE. See the
% GNU General Public License for more details.
%
% You should have received a copy of the GNU General Public License
% along with Background Music. If not, see <http://www.gnu.org/licenses/>.

%
% VolumeIcons.tex
%
% Build with XeTeX:
%   xelatex -jobname=Volume0 '\def\UseOption{}\input{VolumeIcons.tex}'
%   xelatex -jobname=Volume1 '\def\UseOption{w1}\input{VolumeIcons.tex}'
%   xelatex -jobname=Volume2 '\def\UseOption{w1,w2}\input{VolumeIcons.tex}'
%   xelatex -jobname=Volume3 '\def\UseOption{w1,w2,w3}\input{VolumeIcons.tex}'
%   for n in 0 1 2 3; do mv Volume$n.pdf ../BGMApp/BGMApp/Images.xcassets/Volume$n.imageset/; done
%
% Might build correctly with regular LaTeX. I haven't tried it.
%

\documentclass[tikz]{standalone}
\usepackage{tikz}
% "dummyOption" prevents "Package optional Warning: No options were selected,
% so all optional text will be printed" when building Volume0.pdf.
\usepackage[dummyOption]{optional}

\begin{document}
\begin{tikzpicture}

% Speaker (Rounded box and triangle)
\fill[rounded corners=5mm]
    (0mm, 62.5mm) rectangle (25mm, 37.5mm) {};
\draw[rounded corners=2.5mm,fill=black]
    (3mm, 50mm)--(34mm, 76.5mm)--(34mm, 23.5mm)--cycle;

% First sound wave (Curved line)
\opt{w1}{
    \draw[line width=4.3mm,line cap=round]
        (44mm, 36.5mm) to[out=46,in=-46] (44mm, 63.5mm);
}

% Second sound wave (Curved line)
\opt{w2}{
    \draw[line width=4.3mm,line cap=round]
        (57.5mm, 27.5mm) to[out=46,in=-46] (57.5mm, 72.5mm);
}

% Third sound wave (Curved line)
\opt{w3}{
    \draw[line width=4.3mm,line cap=round]
        (72mm, 18.5mm) to[out=46,in=-46] (72mm, 81.5mm);
}

% Always draw a transparent copy of the third wave so the images will all have
% the same width.
\draw[line width=4.3mm,line cap=round,opacity=0]
    (72mm, 18.5mm) to[out=46,in=-46] (72mm, 81.5mm);

\end{tikzpicture}
\end{document}

'
%   xelatex -jobname=Volume1 '\def\UseOption{w1}% This file is part of Background Music.
%
% Background Music is free software: you can redistribute it and/or
% modify it under the terms of the GNU General Public License as
% published by the Free Software Foundation, either version 2 of the
% License, or (at your option) any later version.
%
% Background Music is distributed in the hope that it will be useful,
% but WITHOUT ANY WARRANTY; without even the implied warranty of
% MERCHANTABILITY or FITNESS FOR A PARTICULAR PURPOSE. See the
% GNU General Public License for more details.
%
% You should have received a copy of the GNU General Public License
% along with Background Music. If not, see <http://www.gnu.org/licenses/>.

%
% VolumeIcons.tex
%
% Build with XeTeX:
%   xelatex -jobname=Volume0 '\def\UseOption{}\input{VolumeIcons.tex}'
%   xelatex -jobname=Volume1 '\def\UseOption{w1}\input{VolumeIcons.tex}'
%   xelatex -jobname=Volume2 '\def\UseOption{w1,w2}\input{VolumeIcons.tex}'
%   xelatex -jobname=Volume3 '\def\UseOption{w1,w2,w3}\input{VolumeIcons.tex}'
%   for n in 0 1 2 3; do mv Volume$n.pdf ../BGMApp/BGMApp/Images.xcassets/Volume$n.imageset/; done
%
% Might build correctly with regular LaTeX. I haven't tried it.
%

\documentclass[tikz]{standalone}
\usepackage{tikz}
% "dummyOption" prevents "Package optional Warning: No options were selected,
% so all optional text will be printed" when building Volume0.pdf.
\usepackage[dummyOption]{optional}

\begin{document}
\begin{tikzpicture}

% Speaker (Rounded box and triangle)
\fill[rounded corners=5mm]
    (0mm, 62.5mm) rectangle (25mm, 37.5mm) {};
\draw[rounded corners=2.5mm,fill=black]
    (3mm, 50mm)--(34mm, 76.5mm)--(34mm, 23.5mm)--cycle;

% First sound wave (Curved line)
\opt{w1}{
    \draw[line width=4.3mm,line cap=round]
        (44mm, 36.5mm) to[out=46,in=-46] (44mm, 63.5mm);
}

% Second sound wave (Curved line)
\opt{w2}{
    \draw[line width=4.3mm,line cap=round]
        (57.5mm, 27.5mm) to[out=46,in=-46] (57.5mm, 72.5mm);
}

% Third sound wave (Curved line)
\opt{w3}{
    \draw[line width=4.3mm,line cap=round]
        (72mm, 18.5mm) to[out=46,in=-46] (72mm, 81.5mm);
}

% Always draw a transparent copy of the third wave so the images will all have
% the same width.
\draw[line width=4.3mm,line cap=round,opacity=0]
    (72mm, 18.5mm) to[out=46,in=-46] (72mm, 81.5mm);

\end{tikzpicture}
\end{document}

'
%   xelatex -jobname=Volume2 '\def\UseOption{w1,w2}% This file is part of Background Music.
%
% Background Music is free software: you can redistribute it and/or
% modify it under the terms of the GNU General Public License as
% published by the Free Software Foundation, either version 2 of the
% License, or (at your option) any later version.
%
% Background Music is distributed in the hope that it will be useful,
% but WITHOUT ANY WARRANTY; without even the implied warranty of
% MERCHANTABILITY or FITNESS FOR A PARTICULAR PURPOSE. See the
% GNU General Public License for more details.
%
% You should have received a copy of the GNU General Public License
% along with Background Music. If not, see <http://www.gnu.org/licenses/>.

%
% VolumeIcons.tex
%
% Build with XeTeX:
%   xelatex -jobname=Volume0 '\def\UseOption{}\input{VolumeIcons.tex}'
%   xelatex -jobname=Volume1 '\def\UseOption{w1}\input{VolumeIcons.tex}'
%   xelatex -jobname=Volume2 '\def\UseOption{w1,w2}\input{VolumeIcons.tex}'
%   xelatex -jobname=Volume3 '\def\UseOption{w1,w2,w3}\input{VolumeIcons.tex}'
%   for n in 0 1 2 3; do mv Volume$n.pdf ../BGMApp/BGMApp/Images.xcassets/Volume$n.imageset/; done
%
% Might build correctly with regular LaTeX. I haven't tried it.
%

\documentclass[tikz]{standalone}
\usepackage{tikz}
% "dummyOption" prevents "Package optional Warning: No options were selected,
% so all optional text will be printed" when building Volume0.pdf.
\usepackage[dummyOption]{optional}

\begin{document}
\begin{tikzpicture}

% Speaker (Rounded box and triangle)
\fill[rounded corners=5mm]
    (0mm, 62.5mm) rectangle (25mm, 37.5mm) {};
\draw[rounded corners=2.5mm,fill=black]
    (3mm, 50mm)--(34mm, 76.5mm)--(34mm, 23.5mm)--cycle;

% First sound wave (Curved line)
\opt{w1}{
    \draw[line width=4.3mm,line cap=round]
        (44mm, 36.5mm) to[out=46,in=-46] (44mm, 63.5mm);
}

% Second sound wave (Curved line)
\opt{w2}{
    \draw[line width=4.3mm,line cap=round]
        (57.5mm, 27.5mm) to[out=46,in=-46] (57.5mm, 72.5mm);
}

% Third sound wave (Curved line)
\opt{w3}{
    \draw[line width=4.3mm,line cap=round]
        (72mm, 18.5mm) to[out=46,in=-46] (72mm, 81.5mm);
}

% Always draw a transparent copy of the third wave so the images will all have
% the same width.
\draw[line width=4.3mm,line cap=round,opacity=0]
    (72mm, 18.5mm) to[out=46,in=-46] (72mm, 81.5mm);

\end{tikzpicture}
\end{document}

'
%   xelatex -jobname=Volume3 '\def\UseOption{w1,w2,w3}% This file is part of Background Music.
%
% Background Music is free software: you can redistribute it and/or
% modify it under the terms of the GNU General Public License as
% published by the Free Software Foundation, either version 2 of the
% License, or (at your option) any later version.
%
% Background Music is distributed in the hope that it will be useful,
% but WITHOUT ANY WARRANTY; without even the implied warranty of
% MERCHANTABILITY or FITNESS FOR A PARTICULAR PURPOSE. See the
% GNU General Public License for more details.
%
% You should have received a copy of the GNU General Public License
% along with Background Music. If not, see <http://www.gnu.org/licenses/>.

%
% VolumeIcons.tex
%
% Build with XeTeX:
%   xelatex -jobname=Volume0 '\def\UseOption{}\input{VolumeIcons.tex}'
%   xelatex -jobname=Volume1 '\def\UseOption{w1}\input{VolumeIcons.tex}'
%   xelatex -jobname=Volume2 '\def\UseOption{w1,w2}\input{VolumeIcons.tex}'
%   xelatex -jobname=Volume3 '\def\UseOption{w1,w2,w3}\input{VolumeIcons.tex}'
%   for n in 0 1 2 3; do mv Volume$n.pdf ../BGMApp/BGMApp/Images.xcassets/Volume$n.imageset/; done
%
% Might build correctly with regular LaTeX. I haven't tried it.
%

\documentclass[tikz]{standalone}
\usepackage{tikz}
% "dummyOption" prevents "Package optional Warning: No options were selected,
% so all optional text will be printed" when building Volume0.pdf.
\usepackage[dummyOption]{optional}

\begin{document}
\begin{tikzpicture}

% Speaker (Rounded box and triangle)
\fill[rounded corners=5mm]
    (0mm, 62.5mm) rectangle (25mm, 37.5mm) {};
\draw[rounded corners=2.5mm,fill=black]
    (3mm, 50mm)--(34mm, 76.5mm)--(34mm, 23.5mm)--cycle;

% First sound wave (Curved line)
\opt{w1}{
    \draw[line width=4.3mm,line cap=round]
        (44mm, 36.5mm) to[out=46,in=-46] (44mm, 63.5mm);
}

% Second sound wave (Curved line)
\opt{w2}{
    \draw[line width=4.3mm,line cap=round]
        (57.5mm, 27.5mm) to[out=46,in=-46] (57.5mm, 72.5mm);
}

% Third sound wave (Curved line)
\opt{w3}{
    \draw[line width=4.3mm,line cap=round]
        (72mm, 18.5mm) to[out=46,in=-46] (72mm, 81.5mm);
}

% Always draw a transparent copy of the third wave so the images will all have
% the same width.
\draw[line width=4.3mm,line cap=round,opacity=0]
    (72mm, 18.5mm) to[out=46,in=-46] (72mm, 81.5mm);

\end{tikzpicture}
\end{document}

'
%   for n in 0 1 2 3; do mv Volume$n.pdf ../BGMApp/BGMApp/Images.xcassets/Volume$n.imageset/; done
%
% Might build correctly with regular LaTeX. I haven't tried it.
%

\documentclass[tikz]{standalone}
\usepackage{tikz}
% "dummyOption" prevents "Package optional Warning: No options were selected,
% so all optional text will be printed" when building Volume0.pdf.
\usepackage[dummyOption]{optional}

\begin{document}
\begin{tikzpicture}

% Speaker (Rounded box and triangle)
\fill[rounded corners=5mm]
    (0mm, 62.5mm) rectangle (25mm, 37.5mm) {};
\draw[rounded corners=2.5mm,fill=black]
    (3mm, 50mm)--(34mm, 76.5mm)--(34mm, 23.5mm)--cycle;

% First sound wave (Curved line)
\opt{w1}{
    \draw[line width=4.3mm,line cap=round]
        (44mm, 36.5mm) to[out=46,in=-46] (44mm, 63.5mm);
}

% Second sound wave (Curved line)
\opt{w2}{
    \draw[line width=4.3mm,line cap=round]
        (57.5mm, 27.5mm) to[out=46,in=-46] (57.5mm, 72.5mm);
}

% Third sound wave (Curved line)
\opt{w3}{
    \draw[line width=4.3mm,line cap=round]
        (72mm, 18.5mm) to[out=46,in=-46] (72mm, 81.5mm);
}

% Always draw a transparent copy of the third wave so the images will all have
% the same width.
\draw[line width=4.3mm,line cap=round,opacity=0]
    (72mm, 18.5mm) to[out=46,in=-46] (72mm, 81.5mm);

\end{tikzpicture}
\end{document}

'
%   xelatex -jobname=Volume2 '\def\UseOption{w1,w2}% This file is part of Background Music.
%
% Background Music is free software: you can redistribute it and/or
% modify it under the terms of the GNU General Public License as
% published by the Free Software Foundation, either version 2 of the
% License, or (at your option) any later version.
%
% Background Music is distributed in the hope that it will be useful,
% but WITHOUT ANY WARRANTY; without even the implied warranty of
% MERCHANTABILITY or FITNESS FOR A PARTICULAR PURPOSE. See the
% GNU General Public License for more details.
%
% You should have received a copy of the GNU General Public License
% along with Background Music. If not, see <http://www.gnu.org/licenses/>.

%
% VolumeIcons.tex
%
% Build with XeTeX:
%   xelatex -jobname=Volume0 '\def\UseOption{}% This file is part of Background Music.
%
% Background Music is free software: you can redistribute it and/or
% modify it under the terms of the GNU General Public License as
% published by the Free Software Foundation, either version 2 of the
% License, or (at your option) any later version.
%
% Background Music is distributed in the hope that it will be useful,
% but WITHOUT ANY WARRANTY; without even the implied warranty of
% MERCHANTABILITY or FITNESS FOR A PARTICULAR PURPOSE. See the
% GNU General Public License for more details.
%
% You should have received a copy of the GNU General Public License
% along with Background Music. If not, see <http://www.gnu.org/licenses/>.

%
% VolumeIcons.tex
%
% Build with XeTeX:
%   xelatex -jobname=Volume0 '\def\UseOption{}\input{VolumeIcons.tex}'
%   xelatex -jobname=Volume1 '\def\UseOption{w1}\input{VolumeIcons.tex}'
%   xelatex -jobname=Volume2 '\def\UseOption{w1,w2}\input{VolumeIcons.tex}'
%   xelatex -jobname=Volume3 '\def\UseOption{w1,w2,w3}\input{VolumeIcons.tex}'
%   for n in 0 1 2 3; do mv Volume$n.pdf ../BGMApp/BGMApp/Images.xcassets/Volume$n.imageset/; done
%
% Might build correctly with regular LaTeX. I haven't tried it.
%

\documentclass[tikz]{standalone}
\usepackage{tikz}
% "dummyOption" prevents "Package optional Warning: No options were selected,
% so all optional text will be printed" when building Volume0.pdf.
\usepackage[dummyOption]{optional}

\begin{document}
\begin{tikzpicture}

% Speaker (Rounded box and triangle)
\fill[rounded corners=5mm]
    (0mm, 62.5mm) rectangle (25mm, 37.5mm) {};
\draw[rounded corners=2.5mm,fill=black]
    (3mm, 50mm)--(34mm, 76.5mm)--(34mm, 23.5mm)--cycle;

% First sound wave (Curved line)
\opt{w1}{
    \draw[line width=4.3mm,line cap=round]
        (44mm, 36.5mm) to[out=46,in=-46] (44mm, 63.5mm);
}

% Second sound wave (Curved line)
\opt{w2}{
    \draw[line width=4.3mm,line cap=round]
        (57.5mm, 27.5mm) to[out=46,in=-46] (57.5mm, 72.5mm);
}

% Third sound wave (Curved line)
\opt{w3}{
    \draw[line width=4.3mm,line cap=round]
        (72mm, 18.5mm) to[out=46,in=-46] (72mm, 81.5mm);
}

% Always draw a transparent copy of the third wave so the images will all have
% the same width.
\draw[line width=4.3mm,line cap=round,opacity=0]
    (72mm, 18.5mm) to[out=46,in=-46] (72mm, 81.5mm);

\end{tikzpicture}
\end{document}

'
%   xelatex -jobname=Volume1 '\def\UseOption{w1}% This file is part of Background Music.
%
% Background Music is free software: you can redistribute it and/or
% modify it under the terms of the GNU General Public License as
% published by the Free Software Foundation, either version 2 of the
% License, or (at your option) any later version.
%
% Background Music is distributed in the hope that it will be useful,
% but WITHOUT ANY WARRANTY; without even the implied warranty of
% MERCHANTABILITY or FITNESS FOR A PARTICULAR PURPOSE. See the
% GNU General Public License for more details.
%
% You should have received a copy of the GNU General Public License
% along with Background Music. If not, see <http://www.gnu.org/licenses/>.

%
% VolumeIcons.tex
%
% Build with XeTeX:
%   xelatex -jobname=Volume0 '\def\UseOption{}\input{VolumeIcons.tex}'
%   xelatex -jobname=Volume1 '\def\UseOption{w1}\input{VolumeIcons.tex}'
%   xelatex -jobname=Volume2 '\def\UseOption{w1,w2}\input{VolumeIcons.tex}'
%   xelatex -jobname=Volume3 '\def\UseOption{w1,w2,w3}\input{VolumeIcons.tex}'
%   for n in 0 1 2 3; do mv Volume$n.pdf ../BGMApp/BGMApp/Images.xcassets/Volume$n.imageset/; done
%
% Might build correctly with regular LaTeX. I haven't tried it.
%

\documentclass[tikz]{standalone}
\usepackage{tikz}
% "dummyOption" prevents "Package optional Warning: No options were selected,
% so all optional text will be printed" when building Volume0.pdf.
\usepackage[dummyOption]{optional}

\begin{document}
\begin{tikzpicture}

% Speaker (Rounded box and triangle)
\fill[rounded corners=5mm]
    (0mm, 62.5mm) rectangle (25mm, 37.5mm) {};
\draw[rounded corners=2.5mm,fill=black]
    (3mm, 50mm)--(34mm, 76.5mm)--(34mm, 23.5mm)--cycle;

% First sound wave (Curved line)
\opt{w1}{
    \draw[line width=4.3mm,line cap=round]
        (44mm, 36.5mm) to[out=46,in=-46] (44mm, 63.5mm);
}

% Second sound wave (Curved line)
\opt{w2}{
    \draw[line width=4.3mm,line cap=round]
        (57.5mm, 27.5mm) to[out=46,in=-46] (57.5mm, 72.5mm);
}

% Third sound wave (Curved line)
\opt{w3}{
    \draw[line width=4.3mm,line cap=round]
        (72mm, 18.5mm) to[out=46,in=-46] (72mm, 81.5mm);
}

% Always draw a transparent copy of the third wave so the images will all have
% the same width.
\draw[line width=4.3mm,line cap=round,opacity=0]
    (72mm, 18.5mm) to[out=46,in=-46] (72mm, 81.5mm);

\end{tikzpicture}
\end{document}

'
%   xelatex -jobname=Volume2 '\def\UseOption{w1,w2}% This file is part of Background Music.
%
% Background Music is free software: you can redistribute it and/or
% modify it under the terms of the GNU General Public License as
% published by the Free Software Foundation, either version 2 of the
% License, or (at your option) any later version.
%
% Background Music is distributed in the hope that it will be useful,
% but WITHOUT ANY WARRANTY; without even the implied warranty of
% MERCHANTABILITY or FITNESS FOR A PARTICULAR PURPOSE. See the
% GNU General Public License for more details.
%
% You should have received a copy of the GNU General Public License
% along with Background Music. If not, see <http://www.gnu.org/licenses/>.

%
% VolumeIcons.tex
%
% Build with XeTeX:
%   xelatex -jobname=Volume0 '\def\UseOption{}\input{VolumeIcons.tex}'
%   xelatex -jobname=Volume1 '\def\UseOption{w1}\input{VolumeIcons.tex}'
%   xelatex -jobname=Volume2 '\def\UseOption{w1,w2}\input{VolumeIcons.tex}'
%   xelatex -jobname=Volume3 '\def\UseOption{w1,w2,w3}\input{VolumeIcons.tex}'
%   for n in 0 1 2 3; do mv Volume$n.pdf ../BGMApp/BGMApp/Images.xcassets/Volume$n.imageset/; done
%
% Might build correctly with regular LaTeX. I haven't tried it.
%

\documentclass[tikz]{standalone}
\usepackage{tikz}
% "dummyOption" prevents "Package optional Warning: No options were selected,
% so all optional text will be printed" when building Volume0.pdf.
\usepackage[dummyOption]{optional}

\begin{document}
\begin{tikzpicture}

% Speaker (Rounded box and triangle)
\fill[rounded corners=5mm]
    (0mm, 62.5mm) rectangle (25mm, 37.5mm) {};
\draw[rounded corners=2.5mm,fill=black]
    (3mm, 50mm)--(34mm, 76.5mm)--(34mm, 23.5mm)--cycle;

% First sound wave (Curved line)
\opt{w1}{
    \draw[line width=4.3mm,line cap=round]
        (44mm, 36.5mm) to[out=46,in=-46] (44mm, 63.5mm);
}

% Second sound wave (Curved line)
\opt{w2}{
    \draw[line width=4.3mm,line cap=round]
        (57.5mm, 27.5mm) to[out=46,in=-46] (57.5mm, 72.5mm);
}

% Third sound wave (Curved line)
\opt{w3}{
    \draw[line width=4.3mm,line cap=round]
        (72mm, 18.5mm) to[out=46,in=-46] (72mm, 81.5mm);
}

% Always draw a transparent copy of the third wave so the images will all have
% the same width.
\draw[line width=4.3mm,line cap=round,opacity=0]
    (72mm, 18.5mm) to[out=46,in=-46] (72mm, 81.5mm);

\end{tikzpicture}
\end{document}

'
%   xelatex -jobname=Volume3 '\def\UseOption{w1,w2,w3}% This file is part of Background Music.
%
% Background Music is free software: you can redistribute it and/or
% modify it under the terms of the GNU General Public License as
% published by the Free Software Foundation, either version 2 of the
% License, or (at your option) any later version.
%
% Background Music is distributed in the hope that it will be useful,
% but WITHOUT ANY WARRANTY; without even the implied warranty of
% MERCHANTABILITY or FITNESS FOR A PARTICULAR PURPOSE. See the
% GNU General Public License for more details.
%
% You should have received a copy of the GNU General Public License
% along with Background Music. If not, see <http://www.gnu.org/licenses/>.

%
% VolumeIcons.tex
%
% Build with XeTeX:
%   xelatex -jobname=Volume0 '\def\UseOption{}\input{VolumeIcons.tex}'
%   xelatex -jobname=Volume1 '\def\UseOption{w1}\input{VolumeIcons.tex}'
%   xelatex -jobname=Volume2 '\def\UseOption{w1,w2}\input{VolumeIcons.tex}'
%   xelatex -jobname=Volume3 '\def\UseOption{w1,w2,w3}\input{VolumeIcons.tex}'
%   for n in 0 1 2 3; do mv Volume$n.pdf ../BGMApp/BGMApp/Images.xcassets/Volume$n.imageset/; done
%
% Might build correctly with regular LaTeX. I haven't tried it.
%

\documentclass[tikz]{standalone}
\usepackage{tikz}
% "dummyOption" prevents "Package optional Warning: No options were selected,
% so all optional text will be printed" when building Volume0.pdf.
\usepackage[dummyOption]{optional}

\begin{document}
\begin{tikzpicture}

% Speaker (Rounded box and triangle)
\fill[rounded corners=5mm]
    (0mm, 62.5mm) rectangle (25mm, 37.5mm) {};
\draw[rounded corners=2.5mm,fill=black]
    (3mm, 50mm)--(34mm, 76.5mm)--(34mm, 23.5mm)--cycle;

% First sound wave (Curved line)
\opt{w1}{
    \draw[line width=4.3mm,line cap=round]
        (44mm, 36.5mm) to[out=46,in=-46] (44mm, 63.5mm);
}

% Second sound wave (Curved line)
\opt{w2}{
    \draw[line width=4.3mm,line cap=round]
        (57.5mm, 27.5mm) to[out=46,in=-46] (57.5mm, 72.5mm);
}

% Third sound wave (Curved line)
\opt{w3}{
    \draw[line width=4.3mm,line cap=round]
        (72mm, 18.5mm) to[out=46,in=-46] (72mm, 81.5mm);
}

% Always draw a transparent copy of the third wave so the images will all have
% the same width.
\draw[line width=4.3mm,line cap=round,opacity=0]
    (72mm, 18.5mm) to[out=46,in=-46] (72mm, 81.5mm);

\end{tikzpicture}
\end{document}

'
%   for n in 0 1 2 3; do mv Volume$n.pdf ../BGMApp/BGMApp/Images.xcassets/Volume$n.imageset/; done
%
% Might build correctly with regular LaTeX. I haven't tried it.
%

\documentclass[tikz]{standalone}
\usepackage{tikz}
% "dummyOption" prevents "Package optional Warning: No options were selected,
% so all optional text will be printed" when building Volume0.pdf.
\usepackage[dummyOption]{optional}

\begin{document}
\begin{tikzpicture}

% Speaker (Rounded box and triangle)
\fill[rounded corners=5mm]
    (0mm, 62.5mm) rectangle (25mm, 37.5mm) {};
\draw[rounded corners=2.5mm,fill=black]
    (3mm, 50mm)--(34mm, 76.5mm)--(34mm, 23.5mm)--cycle;

% First sound wave (Curved line)
\opt{w1}{
    \draw[line width=4.3mm,line cap=round]
        (44mm, 36.5mm) to[out=46,in=-46] (44mm, 63.5mm);
}

% Second sound wave (Curved line)
\opt{w2}{
    \draw[line width=4.3mm,line cap=round]
        (57.5mm, 27.5mm) to[out=46,in=-46] (57.5mm, 72.5mm);
}

% Third sound wave (Curved line)
\opt{w3}{
    \draw[line width=4.3mm,line cap=round]
        (72mm, 18.5mm) to[out=46,in=-46] (72mm, 81.5mm);
}

% Always draw a transparent copy of the third wave so the images will all have
% the same width.
\draw[line width=4.3mm,line cap=round,opacity=0]
    (72mm, 18.5mm) to[out=46,in=-46] (72mm, 81.5mm);

\end{tikzpicture}
\end{document}

'
%   xelatex -jobname=Volume3 '\def\UseOption{w1,w2,w3}% This file is part of Background Music.
%
% Background Music is free software: you can redistribute it and/or
% modify it under the terms of the GNU General Public License as
% published by the Free Software Foundation, either version 2 of the
% License, or (at your option) any later version.
%
% Background Music is distributed in the hope that it will be useful,
% but WITHOUT ANY WARRANTY; without even the implied warranty of
% MERCHANTABILITY or FITNESS FOR A PARTICULAR PURPOSE. See the
% GNU General Public License for more details.
%
% You should have received a copy of the GNU General Public License
% along with Background Music. If not, see <http://www.gnu.org/licenses/>.

%
% VolumeIcons.tex
%
% Build with XeTeX:
%   xelatex -jobname=Volume0 '\def\UseOption{}% This file is part of Background Music.
%
% Background Music is free software: you can redistribute it and/or
% modify it under the terms of the GNU General Public License as
% published by the Free Software Foundation, either version 2 of the
% License, or (at your option) any later version.
%
% Background Music is distributed in the hope that it will be useful,
% but WITHOUT ANY WARRANTY; without even the implied warranty of
% MERCHANTABILITY or FITNESS FOR A PARTICULAR PURPOSE. See the
% GNU General Public License for more details.
%
% You should have received a copy of the GNU General Public License
% along with Background Music. If not, see <http://www.gnu.org/licenses/>.

%
% VolumeIcons.tex
%
% Build with XeTeX:
%   xelatex -jobname=Volume0 '\def\UseOption{}\input{VolumeIcons.tex}'
%   xelatex -jobname=Volume1 '\def\UseOption{w1}\input{VolumeIcons.tex}'
%   xelatex -jobname=Volume2 '\def\UseOption{w1,w2}\input{VolumeIcons.tex}'
%   xelatex -jobname=Volume3 '\def\UseOption{w1,w2,w3}\input{VolumeIcons.tex}'
%   for n in 0 1 2 3; do mv Volume$n.pdf ../BGMApp/BGMApp/Images.xcassets/Volume$n.imageset/; done
%
% Might build correctly with regular LaTeX. I haven't tried it.
%

\documentclass[tikz]{standalone}
\usepackage{tikz}
% "dummyOption" prevents "Package optional Warning: No options were selected,
% so all optional text will be printed" when building Volume0.pdf.
\usepackage[dummyOption]{optional}

\begin{document}
\begin{tikzpicture}

% Speaker (Rounded box and triangle)
\fill[rounded corners=5mm]
    (0mm, 62.5mm) rectangle (25mm, 37.5mm) {};
\draw[rounded corners=2.5mm,fill=black]
    (3mm, 50mm)--(34mm, 76.5mm)--(34mm, 23.5mm)--cycle;

% First sound wave (Curved line)
\opt{w1}{
    \draw[line width=4.3mm,line cap=round]
        (44mm, 36.5mm) to[out=46,in=-46] (44mm, 63.5mm);
}

% Second sound wave (Curved line)
\opt{w2}{
    \draw[line width=4.3mm,line cap=round]
        (57.5mm, 27.5mm) to[out=46,in=-46] (57.5mm, 72.5mm);
}

% Third sound wave (Curved line)
\opt{w3}{
    \draw[line width=4.3mm,line cap=round]
        (72mm, 18.5mm) to[out=46,in=-46] (72mm, 81.5mm);
}

% Always draw a transparent copy of the third wave so the images will all have
% the same width.
\draw[line width=4.3mm,line cap=round,opacity=0]
    (72mm, 18.5mm) to[out=46,in=-46] (72mm, 81.5mm);

\end{tikzpicture}
\end{document}

'
%   xelatex -jobname=Volume1 '\def\UseOption{w1}% This file is part of Background Music.
%
% Background Music is free software: you can redistribute it and/or
% modify it under the terms of the GNU General Public License as
% published by the Free Software Foundation, either version 2 of the
% License, or (at your option) any later version.
%
% Background Music is distributed in the hope that it will be useful,
% but WITHOUT ANY WARRANTY; without even the implied warranty of
% MERCHANTABILITY or FITNESS FOR A PARTICULAR PURPOSE. See the
% GNU General Public License for more details.
%
% You should have received a copy of the GNU General Public License
% along with Background Music. If not, see <http://www.gnu.org/licenses/>.

%
% VolumeIcons.tex
%
% Build with XeTeX:
%   xelatex -jobname=Volume0 '\def\UseOption{}\input{VolumeIcons.tex}'
%   xelatex -jobname=Volume1 '\def\UseOption{w1}\input{VolumeIcons.tex}'
%   xelatex -jobname=Volume2 '\def\UseOption{w1,w2}\input{VolumeIcons.tex}'
%   xelatex -jobname=Volume3 '\def\UseOption{w1,w2,w3}\input{VolumeIcons.tex}'
%   for n in 0 1 2 3; do mv Volume$n.pdf ../BGMApp/BGMApp/Images.xcassets/Volume$n.imageset/; done
%
% Might build correctly with regular LaTeX. I haven't tried it.
%

\documentclass[tikz]{standalone}
\usepackage{tikz}
% "dummyOption" prevents "Package optional Warning: No options were selected,
% so all optional text will be printed" when building Volume0.pdf.
\usepackage[dummyOption]{optional}

\begin{document}
\begin{tikzpicture}

% Speaker (Rounded box and triangle)
\fill[rounded corners=5mm]
    (0mm, 62.5mm) rectangle (25mm, 37.5mm) {};
\draw[rounded corners=2.5mm,fill=black]
    (3mm, 50mm)--(34mm, 76.5mm)--(34mm, 23.5mm)--cycle;

% First sound wave (Curved line)
\opt{w1}{
    \draw[line width=4.3mm,line cap=round]
        (44mm, 36.5mm) to[out=46,in=-46] (44mm, 63.5mm);
}

% Second sound wave (Curved line)
\opt{w2}{
    \draw[line width=4.3mm,line cap=round]
        (57.5mm, 27.5mm) to[out=46,in=-46] (57.5mm, 72.5mm);
}

% Third sound wave (Curved line)
\opt{w3}{
    \draw[line width=4.3mm,line cap=round]
        (72mm, 18.5mm) to[out=46,in=-46] (72mm, 81.5mm);
}

% Always draw a transparent copy of the third wave so the images will all have
% the same width.
\draw[line width=4.3mm,line cap=round,opacity=0]
    (72mm, 18.5mm) to[out=46,in=-46] (72mm, 81.5mm);

\end{tikzpicture}
\end{document}

'
%   xelatex -jobname=Volume2 '\def\UseOption{w1,w2}% This file is part of Background Music.
%
% Background Music is free software: you can redistribute it and/or
% modify it under the terms of the GNU General Public License as
% published by the Free Software Foundation, either version 2 of the
% License, or (at your option) any later version.
%
% Background Music is distributed in the hope that it will be useful,
% but WITHOUT ANY WARRANTY; without even the implied warranty of
% MERCHANTABILITY or FITNESS FOR A PARTICULAR PURPOSE. See the
% GNU General Public License for more details.
%
% You should have received a copy of the GNU General Public License
% along with Background Music. If not, see <http://www.gnu.org/licenses/>.

%
% VolumeIcons.tex
%
% Build with XeTeX:
%   xelatex -jobname=Volume0 '\def\UseOption{}\input{VolumeIcons.tex}'
%   xelatex -jobname=Volume1 '\def\UseOption{w1}\input{VolumeIcons.tex}'
%   xelatex -jobname=Volume2 '\def\UseOption{w1,w2}\input{VolumeIcons.tex}'
%   xelatex -jobname=Volume3 '\def\UseOption{w1,w2,w3}\input{VolumeIcons.tex}'
%   for n in 0 1 2 3; do mv Volume$n.pdf ../BGMApp/BGMApp/Images.xcassets/Volume$n.imageset/; done
%
% Might build correctly with regular LaTeX. I haven't tried it.
%

\documentclass[tikz]{standalone}
\usepackage{tikz}
% "dummyOption" prevents "Package optional Warning: No options were selected,
% so all optional text will be printed" when building Volume0.pdf.
\usepackage[dummyOption]{optional}

\begin{document}
\begin{tikzpicture}

% Speaker (Rounded box and triangle)
\fill[rounded corners=5mm]
    (0mm, 62.5mm) rectangle (25mm, 37.5mm) {};
\draw[rounded corners=2.5mm,fill=black]
    (3mm, 50mm)--(34mm, 76.5mm)--(34mm, 23.5mm)--cycle;

% First sound wave (Curved line)
\opt{w1}{
    \draw[line width=4.3mm,line cap=round]
        (44mm, 36.5mm) to[out=46,in=-46] (44mm, 63.5mm);
}

% Second sound wave (Curved line)
\opt{w2}{
    \draw[line width=4.3mm,line cap=round]
        (57.5mm, 27.5mm) to[out=46,in=-46] (57.5mm, 72.5mm);
}

% Third sound wave (Curved line)
\opt{w3}{
    \draw[line width=4.3mm,line cap=round]
        (72mm, 18.5mm) to[out=46,in=-46] (72mm, 81.5mm);
}

% Always draw a transparent copy of the third wave so the images will all have
% the same width.
\draw[line width=4.3mm,line cap=round,opacity=0]
    (72mm, 18.5mm) to[out=46,in=-46] (72mm, 81.5mm);

\end{tikzpicture}
\end{document}

'
%   xelatex -jobname=Volume3 '\def\UseOption{w1,w2,w3}% This file is part of Background Music.
%
% Background Music is free software: you can redistribute it and/or
% modify it under the terms of the GNU General Public License as
% published by the Free Software Foundation, either version 2 of the
% License, or (at your option) any later version.
%
% Background Music is distributed in the hope that it will be useful,
% but WITHOUT ANY WARRANTY; without even the implied warranty of
% MERCHANTABILITY or FITNESS FOR A PARTICULAR PURPOSE. See the
% GNU General Public License for more details.
%
% You should have received a copy of the GNU General Public License
% along with Background Music. If not, see <http://www.gnu.org/licenses/>.

%
% VolumeIcons.tex
%
% Build with XeTeX:
%   xelatex -jobname=Volume0 '\def\UseOption{}\input{VolumeIcons.tex}'
%   xelatex -jobname=Volume1 '\def\UseOption{w1}\input{VolumeIcons.tex}'
%   xelatex -jobname=Volume2 '\def\UseOption{w1,w2}\input{VolumeIcons.tex}'
%   xelatex -jobname=Volume3 '\def\UseOption{w1,w2,w3}\input{VolumeIcons.tex}'
%   for n in 0 1 2 3; do mv Volume$n.pdf ../BGMApp/BGMApp/Images.xcassets/Volume$n.imageset/; done
%
% Might build correctly with regular LaTeX. I haven't tried it.
%

\documentclass[tikz]{standalone}
\usepackage{tikz}
% "dummyOption" prevents "Package optional Warning: No options were selected,
% so all optional text will be printed" when building Volume0.pdf.
\usepackage[dummyOption]{optional}

\begin{document}
\begin{tikzpicture}

% Speaker (Rounded box and triangle)
\fill[rounded corners=5mm]
    (0mm, 62.5mm) rectangle (25mm, 37.5mm) {};
\draw[rounded corners=2.5mm,fill=black]
    (3mm, 50mm)--(34mm, 76.5mm)--(34mm, 23.5mm)--cycle;

% First sound wave (Curved line)
\opt{w1}{
    \draw[line width=4.3mm,line cap=round]
        (44mm, 36.5mm) to[out=46,in=-46] (44mm, 63.5mm);
}

% Second sound wave (Curved line)
\opt{w2}{
    \draw[line width=4.3mm,line cap=round]
        (57.5mm, 27.5mm) to[out=46,in=-46] (57.5mm, 72.5mm);
}

% Third sound wave (Curved line)
\opt{w3}{
    \draw[line width=4.3mm,line cap=round]
        (72mm, 18.5mm) to[out=46,in=-46] (72mm, 81.5mm);
}

% Always draw a transparent copy of the third wave so the images will all have
% the same width.
\draw[line width=4.3mm,line cap=round,opacity=0]
    (72mm, 18.5mm) to[out=46,in=-46] (72mm, 81.5mm);

\end{tikzpicture}
\end{document}

'
%   for n in 0 1 2 3; do mv Volume$n.pdf ../BGMApp/BGMApp/Images.xcassets/Volume$n.imageset/; done
%
% Might build correctly with regular LaTeX. I haven't tried it.
%

\documentclass[tikz]{standalone}
\usepackage{tikz}
% "dummyOption" prevents "Package optional Warning: No options were selected,
% so all optional text will be printed" when building Volume0.pdf.
\usepackage[dummyOption]{optional}

\begin{document}
\begin{tikzpicture}

% Speaker (Rounded box and triangle)
\fill[rounded corners=5mm]
    (0mm, 62.5mm) rectangle (25mm, 37.5mm) {};
\draw[rounded corners=2.5mm,fill=black]
    (3mm, 50mm)--(34mm, 76.5mm)--(34mm, 23.5mm)--cycle;

% First sound wave (Curved line)
\opt{w1}{
    \draw[line width=4.3mm,line cap=round]
        (44mm, 36.5mm) to[out=46,in=-46] (44mm, 63.5mm);
}

% Second sound wave (Curved line)
\opt{w2}{
    \draw[line width=4.3mm,line cap=round]
        (57.5mm, 27.5mm) to[out=46,in=-46] (57.5mm, 72.5mm);
}

% Third sound wave (Curved line)
\opt{w3}{
    \draw[line width=4.3mm,line cap=round]
        (72mm, 18.5mm) to[out=46,in=-46] (72mm, 81.5mm);
}

% Always draw a transparent copy of the third wave so the images will all have
% the same width.
\draw[line width=4.3mm,line cap=round,opacity=0]
    (72mm, 18.5mm) to[out=46,in=-46] (72mm, 81.5mm);

\end{tikzpicture}
\end{document}

'
%   for n in 0 1 2 3; do mv Volume$n.pdf ../BGMApp/BGMApp/Images.xcassets/Volume$n.imageset/; done
%
% Might build correctly with regular LaTeX. I haven't tried it.
%

\documentclass[tikz]{standalone}
\usepackage{tikz}
% "dummyOption" prevents "Package optional Warning: No options were selected,
% so all optional text will be printed" when building Volume0.pdf.
\usepackage[dummyOption]{optional}

\begin{document}
\begin{tikzpicture}

% Speaker (Rounded box and triangle)
\fill[rounded corners=5mm]
    (0mm, 62.5mm) rectangle (25mm, 37.5mm) {};
\draw[rounded corners=2.5mm,fill=black]
    (3mm, 50mm)--(34mm, 76.5mm)--(34mm, 23.5mm)--cycle;

% First sound wave (Curved line)
\opt{w1}{
    \draw[line width=4.3mm,line cap=round]
        (44mm, 36.5mm) to[out=46,in=-46] (44mm, 63.5mm);
}

% Second sound wave (Curved line)
\opt{w2}{
    \draw[line width=4.3mm,line cap=round]
        (57.5mm, 27.5mm) to[out=46,in=-46] (57.5mm, 72.5mm);
}

% Third sound wave (Curved line)
\opt{w3}{
    \draw[line width=4.3mm,line cap=round]
        (72mm, 18.5mm) to[out=46,in=-46] (72mm, 81.5mm);
}

% Always draw a transparent copy of the third wave so the images will all have
% the same width.
\draw[line width=4.3mm,line cap=round,opacity=0]
    (72mm, 18.5mm) to[out=46,in=-46] (72mm, 81.5mm);

\end{tikzpicture}
\end{document}

'
%   xelatex -jobname=Volume2 '\def\UseOption{w1,w2}% This file is part of Background Music.
%
% Background Music is free software: you can redistribute it and/or
% modify it under the terms of the GNU General Public License as
% published by the Free Software Foundation, either version 2 of the
% License, or (at your option) any later version.
%
% Background Music is distributed in the hope that it will be useful,
% but WITHOUT ANY WARRANTY; without even the implied warranty of
% MERCHANTABILITY or FITNESS FOR A PARTICULAR PURPOSE. See the
% GNU General Public License for more details.
%
% You should have received a copy of the GNU General Public License
% along with Background Music. If not, see <http://www.gnu.org/licenses/>.

%
% VolumeIcons.tex
%
% Build with XeTeX:
%   xelatex -jobname=Volume0 '\def\UseOption{}% This file is part of Background Music.
%
% Background Music is free software: you can redistribute it and/or
% modify it under the terms of the GNU General Public License as
% published by the Free Software Foundation, either version 2 of the
% License, or (at your option) any later version.
%
% Background Music is distributed in the hope that it will be useful,
% but WITHOUT ANY WARRANTY; without even the implied warranty of
% MERCHANTABILITY or FITNESS FOR A PARTICULAR PURPOSE. See the
% GNU General Public License for more details.
%
% You should have received a copy of the GNU General Public License
% along with Background Music. If not, see <http://www.gnu.org/licenses/>.

%
% VolumeIcons.tex
%
% Build with XeTeX:
%   xelatex -jobname=Volume0 '\def\UseOption{}% This file is part of Background Music.
%
% Background Music is free software: you can redistribute it and/or
% modify it under the terms of the GNU General Public License as
% published by the Free Software Foundation, either version 2 of the
% License, or (at your option) any later version.
%
% Background Music is distributed in the hope that it will be useful,
% but WITHOUT ANY WARRANTY; without even the implied warranty of
% MERCHANTABILITY or FITNESS FOR A PARTICULAR PURPOSE. See the
% GNU General Public License for more details.
%
% You should have received a copy of the GNU General Public License
% along with Background Music. If not, see <http://www.gnu.org/licenses/>.

%
% VolumeIcons.tex
%
% Build with XeTeX:
%   xelatex -jobname=Volume0 '\def\UseOption{}\input{VolumeIcons.tex}'
%   xelatex -jobname=Volume1 '\def\UseOption{w1}\input{VolumeIcons.tex}'
%   xelatex -jobname=Volume2 '\def\UseOption{w1,w2}\input{VolumeIcons.tex}'
%   xelatex -jobname=Volume3 '\def\UseOption{w1,w2,w3}\input{VolumeIcons.tex}'
%   for n in 0 1 2 3; do mv Volume$n.pdf ../BGMApp/BGMApp/Images.xcassets/Volume$n.imageset/; done
%
% Might build correctly with regular LaTeX. I haven't tried it.
%

\documentclass[tikz]{standalone}
\usepackage{tikz}
% "dummyOption" prevents "Package optional Warning: No options were selected,
% so all optional text will be printed" when building Volume0.pdf.
\usepackage[dummyOption]{optional}

\begin{document}
\begin{tikzpicture}

% Speaker (Rounded box and triangle)
\fill[rounded corners=5mm]
    (0mm, 62.5mm) rectangle (25mm, 37.5mm) {};
\draw[rounded corners=2.5mm,fill=black]
    (3mm, 50mm)--(34mm, 76.5mm)--(34mm, 23.5mm)--cycle;

% First sound wave (Curved line)
\opt{w1}{
    \draw[line width=4.3mm,line cap=round]
        (44mm, 36.5mm) to[out=46,in=-46] (44mm, 63.5mm);
}

% Second sound wave (Curved line)
\opt{w2}{
    \draw[line width=4.3mm,line cap=round]
        (57.5mm, 27.5mm) to[out=46,in=-46] (57.5mm, 72.5mm);
}

% Third sound wave (Curved line)
\opt{w3}{
    \draw[line width=4.3mm,line cap=round]
        (72mm, 18.5mm) to[out=46,in=-46] (72mm, 81.5mm);
}

% Always draw a transparent copy of the third wave so the images will all have
% the same width.
\draw[line width=4.3mm,line cap=round,opacity=0]
    (72mm, 18.5mm) to[out=46,in=-46] (72mm, 81.5mm);

\end{tikzpicture}
\end{document}

'
%   xelatex -jobname=Volume1 '\def\UseOption{w1}% This file is part of Background Music.
%
% Background Music is free software: you can redistribute it and/or
% modify it under the terms of the GNU General Public License as
% published by the Free Software Foundation, either version 2 of the
% License, or (at your option) any later version.
%
% Background Music is distributed in the hope that it will be useful,
% but WITHOUT ANY WARRANTY; without even the implied warranty of
% MERCHANTABILITY or FITNESS FOR A PARTICULAR PURPOSE. See the
% GNU General Public License for more details.
%
% You should have received a copy of the GNU General Public License
% along with Background Music. If not, see <http://www.gnu.org/licenses/>.

%
% VolumeIcons.tex
%
% Build with XeTeX:
%   xelatex -jobname=Volume0 '\def\UseOption{}\input{VolumeIcons.tex}'
%   xelatex -jobname=Volume1 '\def\UseOption{w1}\input{VolumeIcons.tex}'
%   xelatex -jobname=Volume2 '\def\UseOption{w1,w2}\input{VolumeIcons.tex}'
%   xelatex -jobname=Volume3 '\def\UseOption{w1,w2,w3}\input{VolumeIcons.tex}'
%   for n in 0 1 2 3; do mv Volume$n.pdf ../BGMApp/BGMApp/Images.xcassets/Volume$n.imageset/; done
%
% Might build correctly with regular LaTeX. I haven't tried it.
%

\documentclass[tikz]{standalone}
\usepackage{tikz}
% "dummyOption" prevents "Package optional Warning: No options were selected,
% so all optional text will be printed" when building Volume0.pdf.
\usepackage[dummyOption]{optional}

\begin{document}
\begin{tikzpicture}

% Speaker (Rounded box and triangle)
\fill[rounded corners=5mm]
    (0mm, 62.5mm) rectangle (25mm, 37.5mm) {};
\draw[rounded corners=2.5mm,fill=black]
    (3mm, 50mm)--(34mm, 76.5mm)--(34mm, 23.5mm)--cycle;

% First sound wave (Curved line)
\opt{w1}{
    \draw[line width=4.3mm,line cap=round]
        (44mm, 36.5mm) to[out=46,in=-46] (44mm, 63.5mm);
}

% Second sound wave (Curved line)
\opt{w2}{
    \draw[line width=4.3mm,line cap=round]
        (57.5mm, 27.5mm) to[out=46,in=-46] (57.5mm, 72.5mm);
}

% Third sound wave (Curved line)
\opt{w3}{
    \draw[line width=4.3mm,line cap=round]
        (72mm, 18.5mm) to[out=46,in=-46] (72mm, 81.5mm);
}

% Always draw a transparent copy of the third wave so the images will all have
% the same width.
\draw[line width=4.3mm,line cap=round,opacity=0]
    (72mm, 18.5mm) to[out=46,in=-46] (72mm, 81.5mm);

\end{tikzpicture}
\end{document}

'
%   xelatex -jobname=Volume2 '\def\UseOption{w1,w2}% This file is part of Background Music.
%
% Background Music is free software: you can redistribute it and/or
% modify it under the terms of the GNU General Public License as
% published by the Free Software Foundation, either version 2 of the
% License, or (at your option) any later version.
%
% Background Music is distributed in the hope that it will be useful,
% but WITHOUT ANY WARRANTY; without even the implied warranty of
% MERCHANTABILITY or FITNESS FOR A PARTICULAR PURPOSE. See the
% GNU General Public License for more details.
%
% You should have received a copy of the GNU General Public License
% along with Background Music. If not, see <http://www.gnu.org/licenses/>.

%
% VolumeIcons.tex
%
% Build with XeTeX:
%   xelatex -jobname=Volume0 '\def\UseOption{}\input{VolumeIcons.tex}'
%   xelatex -jobname=Volume1 '\def\UseOption{w1}\input{VolumeIcons.tex}'
%   xelatex -jobname=Volume2 '\def\UseOption{w1,w2}\input{VolumeIcons.tex}'
%   xelatex -jobname=Volume3 '\def\UseOption{w1,w2,w3}\input{VolumeIcons.tex}'
%   for n in 0 1 2 3; do mv Volume$n.pdf ../BGMApp/BGMApp/Images.xcassets/Volume$n.imageset/; done
%
% Might build correctly with regular LaTeX. I haven't tried it.
%

\documentclass[tikz]{standalone}
\usepackage{tikz}
% "dummyOption" prevents "Package optional Warning: No options were selected,
% so all optional text will be printed" when building Volume0.pdf.
\usepackage[dummyOption]{optional}

\begin{document}
\begin{tikzpicture}

% Speaker (Rounded box and triangle)
\fill[rounded corners=5mm]
    (0mm, 62.5mm) rectangle (25mm, 37.5mm) {};
\draw[rounded corners=2.5mm,fill=black]
    (3mm, 50mm)--(34mm, 76.5mm)--(34mm, 23.5mm)--cycle;

% First sound wave (Curved line)
\opt{w1}{
    \draw[line width=4.3mm,line cap=round]
        (44mm, 36.5mm) to[out=46,in=-46] (44mm, 63.5mm);
}

% Second sound wave (Curved line)
\opt{w2}{
    \draw[line width=4.3mm,line cap=round]
        (57.5mm, 27.5mm) to[out=46,in=-46] (57.5mm, 72.5mm);
}

% Third sound wave (Curved line)
\opt{w3}{
    \draw[line width=4.3mm,line cap=round]
        (72mm, 18.5mm) to[out=46,in=-46] (72mm, 81.5mm);
}

% Always draw a transparent copy of the third wave so the images will all have
% the same width.
\draw[line width=4.3mm,line cap=round,opacity=0]
    (72mm, 18.5mm) to[out=46,in=-46] (72mm, 81.5mm);

\end{tikzpicture}
\end{document}

'
%   xelatex -jobname=Volume3 '\def\UseOption{w1,w2,w3}% This file is part of Background Music.
%
% Background Music is free software: you can redistribute it and/or
% modify it under the terms of the GNU General Public License as
% published by the Free Software Foundation, either version 2 of the
% License, or (at your option) any later version.
%
% Background Music is distributed in the hope that it will be useful,
% but WITHOUT ANY WARRANTY; without even the implied warranty of
% MERCHANTABILITY or FITNESS FOR A PARTICULAR PURPOSE. See the
% GNU General Public License for more details.
%
% You should have received a copy of the GNU General Public License
% along with Background Music. If not, see <http://www.gnu.org/licenses/>.

%
% VolumeIcons.tex
%
% Build with XeTeX:
%   xelatex -jobname=Volume0 '\def\UseOption{}\input{VolumeIcons.tex}'
%   xelatex -jobname=Volume1 '\def\UseOption{w1}\input{VolumeIcons.tex}'
%   xelatex -jobname=Volume2 '\def\UseOption{w1,w2}\input{VolumeIcons.tex}'
%   xelatex -jobname=Volume3 '\def\UseOption{w1,w2,w3}\input{VolumeIcons.tex}'
%   for n in 0 1 2 3; do mv Volume$n.pdf ../BGMApp/BGMApp/Images.xcassets/Volume$n.imageset/; done
%
% Might build correctly with regular LaTeX. I haven't tried it.
%

\documentclass[tikz]{standalone}
\usepackage{tikz}
% "dummyOption" prevents "Package optional Warning: No options were selected,
% so all optional text will be printed" when building Volume0.pdf.
\usepackage[dummyOption]{optional}

\begin{document}
\begin{tikzpicture}

% Speaker (Rounded box and triangle)
\fill[rounded corners=5mm]
    (0mm, 62.5mm) rectangle (25mm, 37.5mm) {};
\draw[rounded corners=2.5mm,fill=black]
    (3mm, 50mm)--(34mm, 76.5mm)--(34mm, 23.5mm)--cycle;

% First sound wave (Curved line)
\opt{w1}{
    \draw[line width=4.3mm,line cap=round]
        (44mm, 36.5mm) to[out=46,in=-46] (44mm, 63.5mm);
}

% Second sound wave (Curved line)
\opt{w2}{
    \draw[line width=4.3mm,line cap=round]
        (57.5mm, 27.5mm) to[out=46,in=-46] (57.5mm, 72.5mm);
}

% Third sound wave (Curved line)
\opt{w3}{
    \draw[line width=4.3mm,line cap=round]
        (72mm, 18.5mm) to[out=46,in=-46] (72mm, 81.5mm);
}

% Always draw a transparent copy of the third wave so the images will all have
% the same width.
\draw[line width=4.3mm,line cap=round,opacity=0]
    (72mm, 18.5mm) to[out=46,in=-46] (72mm, 81.5mm);

\end{tikzpicture}
\end{document}

'
%   for n in 0 1 2 3; do mv Volume$n.pdf ../BGMApp/BGMApp/Images.xcassets/Volume$n.imageset/; done
%
% Might build correctly with regular LaTeX. I haven't tried it.
%

\documentclass[tikz]{standalone}
\usepackage{tikz}
% "dummyOption" prevents "Package optional Warning: No options were selected,
% so all optional text will be printed" when building Volume0.pdf.
\usepackage[dummyOption]{optional}

\begin{document}
\begin{tikzpicture}

% Speaker (Rounded box and triangle)
\fill[rounded corners=5mm]
    (0mm, 62.5mm) rectangle (25mm, 37.5mm) {};
\draw[rounded corners=2.5mm,fill=black]
    (3mm, 50mm)--(34mm, 76.5mm)--(34mm, 23.5mm)--cycle;

% First sound wave (Curved line)
\opt{w1}{
    \draw[line width=4.3mm,line cap=round]
        (44mm, 36.5mm) to[out=46,in=-46] (44mm, 63.5mm);
}

% Second sound wave (Curved line)
\opt{w2}{
    \draw[line width=4.3mm,line cap=round]
        (57.5mm, 27.5mm) to[out=46,in=-46] (57.5mm, 72.5mm);
}

% Third sound wave (Curved line)
\opt{w3}{
    \draw[line width=4.3mm,line cap=round]
        (72mm, 18.5mm) to[out=46,in=-46] (72mm, 81.5mm);
}

% Always draw a transparent copy of the third wave so the images will all have
% the same width.
\draw[line width=4.3mm,line cap=round,opacity=0]
    (72mm, 18.5mm) to[out=46,in=-46] (72mm, 81.5mm);

\end{tikzpicture}
\end{document}

'
%   xelatex -jobname=Volume1 '\def\UseOption{w1}% This file is part of Background Music.
%
% Background Music is free software: you can redistribute it and/or
% modify it under the terms of the GNU General Public License as
% published by the Free Software Foundation, either version 2 of the
% License, or (at your option) any later version.
%
% Background Music is distributed in the hope that it will be useful,
% but WITHOUT ANY WARRANTY; without even the implied warranty of
% MERCHANTABILITY or FITNESS FOR A PARTICULAR PURPOSE. See the
% GNU General Public License for more details.
%
% You should have received a copy of the GNU General Public License
% along with Background Music. If not, see <http://www.gnu.org/licenses/>.

%
% VolumeIcons.tex
%
% Build with XeTeX:
%   xelatex -jobname=Volume0 '\def\UseOption{}% This file is part of Background Music.
%
% Background Music is free software: you can redistribute it and/or
% modify it under the terms of the GNU General Public License as
% published by the Free Software Foundation, either version 2 of the
% License, or (at your option) any later version.
%
% Background Music is distributed in the hope that it will be useful,
% but WITHOUT ANY WARRANTY; without even the implied warranty of
% MERCHANTABILITY or FITNESS FOR A PARTICULAR PURPOSE. See the
% GNU General Public License for more details.
%
% You should have received a copy of the GNU General Public License
% along with Background Music. If not, see <http://www.gnu.org/licenses/>.

%
% VolumeIcons.tex
%
% Build with XeTeX:
%   xelatex -jobname=Volume0 '\def\UseOption{}\input{VolumeIcons.tex}'
%   xelatex -jobname=Volume1 '\def\UseOption{w1}\input{VolumeIcons.tex}'
%   xelatex -jobname=Volume2 '\def\UseOption{w1,w2}\input{VolumeIcons.tex}'
%   xelatex -jobname=Volume3 '\def\UseOption{w1,w2,w3}\input{VolumeIcons.tex}'
%   for n in 0 1 2 3; do mv Volume$n.pdf ../BGMApp/BGMApp/Images.xcassets/Volume$n.imageset/; done
%
% Might build correctly with regular LaTeX. I haven't tried it.
%

\documentclass[tikz]{standalone}
\usepackage{tikz}
% "dummyOption" prevents "Package optional Warning: No options were selected,
% so all optional text will be printed" when building Volume0.pdf.
\usepackage[dummyOption]{optional}

\begin{document}
\begin{tikzpicture}

% Speaker (Rounded box and triangle)
\fill[rounded corners=5mm]
    (0mm, 62.5mm) rectangle (25mm, 37.5mm) {};
\draw[rounded corners=2.5mm,fill=black]
    (3mm, 50mm)--(34mm, 76.5mm)--(34mm, 23.5mm)--cycle;

% First sound wave (Curved line)
\opt{w1}{
    \draw[line width=4.3mm,line cap=round]
        (44mm, 36.5mm) to[out=46,in=-46] (44mm, 63.5mm);
}

% Second sound wave (Curved line)
\opt{w2}{
    \draw[line width=4.3mm,line cap=round]
        (57.5mm, 27.5mm) to[out=46,in=-46] (57.5mm, 72.5mm);
}

% Third sound wave (Curved line)
\opt{w3}{
    \draw[line width=4.3mm,line cap=round]
        (72mm, 18.5mm) to[out=46,in=-46] (72mm, 81.5mm);
}

% Always draw a transparent copy of the third wave so the images will all have
% the same width.
\draw[line width=4.3mm,line cap=round,opacity=0]
    (72mm, 18.5mm) to[out=46,in=-46] (72mm, 81.5mm);

\end{tikzpicture}
\end{document}

'
%   xelatex -jobname=Volume1 '\def\UseOption{w1}% This file is part of Background Music.
%
% Background Music is free software: you can redistribute it and/or
% modify it under the terms of the GNU General Public License as
% published by the Free Software Foundation, either version 2 of the
% License, or (at your option) any later version.
%
% Background Music is distributed in the hope that it will be useful,
% but WITHOUT ANY WARRANTY; without even the implied warranty of
% MERCHANTABILITY or FITNESS FOR A PARTICULAR PURPOSE. See the
% GNU General Public License for more details.
%
% You should have received a copy of the GNU General Public License
% along with Background Music. If not, see <http://www.gnu.org/licenses/>.

%
% VolumeIcons.tex
%
% Build with XeTeX:
%   xelatex -jobname=Volume0 '\def\UseOption{}\input{VolumeIcons.tex}'
%   xelatex -jobname=Volume1 '\def\UseOption{w1}\input{VolumeIcons.tex}'
%   xelatex -jobname=Volume2 '\def\UseOption{w1,w2}\input{VolumeIcons.tex}'
%   xelatex -jobname=Volume3 '\def\UseOption{w1,w2,w3}\input{VolumeIcons.tex}'
%   for n in 0 1 2 3; do mv Volume$n.pdf ../BGMApp/BGMApp/Images.xcassets/Volume$n.imageset/; done
%
% Might build correctly with regular LaTeX. I haven't tried it.
%

\documentclass[tikz]{standalone}
\usepackage{tikz}
% "dummyOption" prevents "Package optional Warning: No options were selected,
% so all optional text will be printed" when building Volume0.pdf.
\usepackage[dummyOption]{optional}

\begin{document}
\begin{tikzpicture}

% Speaker (Rounded box and triangle)
\fill[rounded corners=5mm]
    (0mm, 62.5mm) rectangle (25mm, 37.5mm) {};
\draw[rounded corners=2.5mm,fill=black]
    (3mm, 50mm)--(34mm, 76.5mm)--(34mm, 23.5mm)--cycle;

% First sound wave (Curved line)
\opt{w1}{
    \draw[line width=4.3mm,line cap=round]
        (44mm, 36.5mm) to[out=46,in=-46] (44mm, 63.5mm);
}

% Second sound wave (Curved line)
\opt{w2}{
    \draw[line width=4.3mm,line cap=round]
        (57.5mm, 27.5mm) to[out=46,in=-46] (57.5mm, 72.5mm);
}

% Third sound wave (Curved line)
\opt{w3}{
    \draw[line width=4.3mm,line cap=round]
        (72mm, 18.5mm) to[out=46,in=-46] (72mm, 81.5mm);
}

% Always draw a transparent copy of the third wave so the images will all have
% the same width.
\draw[line width=4.3mm,line cap=round,opacity=0]
    (72mm, 18.5mm) to[out=46,in=-46] (72mm, 81.5mm);

\end{tikzpicture}
\end{document}

'
%   xelatex -jobname=Volume2 '\def\UseOption{w1,w2}% This file is part of Background Music.
%
% Background Music is free software: you can redistribute it and/or
% modify it under the terms of the GNU General Public License as
% published by the Free Software Foundation, either version 2 of the
% License, or (at your option) any later version.
%
% Background Music is distributed in the hope that it will be useful,
% but WITHOUT ANY WARRANTY; without even the implied warranty of
% MERCHANTABILITY or FITNESS FOR A PARTICULAR PURPOSE. See the
% GNU General Public License for more details.
%
% You should have received a copy of the GNU General Public License
% along with Background Music. If not, see <http://www.gnu.org/licenses/>.

%
% VolumeIcons.tex
%
% Build with XeTeX:
%   xelatex -jobname=Volume0 '\def\UseOption{}\input{VolumeIcons.tex}'
%   xelatex -jobname=Volume1 '\def\UseOption{w1}\input{VolumeIcons.tex}'
%   xelatex -jobname=Volume2 '\def\UseOption{w1,w2}\input{VolumeIcons.tex}'
%   xelatex -jobname=Volume3 '\def\UseOption{w1,w2,w3}\input{VolumeIcons.tex}'
%   for n in 0 1 2 3; do mv Volume$n.pdf ../BGMApp/BGMApp/Images.xcassets/Volume$n.imageset/; done
%
% Might build correctly with regular LaTeX. I haven't tried it.
%

\documentclass[tikz]{standalone}
\usepackage{tikz}
% "dummyOption" prevents "Package optional Warning: No options were selected,
% so all optional text will be printed" when building Volume0.pdf.
\usepackage[dummyOption]{optional}

\begin{document}
\begin{tikzpicture}

% Speaker (Rounded box and triangle)
\fill[rounded corners=5mm]
    (0mm, 62.5mm) rectangle (25mm, 37.5mm) {};
\draw[rounded corners=2.5mm,fill=black]
    (3mm, 50mm)--(34mm, 76.5mm)--(34mm, 23.5mm)--cycle;

% First sound wave (Curved line)
\opt{w1}{
    \draw[line width=4.3mm,line cap=round]
        (44mm, 36.5mm) to[out=46,in=-46] (44mm, 63.5mm);
}

% Second sound wave (Curved line)
\opt{w2}{
    \draw[line width=4.3mm,line cap=round]
        (57.5mm, 27.5mm) to[out=46,in=-46] (57.5mm, 72.5mm);
}

% Third sound wave (Curved line)
\opt{w3}{
    \draw[line width=4.3mm,line cap=round]
        (72mm, 18.5mm) to[out=46,in=-46] (72mm, 81.5mm);
}

% Always draw a transparent copy of the third wave so the images will all have
% the same width.
\draw[line width=4.3mm,line cap=round,opacity=0]
    (72mm, 18.5mm) to[out=46,in=-46] (72mm, 81.5mm);

\end{tikzpicture}
\end{document}

'
%   xelatex -jobname=Volume3 '\def\UseOption{w1,w2,w3}% This file is part of Background Music.
%
% Background Music is free software: you can redistribute it and/or
% modify it under the terms of the GNU General Public License as
% published by the Free Software Foundation, either version 2 of the
% License, or (at your option) any later version.
%
% Background Music is distributed in the hope that it will be useful,
% but WITHOUT ANY WARRANTY; without even the implied warranty of
% MERCHANTABILITY or FITNESS FOR A PARTICULAR PURPOSE. See the
% GNU General Public License for more details.
%
% You should have received a copy of the GNU General Public License
% along with Background Music. If not, see <http://www.gnu.org/licenses/>.

%
% VolumeIcons.tex
%
% Build with XeTeX:
%   xelatex -jobname=Volume0 '\def\UseOption{}\input{VolumeIcons.tex}'
%   xelatex -jobname=Volume1 '\def\UseOption{w1}\input{VolumeIcons.tex}'
%   xelatex -jobname=Volume2 '\def\UseOption{w1,w2}\input{VolumeIcons.tex}'
%   xelatex -jobname=Volume3 '\def\UseOption{w1,w2,w3}\input{VolumeIcons.tex}'
%   for n in 0 1 2 3; do mv Volume$n.pdf ../BGMApp/BGMApp/Images.xcassets/Volume$n.imageset/; done
%
% Might build correctly with regular LaTeX. I haven't tried it.
%

\documentclass[tikz]{standalone}
\usepackage{tikz}
% "dummyOption" prevents "Package optional Warning: No options were selected,
% so all optional text will be printed" when building Volume0.pdf.
\usepackage[dummyOption]{optional}

\begin{document}
\begin{tikzpicture}

% Speaker (Rounded box and triangle)
\fill[rounded corners=5mm]
    (0mm, 62.5mm) rectangle (25mm, 37.5mm) {};
\draw[rounded corners=2.5mm,fill=black]
    (3mm, 50mm)--(34mm, 76.5mm)--(34mm, 23.5mm)--cycle;

% First sound wave (Curved line)
\opt{w1}{
    \draw[line width=4.3mm,line cap=round]
        (44mm, 36.5mm) to[out=46,in=-46] (44mm, 63.5mm);
}

% Second sound wave (Curved line)
\opt{w2}{
    \draw[line width=4.3mm,line cap=round]
        (57.5mm, 27.5mm) to[out=46,in=-46] (57.5mm, 72.5mm);
}

% Third sound wave (Curved line)
\opt{w3}{
    \draw[line width=4.3mm,line cap=round]
        (72mm, 18.5mm) to[out=46,in=-46] (72mm, 81.5mm);
}

% Always draw a transparent copy of the third wave so the images will all have
% the same width.
\draw[line width=4.3mm,line cap=round,opacity=0]
    (72mm, 18.5mm) to[out=46,in=-46] (72mm, 81.5mm);

\end{tikzpicture}
\end{document}

'
%   for n in 0 1 2 3; do mv Volume$n.pdf ../BGMApp/BGMApp/Images.xcassets/Volume$n.imageset/; done
%
% Might build correctly with regular LaTeX. I haven't tried it.
%

\documentclass[tikz]{standalone}
\usepackage{tikz}
% "dummyOption" prevents "Package optional Warning: No options were selected,
% so all optional text will be printed" when building Volume0.pdf.
\usepackage[dummyOption]{optional}

\begin{document}
\begin{tikzpicture}

% Speaker (Rounded box and triangle)
\fill[rounded corners=5mm]
    (0mm, 62.5mm) rectangle (25mm, 37.5mm) {};
\draw[rounded corners=2.5mm,fill=black]
    (3mm, 50mm)--(34mm, 76.5mm)--(34mm, 23.5mm)--cycle;

% First sound wave (Curved line)
\opt{w1}{
    \draw[line width=4.3mm,line cap=round]
        (44mm, 36.5mm) to[out=46,in=-46] (44mm, 63.5mm);
}

% Second sound wave (Curved line)
\opt{w2}{
    \draw[line width=4.3mm,line cap=round]
        (57.5mm, 27.5mm) to[out=46,in=-46] (57.5mm, 72.5mm);
}

% Third sound wave (Curved line)
\opt{w3}{
    \draw[line width=4.3mm,line cap=round]
        (72mm, 18.5mm) to[out=46,in=-46] (72mm, 81.5mm);
}

% Always draw a transparent copy of the third wave so the images will all have
% the same width.
\draw[line width=4.3mm,line cap=round,opacity=0]
    (72mm, 18.5mm) to[out=46,in=-46] (72mm, 81.5mm);

\end{tikzpicture}
\end{document}

'
%   xelatex -jobname=Volume2 '\def\UseOption{w1,w2}% This file is part of Background Music.
%
% Background Music is free software: you can redistribute it and/or
% modify it under the terms of the GNU General Public License as
% published by the Free Software Foundation, either version 2 of the
% License, or (at your option) any later version.
%
% Background Music is distributed in the hope that it will be useful,
% but WITHOUT ANY WARRANTY; without even the implied warranty of
% MERCHANTABILITY or FITNESS FOR A PARTICULAR PURPOSE. See the
% GNU General Public License for more details.
%
% You should have received a copy of the GNU General Public License
% along with Background Music. If not, see <http://www.gnu.org/licenses/>.

%
% VolumeIcons.tex
%
% Build with XeTeX:
%   xelatex -jobname=Volume0 '\def\UseOption{}% This file is part of Background Music.
%
% Background Music is free software: you can redistribute it and/or
% modify it under the terms of the GNU General Public License as
% published by the Free Software Foundation, either version 2 of the
% License, or (at your option) any later version.
%
% Background Music is distributed in the hope that it will be useful,
% but WITHOUT ANY WARRANTY; without even the implied warranty of
% MERCHANTABILITY or FITNESS FOR A PARTICULAR PURPOSE. See the
% GNU General Public License for more details.
%
% You should have received a copy of the GNU General Public License
% along with Background Music. If not, see <http://www.gnu.org/licenses/>.

%
% VolumeIcons.tex
%
% Build with XeTeX:
%   xelatex -jobname=Volume0 '\def\UseOption{}\input{VolumeIcons.tex}'
%   xelatex -jobname=Volume1 '\def\UseOption{w1}\input{VolumeIcons.tex}'
%   xelatex -jobname=Volume2 '\def\UseOption{w1,w2}\input{VolumeIcons.tex}'
%   xelatex -jobname=Volume3 '\def\UseOption{w1,w2,w3}\input{VolumeIcons.tex}'
%   for n in 0 1 2 3; do mv Volume$n.pdf ../BGMApp/BGMApp/Images.xcassets/Volume$n.imageset/; done
%
% Might build correctly with regular LaTeX. I haven't tried it.
%

\documentclass[tikz]{standalone}
\usepackage{tikz}
% "dummyOption" prevents "Package optional Warning: No options were selected,
% so all optional text will be printed" when building Volume0.pdf.
\usepackage[dummyOption]{optional}

\begin{document}
\begin{tikzpicture}

% Speaker (Rounded box and triangle)
\fill[rounded corners=5mm]
    (0mm, 62.5mm) rectangle (25mm, 37.5mm) {};
\draw[rounded corners=2.5mm,fill=black]
    (3mm, 50mm)--(34mm, 76.5mm)--(34mm, 23.5mm)--cycle;

% First sound wave (Curved line)
\opt{w1}{
    \draw[line width=4.3mm,line cap=round]
        (44mm, 36.5mm) to[out=46,in=-46] (44mm, 63.5mm);
}

% Second sound wave (Curved line)
\opt{w2}{
    \draw[line width=4.3mm,line cap=round]
        (57.5mm, 27.5mm) to[out=46,in=-46] (57.5mm, 72.5mm);
}

% Third sound wave (Curved line)
\opt{w3}{
    \draw[line width=4.3mm,line cap=round]
        (72mm, 18.5mm) to[out=46,in=-46] (72mm, 81.5mm);
}

% Always draw a transparent copy of the third wave so the images will all have
% the same width.
\draw[line width=4.3mm,line cap=round,opacity=0]
    (72mm, 18.5mm) to[out=46,in=-46] (72mm, 81.5mm);

\end{tikzpicture}
\end{document}

'
%   xelatex -jobname=Volume1 '\def\UseOption{w1}% This file is part of Background Music.
%
% Background Music is free software: you can redistribute it and/or
% modify it under the terms of the GNU General Public License as
% published by the Free Software Foundation, either version 2 of the
% License, or (at your option) any later version.
%
% Background Music is distributed in the hope that it will be useful,
% but WITHOUT ANY WARRANTY; without even the implied warranty of
% MERCHANTABILITY or FITNESS FOR A PARTICULAR PURPOSE. See the
% GNU General Public License for more details.
%
% You should have received a copy of the GNU General Public License
% along with Background Music. If not, see <http://www.gnu.org/licenses/>.

%
% VolumeIcons.tex
%
% Build with XeTeX:
%   xelatex -jobname=Volume0 '\def\UseOption{}\input{VolumeIcons.tex}'
%   xelatex -jobname=Volume1 '\def\UseOption{w1}\input{VolumeIcons.tex}'
%   xelatex -jobname=Volume2 '\def\UseOption{w1,w2}\input{VolumeIcons.tex}'
%   xelatex -jobname=Volume3 '\def\UseOption{w1,w2,w3}\input{VolumeIcons.tex}'
%   for n in 0 1 2 3; do mv Volume$n.pdf ../BGMApp/BGMApp/Images.xcassets/Volume$n.imageset/; done
%
% Might build correctly with regular LaTeX. I haven't tried it.
%

\documentclass[tikz]{standalone}
\usepackage{tikz}
% "dummyOption" prevents "Package optional Warning: No options were selected,
% so all optional text will be printed" when building Volume0.pdf.
\usepackage[dummyOption]{optional}

\begin{document}
\begin{tikzpicture}

% Speaker (Rounded box and triangle)
\fill[rounded corners=5mm]
    (0mm, 62.5mm) rectangle (25mm, 37.5mm) {};
\draw[rounded corners=2.5mm,fill=black]
    (3mm, 50mm)--(34mm, 76.5mm)--(34mm, 23.5mm)--cycle;

% First sound wave (Curved line)
\opt{w1}{
    \draw[line width=4.3mm,line cap=round]
        (44mm, 36.5mm) to[out=46,in=-46] (44mm, 63.5mm);
}

% Second sound wave (Curved line)
\opt{w2}{
    \draw[line width=4.3mm,line cap=round]
        (57.5mm, 27.5mm) to[out=46,in=-46] (57.5mm, 72.5mm);
}

% Third sound wave (Curved line)
\opt{w3}{
    \draw[line width=4.3mm,line cap=round]
        (72mm, 18.5mm) to[out=46,in=-46] (72mm, 81.5mm);
}

% Always draw a transparent copy of the third wave so the images will all have
% the same width.
\draw[line width=4.3mm,line cap=round,opacity=0]
    (72mm, 18.5mm) to[out=46,in=-46] (72mm, 81.5mm);

\end{tikzpicture}
\end{document}

'
%   xelatex -jobname=Volume2 '\def\UseOption{w1,w2}% This file is part of Background Music.
%
% Background Music is free software: you can redistribute it and/or
% modify it under the terms of the GNU General Public License as
% published by the Free Software Foundation, either version 2 of the
% License, or (at your option) any later version.
%
% Background Music is distributed in the hope that it will be useful,
% but WITHOUT ANY WARRANTY; without even the implied warranty of
% MERCHANTABILITY or FITNESS FOR A PARTICULAR PURPOSE. See the
% GNU General Public License for more details.
%
% You should have received a copy of the GNU General Public License
% along with Background Music. If not, see <http://www.gnu.org/licenses/>.

%
% VolumeIcons.tex
%
% Build with XeTeX:
%   xelatex -jobname=Volume0 '\def\UseOption{}\input{VolumeIcons.tex}'
%   xelatex -jobname=Volume1 '\def\UseOption{w1}\input{VolumeIcons.tex}'
%   xelatex -jobname=Volume2 '\def\UseOption{w1,w2}\input{VolumeIcons.tex}'
%   xelatex -jobname=Volume3 '\def\UseOption{w1,w2,w3}\input{VolumeIcons.tex}'
%   for n in 0 1 2 3; do mv Volume$n.pdf ../BGMApp/BGMApp/Images.xcassets/Volume$n.imageset/; done
%
% Might build correctly with regular LaTeX. I haven't tried it.
%

\documentclass[tikz]{standalone}
\usepackage{tikz}
% "dummyOption" prevents "Package optional Warning: No options were selected,
% so all optional text will be printed" when building Volume0.pdf.
\usepackage[dummyOption]{optional}

\begin{document}
\begin{tikzpicture}

% Speaker (Rounded box and triangle)
\fill[rounded corners=5mm]
    (0mm, 62.5mm) rectangle (25mm, 37.5mm) {};
\draw[rounded corners=2.5mm,fill=black]
    (3mm, 50mm)--(34mm, 76.5mm)--(34mm, 23.5mm)--cycle;

% First sound wave (Curved line)
\opt{w1}{
    \draw[line width=4.3mm,line cap=round]
        (44mm, 36.5mm) to[out=46,in=-46] (44mm, 63.5mm);
}

% Second sound wave (Curved line)
\opt{w2}{
    \draw[line width=4.3mm,line cap=round]
        (57.5mm, 27.5mm) to[out=46,in=-46] (57.5mm, 72.5mm);
}

% Third sound wave (Curved line)
\opt{w3}{
    \draw[line width=4.3mm,line cap=round]
        (72mm, 18.5mm) to[out=46,in=-46] (72mm, 81.5mm);
}

% Always draw a transparent copy of the third wave so the images will all have
% the same width.
\draw[line width=4.3mm,line cap=round,opacity=0]
    (72mm, 18.5mm) to[out=46,in=-46] (72mm, 81.5mm);

\end{tikzpicture}
\end{document}

'
%   xelatex -jobname=Volume3 '\def\UseOption{w1,w2,w3}% This file is part of Background Music.
%
% Background Music is free software: you can redistribute it and/or
% modify it under the terms of the GNU General Public License as
% published by the Free Software Foundation, either version 2 of the
% License, or (at your option) any later version.
%
% Background Music is distributed in the hope that it will be useful,
% but WITHOUT ANY WARRANTY; without even the implied warranty of
% MERCHANTABILITY or FITNESS FOR A PARTICULAR PURPOSE. See the
% GNU General Public License for more details.
%
% You should have received a copy of the GNU General Public License
% along with Background Music. If not, see <http://www.gnu.org/licenses/>.

%
% VolumeIcons.tex
%
% Build with XeTeX:
%   xelatex -jobname=Volume0 '\def\UseOption{}\input{VolumeIcons.tex}'
%   xelatex -jobname=Volume1 '\def\UseOption{w1}\input{VolumeIcons.tex}'
%   xelatex -jobname=Volume2 '\def\UseOption{w1,w2}\input{VolumeIcons.tex}'
%   xelatex -jobname=Volume3 '\def\UseOption{w1,w2,w3}\input{VolumeIcons.tex}'
%   for n in 0 1 2 3; do mv Volume$n.pdf ../BGMApp/BGMApp/Images.xcassets/Volume$n.imageset/; done
%
% Might build correctly with regular LaTeX. I haven't tried it.
%

\documentclass[tikz]{standalone}
\usepackage{tikz}
% "dummyOption" prevents "Package optional Warning: No options were selected,
% so all optional text will be printed" when building Volume0.pdf.
\usepackage[dummyOption]{optional}

\begin{document}
\begin{tikzpicture}

% Speaker (Rounded box and triangle)
\fill[rounded corners=5mm]
    (0mm, 62.5mm) rectangle (25mm, 37.5mm) {};
\draw[rounded corners=2.5mm,fill=black]
    (3mm, 50mm)--(34mm, 76.5mm)--(34mm, 23.5mm)--cycle;

% First sound wave (Curved line)
\opt{w1}{
    \draw[line width=4.3mm,line cap=round]
        (44mm, 36.5mm) to[out=46,in=-46] (44mm, 63.5mm);
}

% Second sound wave (Curved line)
\opt{w2}{
    \draw[line width=4.3mm,line cap=round]
        (57.5mm, 27.5mm) to[out=46,in=-46] (57.5mm, 72.5mm);
}

% Third sound wave (Curved line)
\opt{w3}{
    \draw[line width=4.3mm,line cap=round]
        (72mm, 18.5mm) to[out=46,in=-46] (72mm, 81.5mm);
}

% Always draw a transparent copy of the third wave so the images will all have
% the same width.
\draw[line width=4.3mm,line cap=round,opacity=0]
    (72mm, 18.5mm) to[out=46,in=-46] (72mm, 81.5mm);

\end{tikzpicture}
\end{document}

'
%   for n in 0 1 2 3; do mv Volume$n.pdf ../BGMApp/BGMApp/Images.xcassets/Volume$n.imageset/; done
%
% Might build correctly with regular LaTeX. I haven't tried it.
%

\documentclass[tikz]{standalone}
\usepackage{tikz}
% "dummyOption" prevents "Package optional Warning: No options were selected,
% so all optional text will be printed" when building Volume0.pdf.
\usepackage[dummyOption]{optional}

\begin{document}
\begin{tikzpicture}

% Speaker (Rounded box and triangle)
\fill[rounded corners=5mm]
    (0mm, 62.5mm) rectangle (25mm, 37.5mm) {};
\draw[rounded corners=2.5mm,fill=black]
    (3mm, 50mm)--(34mm, 76.5mm)--(34mm, 23.5mm)--cycle;

% First sound wave (Curved line)
\opt{w1}{
    \draw[line width=4.3mm,line cap=round]
        (44mm, 36.5mm) to[out=46,in=-46] (44mm, 63.5mm);
}

% Second sound wave (Curved line)
\opt{w2}{
    \draw[line width=4.3mm,line cap=round]
        (57.5mm, 27.5mm) to[out=46,in=-46] (57.5mm, 72.5mm);
}

% Third sound wave (Curved line)
\opt{w3}{
    \draw[line width=4.3mm,line cap=round]
        (72mm, 18.5mm) to[out=46,in=-46] (72mm, 81.5mm);
}

% Always draw a transparent copy of the third wave so the images will all have
% the same width.
\draw[line width=4.3mm,line cap=round,opacity=0]
    (72mm, 18.5mm) to[out=46,in=-46] (72mm, 81.5mm);

\end{tikzpicture}
\end{document}

'
%   xelatex -jobname=Volume3 '\def\UseOption{w1,w2,w3}% This file is part of Background Music.
%
% Background Music is free software: you can redistribute it and/or
% modify it under the terms of the GNU General Public License as
% published by the Free Software Foundation, either version 2 of the
% License, or (at your option) any later version.
%
% Background Music is distributed in the hope that it will be useful,
% but WITHOUT ANY WARRANTY; without even the implied warranty of
% MERCHANTABILITY or FITNESS FOR A PARTICULAR PURPOSE. See the
% GNU General Public License for more details.
%
% You should have received a copy of the GNU General Public License
% along with Background Music. If not, see <http://www.gnu.org/licenses/>.

%
% VolumeIcons.tex
%
% Build with XeTeX:
%   xelatex -jobname=Volume0 '\def\UseOption{}% This file is part of Background Music.
%
% Background Music is free software: you can redistribute it and/or
% modify it under the terms of the GNU General Public License as
% published by the Free Software Foundation, either version 2 of the
% License, or (at your option) any later version.
%
% Background Music is distributed in the hope that it will be useful,
% but WITHOUT ANY WARRANTY; without even the implied warranty of
% MERCHANTABILITY or FITNESS FOR A PARTICULAR PURPOSE. See the
% GNU General Public License for more details.
%
% You should have received a copy of the GNU General Public License
% along with Background Music. If not, see <http://www.gnu.org/licenses/>.

%
% VolumeIcons.tex
%
% Build with XeTeX:
%   xelatex -jobname=Volume0 '\def\UseOption{}\input{VolumeIcons.tex}'
%   xelatex -jobname=Volume1 '\def\UseOption{w1}\input{VolumeIcons.tex}'
%   xelatex -jobname=Volume2 '\def\UseOption{w1,w2}\input{VolumeIcons.tex}'
%   xelatex -jobname=Volume3 '\def\UseOption{w1,w2,w3}\input{VolumeIcons.tex}'
%   for n in 0 1 2 3; do mv Volume$n.pdf ../BGMApp/BGMApp/Images.xcassets/Volume$n.imageset/; done
%
% Might build correctly with regular LaTeX. I haven't tried it.
%

\documentclass[tikz]{standalone}
\usepackage{tikz}
% "dummyOption" prevents "Package optional Warning: No options were selected,
% so all optional text will be printed" when building Volume0.pdf.
\usepackage[dummyOption]{optional}

\begin{document}
\begin{tikzpicture}

% Speaker (Rounded box and triangle)
\fill[rounded corners=5mm]
    (0mm, 62.5mm) rectangle (25mm, 37.5mm) {};
\draw[rounded corners=2.5mm,fill=black]
    (3mm, 50mm)--(34mm, 76.5mm)--(34mm, 23.5mm)--cycle;

% First sound wave (Curved line)
\opt{w1}{
    \draw[line width=4.3mm,line cap=round]
        (44mm, 36.5mm) to[out=46,in=-46] (44mm, 63.5mm);
}

% Second sound wave (Curved line)
\opt{w2}{
    \draw[line width=4.3mm,line cap=round]
        (57.5mm, 27.5mm) to[out=46,in=-46] (57.5mm, 72.5mm);
}

% Third sound wave (Curved line)
\opt{w3}{
    \draw[line width=4.3mm,line cap=round]
        (72mm, 18.5mm) to[out=46,in=-46] (72mm, 81.5mm);
}

% Always draw a transparent copy of the third wave so the images will all have
% the same width.
\draw[line width=4.3mm,line cap=round,opacity=0]
    (72mm, 18.5mm) to[out=46,in=-46] (72mm, 81.5mm);

\end{tikzpicture}
\end{document}

'
%   xelatex -jobname=Volume1 '\def\UseOption{w1}% This file is part of Background Music.
%
% Background Music is free software: you can redistribute it and/or
% modify it under the terms of the GNU General Public License as
% published by the Free Software Foundation, either version 2 of the
% License, or (at your option) any later version.
%
% Background Music is distributed in the hope that it will be useful,
% but WITHOUT ANY WARRANTY; without even the implied warranty of
% MERCHANTABILITY or FITNESS FOR A PARTICULAR PURPOSE. See the
% GNU General Public License for more details.
%
% You should have received a copy of the GNU General Public License
% along with Background Music. If not, see <http://www.gnu.org/licenses/>.

%
% VolumeIcons.tex
%
% Build with XeTeX:
%   xelatex -jobname=Volume0 '\def\UseOption{}\input{VolumeIcons.tex}'
%   xelatex -jobname=Volume1 '\def\UseOption{w1}\input{VolumeIcons.tex}'
%   xelatex -jobname=Volume2 '\def\UseOption{w1,w2}\input{VolumeIcons.tex}'
%   xelatex -jobname=Volume3 '\def\UseOption{w1,w2,w3}\input{VolumeIcons.tex}'
%   for n in 0 1 2 3; do mv Volume$n.pdf ../BGMApp/BGMApp/Images.xcassets/Volume$n.imageset/; done
%
% Might build correctly with regular LaTeX. I haven't tried it.
%

\documentclass[tikz]{standalone}
\usepackage{tikz}
% "dummyOption" prevents "Package optional Warning: No options were selected,
% so all optional text will be printed" when building Volume0.pdf.
\usepackage[dummyOption]{optional}

\begin{document}
\begin{tikzpicture}

% Speaker (Rounded box and triangle)
\fill[rounded corners=5mm]
    (0mm, 62.5mm) rectangle (25mm, 37.5mm) {};
\draw[rounded corners=2.5mm,fill=black]
    (3mm, 50mm)--(34mm, 76.5mm)--(34mm, 23.5mm)--cycle;

% First sound wave (Curved line)
\opt{w1}{
    \draw[line width=4.3mm,line cap=round]
        (44mm, 36.5mm) to[out=46,in=-46] (44mm, 63.5mm);
}

% Second sound wave (Curved line)
\opt{w2}{
    \draw[line width=4.3mm,line cap=round]
        (57.5mm, 27.5mm) to[out=46,in=-46] (57.5mm, 72.5mm);
}

% Third sound wave (Curved line)
\opt{w3}{
    \draw[line width=4.3mm,line cap=round]
        (72mm, 18.5mm) to[out=46,in=-46] (72mm, 81.5mm);
}

% Always draw a transparent copy of the third wave so the images will all have
% the same width.
\draw[line width=4.3mm,line cap=round,opacity=0]
    (72mm, 18.5mm) to[out=46,in=-46] (72mm, 81.5mm);

\end{tikzpicture}
\end{document}

'
%   xelatex -jobname=Volume2 '\def\UseOption{w1,w2}% This file is part of Background Music.
%
% Background Music is free software: you can redistribute it and/or
% modify it under the terms of the GNU General Public License as
% published by the Free Software Foundation, either version 2 of the
% License, or (at your option) any later version.
%
% Background Music is distributed in the hope that it will be useful,
% but WITHOUT ANY WARRANTY; without even the implied warranty of
% MERCHANTABILITY or FITNESS FOR A PARTICULAR PURPOSE. See the
% GNU General Public License for more details.
%
% You should have received a copy of the GNU General Public License
% along with Background Music. If not, see <http://www.gnu.org/licenses/>.

%
% VolumeIcons.tex
%
% Build with XeTeX:
%   xelatex -jobname=Volume0 '\def\UseOption{}\input{VolumeIcons.tex}'
%   xelatex -jobname=Volume1 '\def\UseOption{w1}\input{VolumeIcons.tex}'
%   xelatex -jobname=Volume2 '\def\UseOption{w1,w2}\input{VolumeIcons.tex}'
%   xelatex -jobname=Volume3 '\def\UseOption{w1,w2,w3}\input{VolumeIcons.tex}'
%   for n in 0 1 2 3; do mv Volume$n.pdf ../BGMApp/BGMApp/Images.xcassets/Volume$n.imageset/; done
%
% Might build correctly with regular LaTeX. I haven't tried it.
%

\documentclass[tikz]{standalone}
\usepackage{tikz}
% "dummyOption" prevents "Package optional Warning: No options were selected,
% so all optional text will be printed" when building Volume0.pdf.
\usepackage[dummyOption]{optional}

\begin{document}
\begin{tikzpicture}

% Speaker (Rounded box and triangle)
\fill[rounded corners=5mm]
    (0mm, 62.5mm) rectangle (25mm, 37.5mm) {};
\draw[rounded corners=2.5mm,fill=black]
    (3mm, 50mm)--(34mm, 76.5mm)--(34mm, 23.5mm)--cycle;

% First sound wave (Curved line)
\opt{w1}{
    \draw[line width=4.3mm,line cap=round]
        (44mm, 36.5mm) to[out=46,in=-46] (44mm, 63.5mm);
}

% Second sound wave (Curved line)
\opt{w2}{
    \draw[line width=4.3mm,line cap=round]
        (57.5mm, 27.5mm) to[out=46,in=-46] (57.5mm, 72.5mm);
}

% Third sound wave (Curved line)
\opt{w3}{
    \draw[line width=4.3mm,line cap=round]
        (72mm, 18.5mm) to[out=46,in=-46] (72mm, 81.5mm);
}

% Always draw a transparent copy of the third wave so the images will all have
% the same width.
\draw[line width=4.3mm,line cap=round,opacity=0]
    (72mm, 18.5mm) to[out=46,in=-46] (72mm, 81.5mm);

\end{tikzpicture}
\end{document}

'
%   xelatex -jobname=Volume3 '\def\UseOption{w1,w2,w3}% This file is part of Background Music.
%
% Background Music is free software: you can redistribute it and/or
% modify it under the terms of the GNU General Public License as
% published by the Free Software Foundation, either version 2 of the
% License, or (at your option) any later version.
%
% Background Music is distributed in the hope that it will be useful,
% but WITHOUT ANY WARRANTY; without even the implied warranty of
% MERCHANTABILITY or FITNESS FOR A PARTICULAR PURPOSE. See the
% GNU General Public License for more details.
%
% You should have received a copy of the GNU General Public License
% along with Background Music. If not, see <http://www.gnu.org/licenses/>.

%
% VolumeIcons.tex
%
% Build with XeTeX:
%   xelatex -jobname=Volume0 '\def\UseOption{}\input{VolumeIcons.tex}'
%   xelatex -jobname=Volume1 '\def\UseOption{w1}\input{VolumeIcons.tex}'
%   xelatex -jobname=Volume2 '\def\UseOption{w1,w2}\input{VolumeIcons.tex}'
%   xelatex -jobname=Volume3 '\def\UseOption{w1,w2,w3}\input{VolumeIcons.tex}'
%   for n in 0 1 2 3; do mv Volume$n.pdf ../BGMApp/BGMApp/Images.xcassets/Volume$n.imageset/; done
%
% Might build correctly with regular LaTeX. I haven't tried it.
%

\documentclass[tikz]{standalone}
\usepackage{tikz}
% "dummyOption" prevents "Package optional Warning: No options were selected,
% so all optional text will be printed" when building Volume0.pdf.
\usepackage[dummyOption]{optional}

\begin{document}
\begin{tikzpicture}

% Speaker (Rounded box and triangle)
\fill[rounded corners=5mm]
    (0mm, 62.5mm) rectangle (25mm, 37.5mm) {};
\draw[rounded corners=2.5mm,fill=black]
    (3mm, 50mm)--(34mm, 76.5mm)--(34mm, 23.5mm)--cycle;

% First sound wave (Curved line)
\opt{w1}{
    \draw[line width=4.3mm,line cap=round]
        (44mm, 36.5mm) to[out=46,in=-46] (44mm, 63.5mm);
}

% Second sound wave (Curved line)
\opt{w2}{
    \draw[line width=4.3mm,line cap=round]
        (57.5mm, 27.5mm) to[out=46,in=-46] (57.5mm, 72.5mm);
}

% Third sound wave (Curved line)
\opt{w3}{
    \draw[line width=4.3mm,line cap=round]
        (72mm, 18.5mm) to[out=46,in=-46] (72mm, 81.5mm);
}

% Always draw a transparent copy of the third wave so the images will all have
% the same width.
\draw[line width=4.3mm,line cap=round,opacity=0]
    (72mm, 18.5mm) to[out=46,in=-46] (72mm, 81.5mm);

\end{tikzpicture}
\end{document}

'
%   for n in 0 1 2 3; do mv Volume$n.pdf ../BGMApp/BGMApp/Images.xcassets/Volume$n.imageset/; done
%
% Might build correctly with regular LaTeX. I haven't tried it.
%

\documentclass[tikz]{standalone}
\usepackage{tikz}
% "dummyOption" prevents "Package optional Warning: No options were selected,
% so all optional text will be printed" when building Volume0.pdf.
\usepackage[dummyOption]{optional}

\begin{document}
\begin{tikzpicture}

% Speaker (Rounded box and triangle)
\fill[rounded corners=5mm]
    (0mm, 62.5mm) rectangle (25mm, 37.5mm) {};
\draw[rounded corners=2.5mm,fill=black]
    (3mm, 50mm)--(34mm, 76.5mm)--(34mm, 23.5mm)--cycle;

% First sound wave (Curved line)
\opt{w1}{
    \draw[line width=4.3mm,line cap=round]
        (44mm, 36.5mm) to[out=46,in=-46] (44mm, 63.5mm);
}

% Second sound wave (Curved line)
\opt{w2}{
    \draw[line width=4.3mm,line cap=round]
        (57.5mm, 27.5mm) to[out=46,in=-46] (57.5mm, 72.5mm);
}

% Third sound wave (Curved line)
\opt{w3}{
    \draw[line width=4.3mm,line cap=round]
        (72mm, 18.5mm) to[out=46,in=-46] (72mm, 81.5mm);
}

% Always draw a transparent copy of the third wave so the images will all have
% the same width.
\draw[line width=4.3mm,line cap=round,opacity=0]
    (72mm, 18.5mm) to[out=46,in=-46] (72mm, 81.5mm);

\end{tikzpicture}
\end{document}

'
%   for n in 0 1 2 3; do mv Volume$n.pdf ../BGMApp/BGMApp/Images.xcassets/Volume$n.imageset/; done
%
% Might build correctly with regular LaTeX. I haven't tried it.
%

\documentclass[tikz]{standalone}
\usepackage{tikz}
% "dummyOption" prevents "Package optional Warning: No options were selected,
% so all optional text will be printed" when building Volume0.pdf.
\usepackage[dummyOption]{optional}

\begin{document}
\begin{tikzpicture}

% Speaker (Rounded box and triangle)
\fill[rounded corners=5mm]
    (0mm, 62.5mm) rectangle (25mm, 37.5mm) {};
\draw[rounded corners=2.5mm,fill=black]
    (3mm, 50mm)--(34mm, 76.5mm)--(34mm, 23.5mm)--cycle;

% First sound wave (Curved line)
\opt{w1}{
    \draw[line width=4.3mm,line cap=round]
        (44mm, 36.5mm) to[out=46,in=-46] (44mm, 63.5mm);
}

% Second sound wave (Curved line)
\opt{w2}{
    \draw[line width=4.3mm,line cap=round]
        (57.5mm, 27.5mm) to[out=46,in=-46] (57.5mm, 72.5mm);
}

% Third sound wave (Curved line)
\opt{w3}{
    \draw[line width=4.3mm,line cap=round]
        (72mm, 18.5mm) to[out=46,in=-46] (72mm, 81.5mm);
}

% Always draw a transparent copy of the third wave so the images will all have
% the same width.
\draw[line width=4.3mm,line cap=round,opacity=0]
    (72mm, 18.5mm) to[out=46,in=-46] (72mm, 81.5mm);

\end{tikzpicture}
\end{document}

'
%   xelatex -jobname=Volume3 '\def\UseOption{w1,w2,w3}% This file is part of Background Music.
%
% Background Music is free software: you can redistribute it and/or
% modify it under the terms of the GNU General Public License as
% published by the Free Software Foundation, either version 2 of the
% License, or (at your option) any later version.
%
% Background Music is distributed in the hope that it will be useful,
% but WITHOUT ANY WARRANTY; without even the implied warranty of
% MERCHANTABILITY or FITNESS FOR A PARTICULAR PURPOSE. See the
% GNU General Public License for more details.
%
% You should have received a copy of the GNU General Public License
% along with Background Music. If not, see <http://www.gnu.org/licenses/>.

%
% VolumeIcons.tex
%
% Build with XeTeX:
%   xelatex -jobname=Volume0 '\def\UseOption{}% This file is part of Background Music.
%
% Background Music is free software: you can redistribute it and/or
% modify it under the terms of the GNU General Public License as
% published by the Free Software Foundation, either version 2 of the
% License, or (at your option) any later version.
%
% Background Music is distributed in the hope that it will be useful,
% but WITHOUT ANY WARRANTY; without even the implied warranty of
% MERCHANTABILITY or FITNESS FOR A PARTICULAR PURPOSE. See the
% GNU General Public License for more details.
%
% You should have received a copy of the GNU General Public License
% along with Background Music. If not, see <http://www.gnu.org/licenses/>.

%
% VolumeIcons.tex
%
% Build with XeTeX:
%   xelatex -jobname=Volume0 '\def\UseOption{}% This file is part of Background Music.
%
% Background Music is free software: you can redistribute it and/or
% modify it under the terms of the GNU General Public License as
% published by the Free Software Foundation, either version 2 of the
% License, or (at your option) any later version.
%
% Background Music is distributed in the hope that it will be useful,
% but WITHOUT ANY WARRANTY; without even the implied warranty of
% MERCHANTABILITY or FITNESS FOR A PARTICULAR PURPOSE. See the
% GNU General Public License for more details.
%
% You should have received a copy of the GNU General Public License
% along with Background Music. If not, see <http://www.gnu.org/licenses/>.

%
% VolumeIcons.tex
%
% Build with XeTeX:
%   xelatex -jobname=Volume0 '\def\UseOption{}\input{VolumeIcons.tex}'
%   xelatex -jobname=Volume1 '\def\UseOption{w1}\input{VolumeIcons.tex}'
%   xelatex -jobname=Volume2 '\def\UseOption{w1,w2}\input{VolumeIcons.tex}'
%   xelatex -jobname=Volume3 '\def\UseOption{w1,w2,w3}\input{VolumeIcons.tex}'
%   for n in 0 1 2 3; do mv Volume$n.pdf ../BGMApp/BGMApp/Images.xcassets/Volume$n.imageset/; done
%
% Might build correctly with regular LaTeX. I haven't tried it.
%

\documentclass[tikz]{standalone}
\usepackage{tikz}
% "dummyOption" prevents "Package optional Warning: No options were selected,
% so all optional text will be printed" when building Volume0.pdf.
\usepackage[dummyOption]{optional}

\begin{document}
\begin{tikzpicture}

% Speaker (Rounded box and triangle)
\fill[rounded corners=5mm]
    (0mm, 62.5mm) rectangle (25mm, 37.5mm) {};
\draw[rounded corners=2.5mm,fill=black]
    (3mm, 50mm)--(34mm, 76.5mm)--(34mm, 23.5mm)--cycle;

% First sound wave (Curved line)
\opt{w1}{
    \draw[line width=4.3mm,line cap=round]
        (44mm, 36.5mm) to[out=46,in=-46] (44mm, 63.5mm);
}

% Second sound wave (Curved line)
\opt{w2}{
    \draw[line width=4.3mm,line cap=round]
        (57.5mm, 27.5mm) to[out=46,in=-46] (57.5mm, 72.5mm);
}

% Third sound wave (Curved line)
\opt{w3}{
    \draw[line width=4.3mm,line cap=round]
        (72mm, 18.5mm) to[out=46,in=-46] (72mm, 81.5mm);
}

% Always draw a transparent copy of the third wave so the images will all have
% the same width.
\draw[line width=4.3mm,line cap=round,opacity=0]
    (72mm, 18.5mm) to[out=46,in=-46] (72mm, 81.5mm);

\end{tikzpicture}
\end{document}

'
%   xelatex -jobname=Volume1 '\def\UseOption{w1}% This file is part of Background Music.
%
% Background Music is free software: you can redistribute it and/or
% modify it under the terms of the GNU General Public License as
% published by the Free Software Foundation, either version 2 of the
% License, or (at your option) any later version.
%
% Background Music is distributed in the hope that it will be useful,
% but WITHOUT ANY WARRANTY; without even the implied warranty of
% MERCHANTABILITY or FITNESS FOR A PARTICULAR PURPOSE. See the
% GNU General Public License for more details.
%
% You should have received a copy of the GNU General Public License
% along with Background Music. If not, see <http://www.gnu.org/licenses/>.

%
% VolumeIcons.tex
%
% Build with XeTeX:
%   xelatex -jobname=Volume0 '\def\UseOption{}\input{VolumeIcons.tex}'
%   xelatex -jobname=Volume1 '\def\UseOption{w1}\input{VolumeIcons.tex}'
%   xelatex -jobname=Volume2 '\def\UseOption{w1,w2}\input{VolumeIcons.tex}'
%   xelatex -jobname=Volume3 '\def\UseOption{w1,w2,w3}\input{VolumeIcons.tex}'
%   for n in 0 1 2 3; do mv Volume$n.pdf ../BGMApp/BGMApp/Images.xcassets/Volume$n.imageset/; done
%
% Might build correctly with regular LaTeX. I haven't tried it.
%

\documentclass[tikz]{standalone}
\usepackage{tikz}
% "dummyOption" prevents "Package optional Warning: No options were selected,
% so all optional text will be printed" when building Volume0.pdf.
\usepackage[dummyOption]{optional}

\begin{document}
\begin{tikzpicture}

% Speaker (Rounded box and triangle)
\fill[rounded corners=5mm]
    (0mm, 62.5mm) rectangle (25mm, 37.5mm) {};
\draw[rounded corners=2.5mm,fill=black]
    (3mm, 50mm)--(34mm, 76.5mm)--(34mm, 23.5mm)--cycle;

% First sound wave (Curved line)
\opt{w1}{
    \draw[line width=4.3mm,line cap=round]
        (44mm, 36.5mm) to[out=46,in=-46] (44mm, 63.5mm);
}

% Second sound wave (Curved line)
\opt{w2}{
    \draw[line width=4.3mm,line cap=round]
        (57.5mm, 27.5mm) to[out=46,in=-46] (57.5mm, 72.5mm);
}

% Third sound wave (Curved line)
\opt{w3}{
    \draw[line width=4.3mm,line cap=round]
        (72mm, 18.5mm) to[out=46,in=-46] (72mm, 81.5mm);
}

% Always draw a transparent copy of the third wave so the images will all have
% the same width.
\draw[line width=4.3mm,line cap=round,opacity=0]
    (72mm, 18.5mm) to[out=46,in=-46] (72mm, 81.5mm);

\end{tikzpicture}
\end{document}

'
%   xelatex -jobname=Volume2 '\def\UseOption{w1,w2}% This file is part of Background Music.
%
% Background Music is free software: you can redistribute it and/or
% modify it under the terms of the GNU General Public License as
% published by the Free Software Foundation, either version 2 of the
% License, or (at your option) any later version.
%
% Background Music is distributed in the hope that it will be useful,
% but WITHOUT ANY WARRANTY; without even the implied warranty of
% MERCHANTABILITY or FITNESS FOR A PARTICULAR PURPOSE. See the
% GNU General Public License for more details.
%
% You should have received a copy of the GNU General Public License
% along with Background Music. If not, see <http://www.gnu.org/licenses/>.

%
% VolumeIcons.tex
%
% Build with XeTeX:
%   xelatex -jobname=Volume0 '\def\UseOption{}\input{VolumeIcons.tex}'
%   xelatex -jobname=Volume1 '\def\UseOption{w1}\input{VolumeIcons.tex}'
%   xelatex -jobname=Volume2 '\def\UseOption{w1,w2}\input{VolumeIcons.tex}'
%   xelatex -jobname=Volume3 '\def\UseOption{w1,w2,w3}\input{VolumeIcons.tex}'
%   for n in 0 1 2 3; do mv Volume$n.pdf ../BGMApp/BGMApp/Images.xcassets/Volume$n.imageset/; done
%
% Might build correctly with regular LaTeX. I haven't tried it.
%

\documentclass[tikz]{standalone}
\usepackage{tikz}
% "dummyOption" prevents "Package optional Warning: No options were selected,
% so all optional text will be printed" when building Volume0.pdf.
\usepackage[dummyOption]{optional}

\begin{document}
\begin{tikzpicture}

% Speaker (Rounded box and triangle)
\fill[rounded corners=5mm]
    (0mm, 62.5mm) rectangle (25mm, 37.5mm) {};
\draw[rounded corners=2.5mm,fill=black]
    (3mm, 50mm)--(34mm, 76.5mm)--(34mm, 23.5mm)--cycle;

% First sound wave (Curved line)
\opt{w1}{
    \draw[line width=4.3mm,line cap=round]
        (44mm, 36.5mm) to[out=46,in=-46] (44mm, 63.5mm);
}

% Second sound wave (Curved line)
\opt{w2}{
    \draw[line width=4.3mm,line cap=round]
        (57.5mm, 27.5mm) to[out=46,in=-46] (57.5mm, 72.5mm);
}

% Third sound wave (Curved line)
\opt{w3}{
    \draw[line width=4.3mm,line cap=round]
        (72mm, 18.5mm) to[out=46,in=-46] (72mm, 81.5mm);
}

% Always draw a transparent copy of the third wave so the images will all have
% the same width.
\draw[line width=4.3mm,line cap=round,opacity=0]
    (72mm, 18.5mm) to[out=46,in=-46] (72mm, 81.5mm);

\end{tikzpicture}
\end{document}

'
%   xelatex -jobname=Volume3 '\def\UseOption{w1,w2,w3}% This file is part of Background Music.
%
% Background Music is free software: you can redistribute it and/or
% modify it under the terms of the GNU General Public License as
% published by the Free Software Foundation, either version 2 of the
% License, or (at your option) any later version.
%
% Background Music is distributed in the hope that it will be useful,
% but WITHOUT ANY WARRANTY; without even the implied warranty of
% MERCHANTABILITY or FITNESS FOR A PARTICULAR PURPOSE. See the
% GNU General Public License for more details.
%
% You should have received a copy of the GNU General Public License
% along with Background Music. If not, see <http://www.gnu.org/licenses/>.

%
% VolumeIcons.tex
%
% Build with XeTeX:
%   xelatex -jobname=Volume0 '\def\UseOption{}\input{VolumeIcons.tex}'
%   xelatex -jobname=Volume1 '\def\UseOption{w1}\input{VolumeIcons.tex}'
%   xelatex -jobname=Volume2 '\def\UseOption{w1,w2}\input{VolumeIcons.tex}'
%   xelatex -jobname=Volume3 '\def\UseOption{w1,w2,w3}\input{VolumeIcons.tex}'
%   for n in 0 1 2 3; do mv Volume$n.pdf ../BGMApp/BGMApp/Images.xcassets/Volume$n.imageset/; done
%
% Might build correctly with regular LaTeX. I haven't tried it.
%

\documentclass[tikz]{standalone}
\usepackage{tikz}
% "dummyOption" prevents "Package optional Warning: No options were selected,
% so all optional text will be printed" when building Volume0.pdf.
\usepackage[dummyOption]{optional}

\begin{document}
\begin{tikzpicture}

% Speaker (Rounded box and triangle)
\fill[rounded corners=5mm]
    (0mm, 62.5mm) rectangle (25mm, 37.5mm) {};
\draw[rounded corners=2.5mm,fill=black]
    (3mm, 50mm)--(34mm, 76.5mm)--(34mm, 23.5mm)--cycle;

% First sound wave (Curved line)
\opt{w1}{
    \draw[line width=4.3mm,line cap=round]
        (44mm, 36.5mm) to[out=46,in=-46] (44mm, 63.5mm);
}

% Second sound wave (Curved line)
\opt{w2}{
    \draw[line width=4.3mm,line cap=round]
        (57.5mm, 27.5mm) to[out=46,in=-46] (57.5mm, 72.5mm);
}

% Third sound wave (Curved line)
\opt{w3}{
    \draw[line width=4.3mm,line cap=round]
        (72mm, 18.5mm) to[out=46,in=-46] (72mm, 81.5mm);
}

% Always draw a transparent copy of the third wave so the images will all have
% the same width.
\draw[line width=4.3mm,line cap=round,opacity=0]
    (72mm, 18.5mm) to[out=46,in=-46] (72mm, 81.5mm);

\end{tikzpicture}
\end{document}

'
%   for n in 0 1 2 3; do mv Volume$n.pdf ../BGMApp/BGMApp/Images.xcassets/Volume$n.imageset/; done
%
% Might build correctly with regular LaTeX. I haven't tried it.
%

\documentclass[tikz]{standalone}
\usepackage{tikz}
% "dummyOption" prevents "Package optional Warning: No options were selected,
% so all optional text will be printed" when building Volume0.pdf.
\usepackage[dummyOption]{optional}

\begin{document}
\begin{tikzpicture}

% Speaker (Rounded box and triangle)
\fill[rounded corners=5mm]
    (0mm, 62.5mm) rectangle (25mm, 37.5mm) {};
\draw[rounded corners=2.5mm,fill=black]
    (3mm, 50mm)--(34mm, 76.5mm)--(34mm, 23.5mm)--cycle;

% First sound wave (Curved line)
\opt{w1}{
    \draw[line width=4.3mm,line cap=round]
        (44mm, 36.5mm) to[out=46,in=-46] (44mm, 63.5mm);
}

% Second sound wave (Curved line)
\opt{w2}{
    \draw[line width=4.3mm,line cap=round]
        (57.5mm, 27.5mm) to[out=46,in=-46] (57.5mm, 72.5mm);
}

% Third sound wave (Curved line)
\opt{w3}{
    \draw[line width=4.3mm,line cap=round]
        (72mm, 18.5mm) to[out=46,in=-46] (72mm, 81.5mm);
}

% Always draw a transparent copy of the third wave so the images will all have
% the same width.
\draw[line width=4.3mm,line cap=round,opacity=0]
    (72mm, 18.5mm) to[out=46,in=-46] (72mm, 81.5mm);

\end{tikzpicture}
\end{document}

'
%   xelatex -jobname=Volume1 '\def\UseOption{w1}% This file is part of Background Music.
%
% Background Music is free software: you can redistribute it and/or
% modify it under the terms of the GNU General Public License as
% published by the Free Software Foundation, either version 2 of the
% License, or (at your option) any later version.
%
% Background Music is distributed in the hope that it will be useful,
% but WITHOUT ANY WARRANTY; without even the implied warranty of
% MERCHANTABILITY or FITNESS FOR A PARTICULAR PURPOSE. See the
% GNU General Public License for more details.
%
% You should have received a copy of the GNU General Public License
% along with Background Music. If not, see <http://www.gnu.org/licenses/>.

%
% VolumeIcons.tex
%
% Build with XeTeX:
%   xelatex -jobname=Volume0 '\def\UseOption{}% This file is part of Background Music.
%
% Background Music is free software: you can redistribute it and/or
% modify it under the terms of the GNU General Public License as
% published by the Free Software Foundation, either version 2 of the
% License, or (at your option) any later version.
%
% Background Music is distributed in the hope that it will be useful,
% but WITHOUT ANY WARRANTY; without even the implied warranty of
% MERCHANTABILITY or FITNESS FOR A PARTICULAR PURPOSE. See the
% GNU General Public License for more details.
%
% You should have received a copy of the GNU General Public License
% along with Background Music. If not, see <http://www.gnu.org/licenses/>.

%
% VolumeIcons.tex
%
% Build with XeTeX:
%   xelatex -jobname=Volume0 '\def\UseOption{}\input{VolumeIcons.tex}'
%   xelatex -jobname=Volume1 '\def\UseOption{w1}\input{VolumeIcons.tex}'
%   xelatex -jobname=Volume2 '\def\UseOption{w1,w2}\input{VolumeIcons.tex}'
%   xelatex -jobname=Volume3 '\def\UseOption{w1,w2,w3}\input{VolumeIcons.tex}'
%   for n in 0 1 2 3; do mv Volume$n.pdf ../BGMApp/BGMApp/Images.xcassets/Volume$n.imageset/; done
%
% Might build correctly with regular LaTeX. I haven't tried it.
%

\documentclass[tikz]{standalone}
\usepackage{tikz}
% "dummyOption" prevents "Package optional Warning: No options were selected,
% so all optional text will be printed" when building Volume0.pdf.
\usepackage[dummyOption]{optional}

\begin{document}
\begin{tikzpicture}

% Speaker (Rounded box and triangle)
\fill[rounded corners=5mm]
    (0mm, 62.5mm) rectangle (25mm, 37.5mm) {};
\draw[rounded corners=2.5mm,fill=black]
    (3mm, 50mm)--(34mm, 76.5mm)--(34mm, 23.5mm)--cycle;

% First sound wave (Curved line)
\opt{w1}{
    \draw[line width=4.3mm,line cap=round]
        (44mm, 36.5mm) to[out=46,in=-46] (44mm, 63.5mm);
}

% Second sound wave (Curved line)
\opt{w2}{
    \draw[line width=4.3mm,line cap=round]
        (57.5mm, 27.5mm) to[out=46,in=-46] (57.5mm, 72.5mm);
}

% Third sound wave (Curved line)
\opt{w3}{
    \draw[line width=4.3mm,line cap=round]
        (72mm, 18.5mm) to[out=46,in=-46] (72mm, 81.5mm);
}

% Always draw a transparent copy of the third wave so the images will all have
% the same width.
\draw[line width=4.3mm,line cap=round,opacity=0]
    (72mm, 18.5mm) to[out=46,in=-46] (72mm, 81.5mm);

\end{tikzpicture}
\end{document}

'
%   xelatex -jobname=Volume1 '\def\UseOption{w1}% This file is part of Background Music.
%
% Background Music is free software: you can redistribute it and/or
% modify it under the terms of the GNU General Public License as
% published by the Free Software Foundation, either version 2 of the
% License, or (at your option) any later version.
%
% Background Music is distributed in the hope that it will be useful,
% but WITHOUT ANY WARRANTY; without even the implied warranty of
% MERCHANTABILITY or FITNESS FOR A PARTICULAR PURPOSE. See the
% GNU General Public License for more details.
%
% You should have received a copy of the GNU General Public License
% along with Background Music. If not, see <http://www.gnu.org/licenses/>.

%
% VolumeIcons.tex
%
% Build with XeTeX:
%   xelatex -jobname=Volume0 '\def\UseOption{}\input{VolumeIcons.tex}'
%   xelatex -jobname=Volume1 '\def\UseOption{w1}\input{VolumeIcons.tex}'
%   xelatex -jobname=Volume2 '\def\UseOption{w1,w2}\input{VolumeIcons.tex}'
%   xelatex -jobname=Volume3 '\def\UseOption{w1,w2,w3}\input{VolumeIcons.tex}'
%   for n in 0 1 2 3; do mv Volume$n.pdf ../BGMApp/BGMApp/Images.xcassets/Volume$n.imageset/; done
%
% Might build correctly with regular LaTeX. I haven't tried it.
%

\documentclass[tikz]{standalone}
\usepackage{tikz}
% "dummyOption" prevents "Package optional Warning: No options were selected,
% so all optional text will be printed" when building Volume0.pdf.
\usepackage[dummyOption]{optional}

\begin{document}
\begin{tikzpicture}

% Speaker (Rounded box and triangle)
\fill[rounded corners=5mm]
    (0mm, 62.5mm) rectangle (25mm, 37.5mm) {};
\draw[rounded corners=2.5mm,fill=black]
    (3mm, 50mm)--(34mm, 76.5mm)--(34mm, 23.5mm)--cycle;

% First sound wave (Curved line)
\opt{w1}{
    \draw[line width=4.3mm,line cap=round]
        (44mm, 36.5mm) to[out=46,in=-46] (44mm, 63.5mm);
}

% Second sound wave (Curved line)
\opt{w2}{
    \draw[line width=4.3mm,line cap=round]
        (57.5mm, 27.5mm) to[out=46,in=-46] (57.5mm, 72.5mm);
}

% Third sound wave (Curved line)
\opt{w3}{
    \draw[line width=4.3mm,line cap=round]
        (72mm, 18.5mm) to[out=46,in=-46] (72mm, 81.5mm);
}

% Always draw a transparent copy of the third wave so the images will all have
% the same width.
\draw[line width=4.3mm,line cap=round,opacity=0]
    (72mm, 18.5mm) to[out=46,in=-46] (72mm, 81.5mm);

\end{tikzpicture}
\end{document}

'
%   xelatex -jobname=Volume2 '\def\UseOption{w1,w2}% This file is part of Background Music.
%
% Background Music is free software: you can redistribute it and/or
% modify it under the terms of the GNU General Public License as
% published by the Free Software Foundation, either version 2 of the
% License, or (at your option) any later version.
%
% Background Music is distributed in the hope that it will be useful,
% but WITHOUT ANY WARRANTY; without even the implied warranty of
% MERCHANTABILITY or FITNESS FOR A PARTICULAR PURPOSE. See the
% GNU General Public License for more details.
%
% You should have received a copy of the GNU General Public License
% along with Background Music. If not, see <http://www.gnu.org/licenses/>.

%
% VolumeIcons.tex
%
% Build with XeTeX:
%   xelatex -jobname=Volume0 '\def\UseOption{}\input{VolumeIcons.tex}'
%   xelatex -jobname=Volume1 '\def\UseOption{w1}\input{VolumeIcons.tex}'
%   xelatex -jobname=Volume2 '\def\UseOption{w1,w2}\input{VolumeIcons.tex}'
%   xelatex -jobname=Volume3 '\def\UseOption{w1,w2,w3}\input{VolumeIcons.tex}'
%   for n in 0 1 2 3; do mv Volume$n.pdf ../BGMApp/BGMApp/Images.xcassets/Volume$n.imageset/; done
%
% Might build correctly with regular LaTeX. I haven't tried it.
%

\documentclass[tikz]{standalone}
\usepackage{tikz}
% "dummyOption" prevents "Package optional Warning: No options were selected,
% so all optional text will be printed" when building Volume0.pdf.
\usepackage[dummyOption]{optional}

\begin{document}
\begin{tikzpicture}

% Speaker (Rounded box and triangle)
\fill[rounded corners=5mm]
    (0mm, 62.5mm) rectangle (25mm, 37.5mm) {};
\draw[rounded corners=2.5mm,fill=black]
    (3mm, 50mm)--(34mm, 76.5mm)--(34mm, 23.5mm)--cycle;

% First sound wave (Curved line)
\opt{w1}{
    \draw[line width=4.3mm,line cap=round]
        (44mm, 36.5mm) to[out=46,in=-46] (44mm, 63.5mm);
}

% Second sound wave (Curved line)
\opt{w2}{
    \draw[line width=4.3mm,line cap=round]
        (57.5mm, 27.5mm) to[out=46,in=-46] (57.5mm, 72.5mm);
}

% Third sound wave (Curved line)
\opt{w3}{
    \draw[line width=4.3mm,line cap=round]
        (72mm, 18.5mm) to[out=46,in=-46] (72mm, 81.5mm);
}

% Always draw a transparent copy of the third wave so the images will all have
% the same width.
\draw[line width=4.3mm,line cap=round,opacity=0]
    (72mm, 18.5mm) to[out=46,in=-46] (72mm, 81.5mm);

\end{tikzpicture}
\end{document}

'
%   xelatex -jobname=Volume3 '\def\UseOption{w1,w2,w3}% This file is part of Background Music.
%
% Background Music is free software: you can redistribute it and/or
% modify it under the terms of the GNU General Public License as
% published by the Free Software Foundation, either version 2 of the
% License, or (at your option) any later version.
%
% Background Music is distributed in the hope that it will be useful,
% but WITHOUT ANY WARRANTY; without even the implied warranty of
% MERCHANTABILITY or FITNESS FOR A PARTICULAR PURPOSE. See the
% GNU General Public License for more details.
%
% You should have received a copy of the GNU General Public License
% along with Background Music. If not, see <http://www.gnu.org/licenses/>.

%
% VolumeIcons.tex
%
% Build with XeTeX:
%   xelatex -jobname=Volume0 '\def\UseOption{}\input{VolumeIcons.tex}'
%   xelatex -jobname=Volume1 '\def\UseOption{w1}\input{VolumeIcons.tex}'
%   xelatex -jobname=Volume2 '\def\UseOption{w1,w2}\input{VolumeIcons.tex}'
%   xelatex -jobname=Volume3 '\def\UseOption{w1,w2,w3}\input{VolumeIcons.tex}'
%   for n in 0 1 2 3; do mv Volume$n.pdf ../BGMApp/BGMApp/Images.xcassets/Volume$n.imageset/; done
%
% Might build correctly with regular LaTeX. I haven't tried it.
%

\documentclass[tikz]{standalone}
\usepackage{tikz}
% "dummyOption" prevents "Package optional Warning: No options were selected,
% so all optional text will be printed" when building Volume0.pdf.
\usepackage[dummyOption]{optional}

\begin{document}
\begin{tikzpicture}

% Speaker (Rounded box and triangle)
\fill[rounded corners=5mm]
    (0mm, 62.5mm) rectangle (25mm, 37.5mm) {};
\draw[rounded corners=2.5mm,fill=black]
    (3mm, 50mm)--(34mm, 76.5mm)--(34mm, 23.5mm)--cycle;

% First sound wave (Curved line)
\opt{w1}{
    \draw[line width=4.3mm,line cap=round]
        (44mm, 36.5mm) to[out=46,in=-46] (44mm, 63.5mm);
}

% Second sound wave (Curved line)
\opt{w2}{
    \draw[line width=4.3mm,line cap=round]
        (57.5mm, 27.5mm) to[out=46,in=-46] (57.5mm, 72.5mm);
}

% Third sound wave (Curved line)
\opt{w3}{
    \draw[line width=4.3mm,line cap=round]
        (72mm, 18.5mm) to[out=46,in=-46] (72mm, 81.5mm);
}

% Always draw a transparent copy of the third wave so the images will all have
% the same width.
\draw[line width=4.3mm,line cap=round,opacity=0]
    (72mm, 18.5mm) to[out=46,in=-46] (72mm, 81.5mm);

\end{tikzpicture}
\end{document}

'
%   for n in 0 1 2 3; do mv Volume$n.pdf ../BGMApp/BGMApp/Images.xcassets/Volume$n.imageset/; done
%
% Might build correctly with regular LaTeX. I haven't tried it.
%

\documentclass[tikz]{standalone}
\usepackage{tikz}
% "dummyOption" prevents "Package optional Warning: No options were selected,
% so all optional text will be printed" when building Volume0.pdf.
\usepackage[dummyOption]{optional}

\begin{document}
\begin{tikzpicture}

% Speaker (Rounded box and triangle)
\fill[rounded corners=5mm]
    (0mm, 62.5mm) rectangle (25mm, 37.5mm) {};
\draw[rounded corners=2.5mm,fill=black]
    (3mm, 50mm)--(34mm, 76.5mm)--(34mm, 23.5mm)--cycle;

% First sound wave (Curved line)
\opt{w1}{
    \draw[line width=4.3mm,line cap=round]
        (44mm, 36.5mm) to[out=46,in=-46] (44mm, 63.5mm);
}

% Second sound wave (Curved line)
\opt{w2}{
    \draw[line width=4.3mm,line cap=round]
        (57.5mm, 27.5mm) to[out=46,in=-46] (57.5mm, 72.5mm);
}

% Third sound wave (Curved line)
\opt{w3}{
    \draw[line width=4.3mm,line cap=round]
        (72mm, 18.5mm) to[out=46,in=-46] (72mm, 81.5mm);
}

% Always draw a transparent copy of the third wave so the images will all have
% the same width.
\draw[line width=4.3mm,line cap=round,opacity=0]
    (72mm, 18.5mm) to[out=46,in=-46] (72mm, 81.5mm);

\end{tikzpicture}
\end{document}

'
%   xelatex -jobname=Volume2 '\def\UseOption{w1,w2}% This file is part of Background Music.
%
% Background Music is free software: you can redistribute it and/or
% modify it under the terms of the GNU General Public License as
% published by the Free Software Foundation, either version 2 of the
% License, or (at your option) any later version.
%
% Background Music is distributed in the hope that it will be useful,
% but WITHOUT ANY WARRANTY; without even the implied warranty of
% MERCHANTABILITY or FITNESS FOR A PARTICULAR PURPOSE. See the
% GNU General Public License for more details.
%
% You should have received a copy of the GNU General Public License
% along with Background Music. If not, see <http://www.gnu.org/licenses/>.

%
% VolumeIcons.tex
%
% Build with XeTeX:
%   xelatex -jobname=Volume0 '\def\UseOption{}% This file is part of Background Music.
%
% Background Music is free software: you can redistribute it and/or
% modify it under the terms of the GNU General Public License as
% published by the Free Software Foundation, either version 2 of the
% License, or (at your option) any later version.
%
% Background Music is distributed in the hope that it will be useful,
% but WITHOUT ANY WARRANTY; without even the implied warranty of
% MERCHANTABILITY or FITNESS FOR A PARTICULAR PURPOSE. See the
% GNU General Public License for more details.
%
% You should have received a copy of the GNU General Public License
% along with Background Music. If not, see <http://www.gnu.org/licenses/>.

%
% VolumeIcons.tex
%
% Build with XeTeX:
%   xelatex -jobname=Volume0 '\def\UseOption{}\input{VolumeIcons.tex}'
%   xelatex -jobname=Volume1 '\def\UseOption{w1}\input{VolumeIcons.tex}'
%   xelatex -jobname=Volume2 '\def\UseOption{w1,w2}\input{VolumeIcons.tex}'
%   xelatex -jobname=Volume3 '\def\UseOption{w1,w2,w3}\input{VolumeIcons.tex}'
%   for n in 0 1 2 3; do mv Volume$n.pdf ../BGMApp/BGMApp/Images.xcassets/Volume$n.imageset/; done
%
% Might build correctly with regular LaTeX. I haven't tried it.
%

\documentclass[tikz]{standalone}
\usepackage{tikz}
% "dummyOption" prevents "Package optional Warning: No options were selected,
% so all optional text will be printed" when building Volume0.pdf.
\usepackage[dummyOption]{optional}

\begin{document}
\begin{tikzpicture}

% Speaker (Rounded box and triangle)
\fill[rounded corners=5mm]
    (0mm, 62.5mm) rectangle (25mm, 37.5mm) {};
\draw[rounded corners=2.5mm,fill=black]
    (3mm, 50mm)--(34mm, 76.5mm)--(34mm, 23.5mm)--cycle;

% First sound wave (Curved line)
\opt{w1}{
    \draw[line width=4.3mm,line cap=round]
        (44mm, 36.5mm) to[out=46,in=-46] (44mm, 63.5mm);
}

% Second sound wave (Curved line)
\opt{w2}{
    \draw[line width=4.3mm,line cap=round]
        (57.5mm, 27.5mm) to[out=46,in=-46] (57.5mm, 72.5mm);
}

% Third sound wave (Curved line)
\opt{w3}{
    \draw[line width=4.3mm,line cap=round]
        (72mm, 18.5mm) to[out=46,in=-46] (72mm, 81.5mm);
}

% Always draw a transparent copy of the third wave so the images will all have
% the same width.
\draw[line width=4.3mm,line cap=round,opacity=0]
    (72mm, 18.5mm) to[out=46,in=-46] (72mm, 81.5mm);

\end{tikzpicture}
\end{document}

'
%   xelatex -jobname=Volume1 '\def\UseOption{w1}% This file is part of Background Music.
%
% Background Music is free software: you can redistribute it and/or
% modify it under the terms of the GNU General Public License as
% published by the Free Software Foundation, either version 2 of the
% License, or (at your option) any later version.
%
% Background Music is distributed in the hope that it will be useful,
% but WITHOUT ANY WARRANTY; without even the implied warranty of
% MERCHANTABILITY or FITNESS FOR A PARTICULAR PURPOSE. See the
% GNU General Public License for more details.
%
% You should have received a copy of the GNU General Public License
% along with Background Music. If not, see <http://www.gnu.org/licenses/>.

%
% VolumeIcons.tex
%
% Build with XeTeX:
%   xelatex -jobname=Volume0 '\def\UseOption{}\input{VolumeIcons.tex}'
%   xelatex -jobname=Volume1 '\def\UseOption{w1}\input{VolumeIcons.tex}'
%   xelatex -jobname=Volume2 '\def\UseOption{w1,w2}\input{VolumeIcons.tex}'
%   xelatex -jobname=Volume3 '\def\UseOption{w1,w2,w3}\input{VolumeIcons.tex}'
%   for n in 0 1 2 3; do mv Volume$n.pdf ../BGMApp/BGMApp/Images.xcassets/Volume$n.imageset/; done
%
% Might build correctly with regular LaTeX. I haven't tried it.
%

\documentclass[tikz]{standalone}
\usepackage{tikz}
% "dummyOption" prevents "Package optional Warning: No options were selected,
% so all optional text will be printed" when building Volume0.pdf.
\usepackage[dummyOption]{optional}

\begin{document}
\begin{tikzpicture}

% Speaker (Rounded box and triangle)
\fill[rounded corners=5mm]
    (0mm, 62.5mm) rectangle (25mm, 37.5mm) {};
\draw[rounded corners=2.5mm,fill=black]
    (3mm, 50mm)--(34mm, 76.5mm)--(34mm, 23.5mm)--cycle;

% First sound wave (Curved line)
\opt{w1}{
    \draw[line width=4.3mm,line cap=round]
        (44mm, 36.5mm) to[out=46,in=-46] (44mm, 63.5mm);
}

% Second sound wave (Curved line)
\opt{w2}{
    \draw[line width=4.3mm,line cap=round]
        (57.5mm, 27.5mm) to[out=46,in=-46] (57.5mm, 72.5mm);
}

% Third sound wave (Curved line)
\opt{w3}{
    \draw[line width=4.3mm,line cap=round]
        (72mm, 18.5mm) to[out=46,in=-46] (72mm, 81.5mm);
}

% Always draw a transparent copy of the third wave so the images will all have
% the same width.
\draw[line width=4.3mm,line cap=round,opacity=0]
    (72mm, 18.5mm) to[out=46,in=-46] (72mm, 81.5mm);

\end{tikzpicture}
\end{document}

'
%   xelatex -jobname=Volume2 '\def\UseOption{w1,w2}% This file is part of Background Music.
%
% Background Music is free software: you can redistribute it and/or
% modify it under the terms of the GNU General Public License as
% published by the Free Software Foundation, either version 2 of the
% License, or (at your option) any later version.
%
% Background Music is distributed in the hope that it will be useful,
% but WITHOUT ANY WARRANTY; without even the implied warranty of
% MERCHANTABILITY or FITNESS FOR A PARTICULAR PURPOSE. See the
% GNU General Public License for more details.
%
% You should have received a copy of the GNU General Public License
% along with Background Music. If not, see <http://www.gnu.org/licenses/>.

%
% VolumeIcons.tex
%
% Build with XeTeX:
%   xelatex -jobname=Volume0 '\def\UseOption{}\input{VolumeIcons.tex}'
%   xelatex -jobname=Volume1 '\def\UseOption{w1}\input{VolumeIcons.tex}'
%   xelatex -jobname=Volume2 '\def\UseOption{w1,w2}\input{VolumeIcons.tex}'
%   xelatex -jobname=Volume3 '\def\UseOption{w1,w2,w3}\input{VolumeIcons.tex}'
%   for n in 0 1 2 3; do mv Volume$n.pdf ../BGMApp/BGMApp/Images.xcassets/Volume$n.imageset/; done
%
% Might build correctly with regular LaTeX. I haven't tried it.
%

\documentclass[tikz]{standalone}
\usepackage{tikz}
% "dummyOption" prevents "Package optional Warning: No options were selected,
% so all optional text will be printed" when building Volume0.pdf.
\usepackage[dummyOption]{optional}

\begin{document}
\begin{tikzpicture}

% Speaker (Rounded box and triangle)
\fill[rounded corners=5mm]
    (0mm, 62.5mm) rectangle (25mm, 37.5mm) {};
\draw[rounded corners=2.5mm,fill=black]
    (3mm, 50mm)--(34mm, 76.5mm)--(34mm, 23.5mm)--cycle;

% First sound wave (Curved line)
\opt{w1}{
    \draw[line width=4.3mm,line cap=round]
        (44mm, 36.5mm) to[out=46,in=-46] (44mm, 63.5mm);
}

% Second sound wave (Curved line)
\opt{w2}{
    \draw[line width=4.3mm,line cap=round]
        (57.5mm, 27.5mm) to[out=46,in=-46] (57.5mm, 72.5mm);
}

% Third sound wave (Curved line)
\opt{w3}{
    \draw[line width=4.3mm,line cap=round]
        (72mm, 18.5mm) to[out=46,in=-46] (72mm, 81.5mm);
}

% Always draw a transparent copy of the third wave so the images will all have
% the same width.
\draw[line width=4.3mm,line cap=round,opacity=0]
    (72mm, 18.5mm) to[out=46,in=-46] (72mm, 81.5mm);

\end{tikzpicture}
\end{document}

'
%   xelatex -jobname=Volume3 '\def\UseOption{w1,w2,w3}% This file is part of Background Music.
%
% Background Music is free software: you can redistribute it and/or
% modify it under the terms of the GNU General Public License as
% published by the Free Software Foundation, either version 2 of the
% License, or (at your option) any later version.
%
% Background Music is distributed in the hope that it will be useful,
% but WITHOUT ANY WARRANTY; without even the implied warranty of
% MERCHANTABILITY or FITNESS FOR A PARTICULAR PURPOSE. See the
% GNU General Public License for more details.
%
% You should have received a copy of the GNU General Public License
% along with Background Music. If not, see <http://www.gnu.org/licenses/>.

%
% VolumeIcons.tex
%
% Build with XeTeX:
%   xelatex -jobname=Volume0 '\def\UseOption{}\input{VolumeIcons.tex}'
%   xelatex -jobname=Volume1 '\def\UseOption{w1}\input{VolumeIcons.tex}'
%   xelatex -jobname=Volume2 '\def\UseOption{w1,w2}\input{VolumeIcons.tex}'
%   xelatex -jobname=Volume3 '\def\UseOption{w1,w2,w3}\input{VolumeIcons.tex}'
%   for n in 0 1 2 3; do mv Volume$n.pdf ../BGMApp/BGMApp/Images.xcassets/Volume$n.imageset/; done
%
% Might build correctly with regular LaTeX. I haven't tried it.
%

\documentclass[tikz]{standalone}
\usepackage{tikz}
% "dummyOption" prevents "Package optional Warning: No options were selected,
% so all optional text will be printed" when building Volume0.pdf.
\usepackage[dummyOption]{optional}

\begin{document}
\begin{tikzpicture}

% Speaker (Rounded box and triangle)
\fill[rounded corners=5mm]
    (0mm, 62.5mm) rectangle (25mm, 37.5mm) {};
\draw[rounded corners=2.5mm,fill=black]
    (3mm, 50mm)--(34mm, 76.5mm)--(34mm, 23.5mm)--cycle;

% First sound wave (Curved line)
\opt{w1}{
    \draw[line width=4.3mm,line cap=round]
        (44mm, 36.5mm) to[out=46,in=-46] (44mm, 63.5mm);
}

% Second sound wave (Curved line)
\opt{w2}{
    \draw[line width=4.3mm,line cap=round]
        (57.5mm, 27.5mm) to[out=46,in=-46] (57.5mm, 72.5mm);
}

% Third sound wave (Curved line)
\opt{w3}{
    \draw[line width=4.3mm,line cap=round]
        (72mm, 18.5mm) to[out=46,in=-46] (72mm, 81.5mm);
}

% Always draw a transparent copy of the third wave so the images will all have
% the same width.
\draw[line width=4.3mm,line cap=round,opacity=0]
    (72mm, 18.5mm) to[out=46,in=-46] (72mm, 81.5mm);

\end{tikzpicture}
\end{document}

'
%   for n in 0 1 2 3; do mv Volume$n.pdf ../BGMApp/BGMApp/Images.xcassets/Volume$n.imageset/; done
%
% Might build correctly with regular LaTeX. I haven't tried it.
%

\documentclass[tikz]{standalone}
\usepackage{tikz}
% "dummyOption" prevents "Package optional Warning: No options were selected,
% so all optional text will be printed" when building Volume0.pdf.
\usepackage[dummyOption]{optional}

\begin{document}
\begin{tikzpicture}

% Speaker (Rounded box and triangle)
\fill[rounded corners=5mm]
    (0mm, 62.5mm) rectangle (25mm, 37.5mm) {};
\draw[rounded corners=2.5mm,fill=black]
    (3mm, 50mm)--(34mm, 76.5mm)--(34mm, 23.5mm)--cycle;

% First sound wave (Curved line)
\opt{w1}{
    \draw[line width=4.3mm,line cap=round]
        (44mm, 36.5mm) to[out=46,in=-46] (44mm, 63.5mm);
}

% Second sound wave (Curved line)
\opt{w2}{
    \draw[line width=4.3mm,line cap=round]
        (57.5mm, 27.5mm) to[out=46,in=-46] (57.5mm, 72.5mm);
}

% Third sound wave (Curved line)
\opt{w3}{
    \draw[line width=4.3mm,line cap=round]
        (72mm, 18.5mm) to[out=46,in=-46] (72mm, 81.5mm);
}

% Always draw a transparent copy of the third wave so the images will all have
% the same width.
\draw[line width=4.3mm,line cap=round,opacity=0]
    (72mm, 18.5mm) to[out=46,in=-46] (72mm, 81.5mm);

\end{tikzpicture}
\end{document}

'
%   xelatex -jobname=Volume3 '\def\UseOption{w1,w2,w3}% This file is part of Background Music.
%
% Background Music is free software: you can redistribute it and/or
% modify it under the terms of the GNU General Public License as
% published by the Free Software Foundation, either version 2 of the
% License, or (at your option) any later version.
%
% Background Music is distributed in the hope that it will be useful,
% but WITHOUT ANY WARRANTY; without even the implied warranty of
% MERCHANTABILITY or FITNESS FOR A PARTICULAR PURPOSE. See the
% GNU General Public License for more details.
%
% You should have received a copy of the GNU General Public License
% along with Background Music. If not, see <http://www.gnu.org/licenses/>.

%
% VolumeIcons.tex
%
% Build with XeTeX:
%   xelatex -jobname=Volume0 '\def\UseOption{}% This file is part of Background Music.
%
% Background Music is free software: you can redistribute it and/or
% modify it under the terms of the GNU General Public License as
% published by the Free Software Foundation, either version 2 of the
% License, or (at your option) any later version.
%
% Background Music is distributed in the hope that it will be useful,
% but WITHOUT ANY WARRANTY; without even the implied warranty of
% MERCHANTABILITY or FITNESS FOR A PARTICULAR PURPOSE. See the
% GNU General Public License for more details.
%
% You should have received a copy of the GNU General Public License
% along with Background Music. If not, see <http://www.gnu.org/licenses/>.

%
% VolumeIcons.tex
%
% Build with XeTeX:
%   xelatex -jobname=Volume0 '\def\UseOption{}\input{VolumeIcons.tex}'
%   xelatex -jobname=Volume1 '\def\UseOption{w1}\input{VolumeIcons.tex}'
%   xelatex -jobname=Volume2 '\def\UseOption{w1,w2}\input{VolumeIcons.tex}'
%   xelatex -jobname=Volume3 '\def\UseOption{w1,w2,w3}\input{VolumeIcons.tex}'
%   for n in 0 1 2 3; do mv Volume$n.pdf ../BGMApp/BGMApp/Images.xcassets/Volume$n.imageset/; done
%
% Might build correctly with regular LaTeX. I haven't tried it.
%

\documentclass[tikz]{standalone}
\usepackage{tikz}
% "dummyOption" prevents "Package optional Warning: No options were selected,
% so all optional text will be printed" when building Volume0.pdf.
\usepackage[dummyOption]{optional}

\begin{document}
\begin{tikzpicture}

% Speaker (Rounded box and triangle)
\fill[rounded corners=5mm]
    (0mm, 62.5mm) rectangle (25mm, 37.5mm) {};
\draw[rounded corners=2.5mm,fill=black]
    (3mm, 50mm)--(34mm, 76.5mm)--(34mm, 23.5mm)--cycle;

% First sound wave (Curved line)
\opt{w1}{
    \draw[line width=4.3mm,line cap=round]
        (44mm, 36.5mm) to[out=46,in=-46] (44mm, 63.5mm);
}

% Second sound wave (Curved line)
\opt{w2}{
    \draw[line width=4.3mm,line cap=round]
        (57.5mm, 27.5mm) to[out=46,in=-46] (57.5mm, 72.5mm);
}

% Third sound wave (Curved line)
\opt{w3}{
    \draw[line width=4.3mm,line cap=round]
        (72mm, 18.5mm) to[out=46,in=-46] (72mm, 81.5mm);
}

% Always draw a transparent copy of the third wave so the images will all have
% the same width.
\draw[line width=4.3mm,line cap=round,opacity=0]
    (72mm, 18.5mm) to[out=46,in=-46] (72mm, 81.5mm);

\end{tikzpicture}
\end{document}

'
%   xelatex -jobname=Volume1 '\def\UseOption{w1}% This file is part of Background Music.
%
% Background Music is free software: you can redistribute it and/or
% modify it under the terms of the GNU General Public License as
% published by the Free Software Foundation, either version 2 of the
% License, or (at your option) any later version.
%
% Background Music is distributed in the hope that it will be useful,
% but WITHOUT ANY WARRANTY; without even the implied warranty of
% MERCHANTABILITY or FITNESS FOR A PARTICULAR PURPOSE. See the
% GNU General Public License for more details.
%
% You should have received a copy of the GNU General Public License
% along with Background Music. If not, see <http://www.gnu.org/licenses/>.

%
% VolumeIcons.tex
%
% Build with XeTeX:
%   xelatex -jobname=Volume0 '\def\UseOption{}\input{VolumeIcons.tex}'
%   xelatex -jobname=Volume1 '\def\UseOption{w1}\input{VolumeIcons.tex}'
%   xelatex -jobname=Volume2 '\def\UseOption{w1,w2}\input{VolumeIcons.tex}'
%   xelatex -jobname=Volume3 '\def\UseOption{w1,w2,w3}\input{VolumeIcons.tex}'
%   for n in 0 1 2 3; do mv Volume$n.pdf ../BGMApp/BGMApp/Images.xcassets/Volume$n.imageset/; done
%
% Might build correctly with regular LaTeX. I haven't tried it.
%

\documentclass[tikz]{standalone}
\usepackage{tikz}
% "dummyOption" prevents "Package optional Warning: No options were selected,
% so all optional text will be printed" when building Volume0.pdf.
\usepackage[dummyOption]{optional}

\begin{document}
\begin{tikzpicture}

% Speaker (Rounded box and triangle)
\fill[rounded corners=5mm]
    (0mm, 62.5mm) rectangle (25mm, 37.5mm) {};
\draw[rounded corners=2.5mm,fill=black]
    (3mm, 50mm)--(34mm, 76.5mm)--(34mm, 23.5mm)--cycle;

% First sound wave (Curved line)
\opt{w1}{
    \draw[line width=4.3mm,line cap=round]
        (44mm, 36.5mm) to[out=46,in=-46] (44mm, 63.5mm);
}

% Second sound wave (Curved line)
\opt{w2}{
    \draw[line width=4.3mm,line cap=round]
        (57.5mm, 27.5mm) to[out=46,in=-46] (57.5mm, 72.5mm);
}

% Third sound wave (Curved line)
\opt{w3}{
    \draw[line width=4.3mm,line cap=round]
        (72mm, 18.5mm) to[out=46,in=-46] (72mm, 81.5mm);
}

% Always draw a transparent copy of the third wave so the images will all have
% the same width.
\draw[line width=4.3mm,line cap=round,opacity=0]
    (72mm, 18.5mm) to[out=46,in=-46] (72mm, 81.5mm);

\end{tikzpicture}
\end{document}

'
%   xelatex -jobname=Volume2 '\def\UseOption{w1,w2}% This file is part of Background Music.
%
% Background Music is free software: you can redistribute it and/or
% modify it under the terms of the GNU General Public License as
% published by the Free Software Foundation, either version 2 of the
% License, or (at your option) any later version.
%
% Background Music is distributed in the hope that it will be useful,
% but WITHOUT ANY WARRANTY; without even the implied warranty of
% MERCHANTABILITY or FITNESS FOR A PARTICULAR PURPOSE. See the
% GNU General Public License for more details.
%
% You should have received a copy of the GNU General Public License
% along with Background Music. If not, see <http://www.gnu.org/licenses/>.

%
% VolumeIcons.tex
%
% Build with XeTeX:
%   xelatex -jobname=Volume0 '\def\UseOption{}\input{VolumeIcons.tex}'
%   xelatex -jobname=Volume1 '\def\UseOption{w1}\input{VolumeIcons.tex}'
%   xelatex -jobname=Volume2 '\def\UseOption{w1,w2}\input{VolumeIcons.tex}'
%   xelatex -jobname=Volume3 '\def\UseOption{w1,w2,w3}\input{VolumeIcons.tex}'
%   for n in 0 1 2 3; do mv Volume$n.pdf ../BGMApp/BGMApp/Images.xcassets/Volume$n.imageset/; done
%
% Might build correctly with regular LaTeX. I haven't tried it.
%

\documentclass[tikz]{standalone}
\usepackage{tikz}
% "dummyOption" prevents "Package optional Warning: No options were selected,
% so all optional text will be printed" when building Volume0.pdf.
\usepackage[dummyOption]{optional}

\begin{document}
\begin{tikzpicture}

% Speaker (Rounded box and triangle)
\fill[rounded corners=5mm]
    (0mm, 62.5mm) rectangle (25mm, 37.5mm) {};
\draw[rounded corners=2.5mm,fill=black]
    (3mm, 50mm)--(34mm, 76.5mm)--(34mm, 23.5mm)--cycle;

% First sound wave (Curved line)
\opt{w1}{
    \draw[line width=4.3mm,line cap=round]
        (44mm, 36.5mm) to[out=46,in=-46] (44mm, 63.5mm);
}

% Second sound wave (Curved line)
\opt{w2}{
    \draw[line width=4.3mm,line cap=round]
        (57.5mm, 27.5mm) to[out=46,in=-46] (57.5mm, 72.5mm);
}

% Third sound wave (Curved line)
\opt{w3}{
    \draw[line width=4.3mm,line cap=round]
        (72mm, 18.5mm) to[out=46,in=-46] (72mm, 81.5mm);
}

% Always draw a transparent copy of the third wave so the images will all have
% the same width.
\draw[line width=4.3mm,line cap=round,opacity=0]
    (72mm, 18.5mm) to[out=46,in=-46] (72mm, 81.5mm);

\end{tikzpicture}
\end{document}

'
%   xelatex -jobname=Volume3 '\def\UseOption{w1,w2,w3}% This file is part of Background Music.
%
% Background Music is free software: you can redistribute it and/or
% modify it under the terms of the GNU General Public License as
% published by the Free Software Foundation, either version 2 of the
% License, or (at your option) any later version.
%
% Background Music is distributed in the hope that it will be useful,
% but WITHOUT ANY WARRANTY; without even the implied warranty of
% MERCHANTABILITY or FITNESS FOR A PARTICULAR PURPOSE. See the
% GNU General Public License for more details.
%
% You should have received a copy of the GNU General Public License
% along with Background Music. If not, see <http://www.gnu.org/licenses/>.

%
% VolumeIcons.tex
%
% Build with XeTeX:
%   xelatex -jobname=Volume0 '\def\UseOption{}\input{VolumeIcons.tex}'
%   xelatex -jobname=Volume1 '\def\UseOption{w1}\input{VolumeIcons.tex}'
%   xelatex -jobname=Volume2 '\def\UseOption{w1,w2}\input{VolumeIcons.tex}'
%   xelatex -jobname=Volume3 '\def\UseOption{w1,w2,w3}\input{VolumeIcons.tex}'
%   for n in 0 1 2 3; do mv Volume$n.pdf ../BGMApp/BGMApp/Images.xcassets/Volume$n.imageset/; done
%
% Might build correctly with regular LaTeX. I haven't tried it.
%

\documentclass[tikz]{standalone}
\usepackage{tikz}
% "dummyOption" prevents "Package optional Warning: No options were selected,
% so all optional text will be printed" when building Volume0.pdf.
\usepackage[dummyOption]{optional}

\begin{document}
\begin{tikzpicture}

% Speaker (Rounded box and triangle)
\fill[rounded corners=5mm]
    (0mm, 62.5mm) rectangle (25mm, 37.5mm) {};
\draw[rounded corners=2.5mm,fill=black]
    (3mm, 50mm)--(34mm, 76.5mm)--(34mm, 23.5mm)--cycle;

% First sound wave (Curved line)
\opt{w1}{
    \draw[line width=4.3mm,line cap=round]
        (44mm, 36.5mm) to[out=46,in=-46] (44mm, 63.5mm);
}

% Second sound wave (Curved line)
\opt{w2}{
    \draw[line width=4.3mm,line cap=round]
        (57.5mm, 27.5mm) to[out=46,in=-46] (57.5mm, 72.5mm);
}

% Third sound wave (Curved line)
\opt{w3}{
    \draw[line width=4.3mm,line cap=round]
        (72mm, 18.5mm) to[out=46,in=-46] (72mm, 81.5mm);
}

% Always draw a transparent copy of the third wave so the images will all have
% the same width.
\draw[line width=4.3mm,line cap=round,opacity=0]
    (72mm, 18.5mm) to[out=46,in=-46] (72mm, 81.5mm);

\end{tikzpicture}
\end{document}

'
%   for n in 0 1 2 3; do mv Volume$n.pdf ../BGMApp/BGMApp/Images.xcassets/Volume$n.imageset/; done
%
% Might build correctly with regular LaTeX. I haven't tried it.
%

\documentclass[tikz]{standalone}
\usepackage{tikz}
% "dummyOption" prevents "Package optional Warning: No options were selected,
% so all optional text will be printed" when building Volume0.pdf.
\usepackage[dummyOption]{optional}

\begin{document}
\begin{tikzpicture}

% Speaker (Rounded box and triangle)
\fill[rounded corners=5mm]
    (0mm, 62.5mm) rectangle (25mm, 37.5mm) {};
\draw[rounded corners=2.5mm,fill=black]
    (3mm, 50mm)--(34mm, 76.5mm)--(34mm, 23.5mm)--cycle;

% First sound wave (Curved line)
\opt{w1}{
    \draw[line width=4.3mm,line cap=round]
        (44mm, 36.5mm) to[out=46,in=-46] (44mm, 63.5mm);
}

% Second sound wave (Curved line)
\opt{w2}{
    \draw[line width=4.3mm,line cap=round]
        (57.5mm, 27.5mm) to[out=46,in=-46] (57.5mm, 72.5mm);
}

% Third sound wave (Curved line)
\opt{w3}{
    \draw[line width=4.3mm,line cap=round]
        (72mm, 18.5mm) to[out=46,in=-46] (72mm, 81.5mm);
}

% Always draw a transparent copy of the third wave so the images will all have
% the same width.
\draw[line width=4.3mm,line cap=round,opacity=0]
    (72mm, 18.5mm) to[out=46,in=-46] (72mm, 81.5mm);

\end{tikzpicture}
\end{document}

'
%   for n in 0 1 2 3; do mv Volume$n.pdf ../BGMApp/BGMApp/Images.xcassets/Volume$n.imageset/; done
%
% Might build correctly with regular LaTeX. I haven't tried it.
%

\documentclass[tikz]{standalone}
\usepackage{tikz}
% "dummyOption" prevents "Package optional Warning: No options were selected,
% so all optional text will be printed" when building Volume0.pdf.
\usepackage[dummyOption]{optional}

\begin{document}
\begin{tikzpicture}

% Speaker (Rounded box and triangle)
\fill[rounded corners=5mm]
    (0mm, 62.5mm) rectangle (25mm, 37.5mm) {};
\draw[rounded corners=2.5mm,fill=black]
    (3mm, 50mm)--(34mm, 76.5mm)--(34mm, 23.5mm)--cycle;

% First sound wave (Curved line)
\opt{w1}{
    \draw[line width=4.3mm,line cap=round]
        (44mm, 36.5mm) to[out=46,in=-46] (44mm, 63.5mm);
}

% Second sound wave (Curved line)
\opt{w2}{
    \draw[line width=4.3mm,line cap=round]
        (57.5mm, 27.5mm) to[out=46,in=-46] (57.5mm, 72.5mm);
}

% Third sound wave (Curved line)
\opt{w3}{
    \draw[line width=4.3mm,line cap=round]
        (72mm, 18.5mm) to[out=46,in=-46] (72mm, 81.5mm);
}

% Always draw a transparent copy of the third wave so the images will all have
% the same width.
\draw[line width=4.3mm,line cap=round,opacity=0]
    (72mm, 18.5mm) to[out=46,in=-46] (72mm, 81.5mm);

\end{tikzpicture}
\end{document}

'
%   for n in 0 1 2 3; do mv Volume$n.pdf ../BGMApp/BGMApp/Images.xcassets/Volume$n.imageset/; done
%
% Might build correctly with regular LaTeX. I haven't tried it.
%

\documentclass[tikz]{standalone}
\usepackage{tikz}
% "dummyOption" prevents "Package optional Warning: No options were selected,
% so all optional text will be printed" when building Volume0.pdf.
\usepackage[dummyOption]{optional}

\begin{document}
\begin{tikzpicture}

% Speaker (Rounded box and triangle)
\fill[rounded corners=5mm]
    (0mm, 62.5mm) rectangle (25mm, 37.5mm) {};
\draw[rounded corners=2.5mm,fill=black]
    (3mm, 50mm)--(34mm, 76.5mm)--(34mm, 23.5mm)--cycle;

% First sound wave (Curved line)
\opt{w1}{
    \draw[line width=4.3mm,line cap=round]
        (44mm, 36.5mm) to[out=46,in=-46] (44mm, 63.5mm);
}

% Second sound wave (Curved line)
\opt{w2}{
    \draw[line width=4.3mm,line cap=round]
        (57.5mm, 27.5mm) to[out=46,in=-46] (57.5mm, 72.5mm);
}

% Third sound wave (Curved line)
\opt{w3}{
    \draw[line width=4.3mm,line cap=round]
        (72mm, 18.5mm) to[out=46,in=-46] (72mm, 81.5mm);
}

% Always draw a transparent copy of the third wave so the images will all have
% the same width.
\draw[line width=4.3mm,line cap=round,opacity=0]
    (72mm, 18.5mm) to[out=46,in=-46] (72mm, 81.5mm);

\end{tikzpicture}
\end{document}

